\documentclass[11pt]{amsart}
\usepackage{amsmath,amssymb,amsthm,mathtools}
\usepackage[margin=1in]{geometry}
\usepackage{booktabs,tabularx}

\emergencystretch=4em
\tolerance=3000
\usepackage{hyperref}
\hypersetup{colorlinks=true,linkcolor=blue,citecolor=blue,urlcolor=blue}

\newtheorem{theorem}{Theorem}[section]
\newtheorem{lemma}[theorem]{Lemma}
\newtheorem{proposition}[theorem]{Proposition}
\newtheorem{corollary}[theorem]{Corollary}
\theoremstyle{definition}
\newtheorem{definition}[theorem]{Definition}
\newtheorem{remark}[theorem]{Remark}
\newtheorem{problem}[theorem]{Open Problem}
\theoremstyle{remark}
\newtheorem*{acknowledgement}{Acknowledgement}
\numberwithin{equation}{section}

\newcommand{\Q}{\mathbb{Q}}
\newcommand{\R}{\mathbb{R}}
\newcommand{\C}{\mathbb{C}}
\newcommand{\Z}{\mathbb{Z}}
\newcommand{\N}{\mathbb{N}}
\newcommand{\Qbar}{\overline{\mathbb{Q}}}
\newcommand{\trdeg}{\operatorname{trdeg}}
\newcommand{\irr}{\mu}
\newcommand{\ee}{e}
\DeclareMathOperator{\wpfun}{\wp}

\title[On the Irrationality of $e+\pi$]{Continued-Fraction Windows, Quadratic Orbit Spectra for Rational Affine Images of $e$, and Obstructions to the Irrationality of $e+\pi$}
\author{Anonymous author}
\thanks{The manuscript contains proofs, conditional deductions, and explicitly marked open steps; no priority is claimed for elementary combinations of classical tools. \textbf{This paper does not prove the irrationality of $e+\pi$; that problem remains open.} Every external theorem is cited with its verification status; results depending on unreviewed preprints are explicitly conditional. A summary theorem--source--application audit table is given in Appendix~\ref{app:audit}, with the declaration-level audit in the accompanying archive.}
\date{Integrated revision, November 22, 2026}
\subjclass[2020]{Primary 11J82; Secondary 11J85, 11J72, 11A55.}
\keywords{Irrationality, transcendence, $e+\pi$, irrationality measure, Schanuel's conjecture, Flint Hills series, continued fractions.}

\begin{document}

\begin{abstract}
The rationality of $e+\pi$ is open. We study the continued-fraction
windows indexed by $H_j=\lfloor Q_j\sqrt j\rfloor$, where $Q_j$ is the
denominator of the convergent to $e$ immediately preceding the partial
quotient $2j$. If $e+x=A/B$ with $B>0$, each window with $j>9B^2/4$
contains an odd number of new convergents, at least three, and the
window counts have period $2B^2$. For a rational affine input
$x=(A+Ce)/B$, $C\ne0$, the counts and the digit suffixes have period
$2B|C|$, and the first digit increases by an explicit amount at that lag.
In a primitive representation, five consecutive single-digit windows
in the locked range characterize $|C|=1$, $B\in\{2,4\}$ and odd $A$.
Four windows do not suffice. For $x=A/B-e$ we also bound the lengths
of runs of ones, obtain a lower bound for approximation by the
quadratic orbit of the golden ratio, and determine the exact
bounded golden-ratio orbit limit spectrum by a finite exact algorithm for
every reduced $A/B$. In particular, we establish complete Case~1 and Case~2
classifications of the exact constants $c_\varphi$, all extremal families,
and the second smallest limit points for all coprime denominator classes
$B=1,\dots,6$ and for the multi-scale affine targets $(1+e)/6$, $(1+e)/3$,
$(1+2e)/3$, and $22/7-e$.
We characterize the rational affine span of $e$ using
either reduced denominator quotient, and give signed-residual,
Dirichlet and factorial criteria together with their unresolved inputs
at $\pi$.
\end{abstract}

\maketitle

\section{Introduction}\label{sec:intro}

Let $\ee=\exp(1)$ and let $\pi$ be the circle constant. The problem addressed in this paper is:

\begin{problem}\label{prob:main}
Is $e+\pi$ irrational? Is it transcendental? Are $e$ and $\pi$ algebraically independent over $\Q$?
\end{problem}

The transcendence of $\ee$ was proved by Hermite in 1873 \cite{Hermite1873}, and that of $\pi$ by Lindemann in 1882 \cite{Lindemann1882}. Nevertheless, essentially nothing is known about the arithmetic nature of sums, products, differences, or quotients of these two constants. The OEIS Wiki records the striking fact that there is no pair of nonzero integers $m,n$ for which the rationality or irrationality of $m\pi+n\ee$ is known \cite{OEIS}. One elementary unconditional result in the immediate vicinity of Problem~\ref{prob:main} is the disjunctive statement that at least one of $e+\pi$ and $e\pi$ is transcendental (Theorem~\ref{thm:vieta} below).

\subsection{Contributions} The principal arithmetic results are the affine window period theorem (Theorem~\ref{thm:affine-window}), the bounded golden-orbit finite spectrum theorem (Theorem~\ref{thm:golden-finite-shift}), and the ideal-class trace density theorem (Theorem~\ref{thm:trace-density-one}). The contributions are:
\begin{enumerate}
\item Window periods and the converse to denominator locking (Section~\ref{sec:windows}): the Dirichlet-scale denominator restriction (Proposition~\ref{prop:dirichlet-divisor}), the two-primitive-witness test (Proposition~\ref{prop:two-window}), the signed rounding parity test (Proposition~\ref{prop:signed-round}), the explicit period $2B^2$ forced by rationality (Theorem~\ref{thm:explicit-period}), the complete digit shift and its corollary tests (Theorem~\ref{thm:digit-shift}, Corollary~\ref{cor:digit-statistics}), the converse to denominator locking (Theorem~\ref{thm:gcd-converse}), the slope isolation (Corollary~\ref{cor:recover-slope}), and the affine window period $2B|C|$ with its digit increment (Theorem~\ref{thm:affine-window}); for a primitive representation, five consecutive single-digit windows in the locked range characterize $|C|=1$, $B\in\{2,4\}$ and odd $A$, while four windows do not suffice.
\item An effective finite evaluation of the bounded golden-orbit limit
spectrum for nonconstant rational affine images of $e$
(Theorem~\ref{thm:golden-finite-shift}), together with exact extremal
classifications for $A/B-e$ with $1\le B\le6$ and for several mixed-scale
targets (Corollary~\ref{cor:golden-finite-shift-examples} and
Theorems~\ref{thm:golden-b5-exact}--\ref{thm:golden-multiscale-exact}).
For any real quadratic orbit anchor, the finite-phase theorem, trace and
residue tests, and uniform convergence estimate give an effective extension
(Theorems~\ref{thm:general-quadratic-phase}, \ref{thm:pure-trace-test}
and~\ref{thm:phase-convergence-rate}).
\item The ideal-class trace test and uniform interval counting
(Section~\ref{sec:quadratic}):
Theorems~\ref{thm:trace-all-norms}--\ref{thm:trace-common-progression}
factor the full trace test through invertible ideal classes and
give common progressions for orbit spectra in the same field.
Theorem~\ref{thm:trace-density-one} and
Corollary~\ref{cor:trace-common-density} give location-uniform density one in
the admissible trace lattice and the leading common count.
\item Elementary reductions and conditional implications (Section~\ref{sec:new}): the measure reduction (Theorem~\ref{thm:measure-reduction}), the continued-fraction necessary condition (Theorem~\ref{thm:cf}), the series rigidity propositions (Propositions~\ref{prop:window} and~\ref{prop:recurrence}), the four-exponentials obstruction (Proposition~\ref{prop:four-exp}), the translation identity for the $\irr(\pi)$ program (Theorem~\ref{thm:translation}), the scaling-closure proposition (Proposition~\ref{prop:scaling}), and the conditional Flint Hills bridge (Theorem~\ref{thm:flint}), together with the exclusion map, the surviving targets, and the adaptive-window and factorial digit-quantization criteria (Theorems~\ref{thm:adaptive-strip} and~\ref{thm:factorial-quantization}), with the required infinite behavior at $\pi$ still open.
\item A systematic obstruction analysis (\S\ref{sec:obstructions}) of four proof routes and four candidate bridge identities, with the limitations of each examined construction stated explicitly. These observations do not exclude unexamined constructions or prove that an entire method is impossible.
\item Proofs of selected classical unconditional results (Section~\ref{sec:unconditional}), including a compressed technical baseline comparing Fourier's and Niven's irrationality proofs.
\item Machine verification and rigorous certificates (Appendix~\ref{sec:computation}): twelve symbolic verifications plus a separate cosine-integral check, and a 450-digit interval-arithmetic certificate.
\item A summary audit table (Appendix~\ref{app:audit}) recording, for the principal external facts, their statement, source, verification method, status, and application sites; the full declaration-level audit is in the archive (\texttt{LEMMA\_AUDIT.md}, \texttt{PROOF\_GAPS.md}, \texttt{FINAL\_REVIEW.md}, and \texttt{verification/\allowbreak audit\_\allowbreak final\_\allowbreak 20261122.log}).
\item A machine-verified formalization in Lean~4 (Appendix~\ref{sec:lean}): a non-degenerate definition of the irrationality measure and kernel-checked derivations of the registered declarations listed in the archive with no \texttt{sorryAx} in the audited declarations. The \emph{Flint Hills case} of Meiburg's criterion is fully formalized (circle spacing, dyadic-block counting, Jordan bound, geometric summation), replacing the preprint axiom by a theorem; the Alekseyev direction (convergence $\Rightarrow\irr(\pi)\le5/2$) is now machine-verified via a nearest-integer argument, with the earlier density obstruction preserved as a historical record.
\item A kernel-checked quantitative reduction of the escape target to a lattice-escape subproblem P1 (\S\ref{sec:lean-reduction}): match $\Rightarrow$ super-approximation $\Rightarrow$ upper envelope tending to $1/8$ $\Rightarrow$ collision with the classification gap $\Rightarrow$ $B\pi\in L$. This is a change of target, not a reduction in difficulty; P1 is of the same family as the original problem and remains open.
\item A reorganization of the four combinations under two expected-false necessary conditions (Section~\ref{sec:expected-false}): ``$\pi$ is a Hurwitzian number'' (by Perron's theorem) and ``$\pi\in\mathbf{E}$'' (via the Fischler--Rivoal $E/G$-value rings). The routes $\irr(\pi)>2$ and Flint Hills divergence are flagged as expected false and abandoned as attack routes.
\end{enumerate}

The golden-orbit finite-spectrum theorem
(Theorem~\ref{thm:golden-finite-shift}) is the principal arithmetic
result: it uses the periodic window data to prove that every bounded
asymptotic error comes from one of three explicit, realizable families.
Bugeaud's M\"obius formula
\cite[eq.~(3.4)]{BugeaudQuad} supplies the local error identity. The finite
search and the completeness of its output use the rational-shift window
structure specific to this paper. The companion paper studies the metric
spectrum of the distinct quantity $\limsup_j j\|Q_jx\|$, including the
Hausdorff gauges of its exact levels \cite{CompanionGauge}.

\subsection{Honesty statement} This paper does not resolve Problem~\ref{prob:main}. No contradiction to $e+\pi\in\Q$ is obtained from the results used here. This is a statement about the present arguments, not an independence theorem or a proof that existing mathematics cannot resolve the problem. Results that depend on unreviewed preprints (Bertolin \cite{Bertolin2026}, Fischler--Rivoal \cite{FischlerRivoal2026}, and the external claim \cite{OpenAI2026}) are stated conditionally and flagged as such; Bai's preprint \cite{Bai2026} is used only for a numerical comparison. Meiburg's criterion \cite{Meiburg2022} was originally used conditionally; its Flint Hills case is now formalized in Lean (Appendix~\ref{sec:lean}), so the formal version of the bridge no longer rests on that preprint, while the paper's general statement of Meiburg's criterion remains preprint-dependent. The measure-theoretic intuition that $e+\pi$ is ``almost certainly'' transcendental is used only heuristically and never as evidence. One attribution error made in an intermediate round (concerning the Flint Hills criterion) was detected by primary-source verification and is corrected in Section~\ref{sec:flint-bridge}; the correction is part of the permanent record.

\subsection{Organization} Section~\ref{sec:prelim} fixes notation, states the classical theorems used, and records the unconditional baseline. Section~\ref{sec:windows} establishes the window periods and the converse to denominator locking for rational affine images of $e$. Section~\ref{sec:golden} evaluates the bounded golden-orbit limit spectrum. Section~\ref{sec:quadratic} treats general real quadratic orbits, the ideal-class trace test, and uniform interval counting. Section~\ref{sec:new} assembles the reductions, the systematic obstruction analysis, the conditional frameworks, the exclusion map, and the surviving targets. Section~\ref{sec:conclusion} states the research priorities and limitations. The computational certificates, the Lean~4 formalization, the audit table, and reproducibility instructions are collected in Appendices~\ref{sec:computation}--\ref{app:repro}.

\section{Preliminaries, classical theorems, and the unconditional baseline}\label{sec:prelim}

\begin{definition}
A complex number is \emph{algebraic} if it is a root of a nonzero polynomial with rational coefficients; otherwise it is \emph{transcendental}. We write $\Qbar$ for the field of algebraic numbers. A finite set $\{x_1,\dots,x_n\}\subset\C$ is \emph{algebraically independent over $\Q$} if no nonzero polynomial $P\in\Q[X_1,\dots,X_n]$ satisfies $P(x_1,\dots,x_n)=0$. The \emph{transcendence degree} $\trdeg_\Q K$ of a field extension $K/\Q$ is the cardinality of a maximal algebraically independent subset.
\end{definition}

\begin{definition}[Irrationality measure]\label{def:measure}
For an irrational $x\in\R$, define
\[
\irr(x)=\sup\{\nu\ge0:\ |x-p/q|<q^{-\nu}
\text{ for infinitely many }(p,q)\in\Z\times\N_{>0}\},
\]
with supremum in $[0,\infty]$. Equivalently, this is the infimum of positive exponents for which only finitely many such pairs exist. Exact equality $x=p/q$ cannot occur because $x$ is irrational.
\end{definition}

For irrational arguments we use $\irr(x)\ge2$ (Dirichlet) and $\irr(x)=2$ for algebraic $x$ (Roth \cite{Roth1955}). We make no convention for rational arguments in the mathematical discussion. The Lean definition in \S\ref{sec:lean-def} extends the displayed supremum to all reals; because it includes exact approximations with unreduced pairs, its value at a rational is $+\infty$, not the alternative conventional value $1$. All arithmetic-measure applications here concern irrational arguments.

\begin{definition}
For an irrational $x$, its \emph{simple continued fraction} is $x=[a_0;a_1,a_2,\dots]$ with convergents $p_n/q_n$. We use the classical bound $|x-p_n/q_n|<1/q_n^2$.
\end{definition}

The following classical theorems are used as black boxes; each is cited with its source and is not reproved here.

\begin{theorem}[Hermite--Lindemann; {\cite{Hermite1873,Lindemann1882}}]\label{thm:HL}
If $\alpha\in\Qbar\setminus\{0\}$, then $\ee^\alpha$ is transcendental. In particular, $\ee$ and $\pi$ are transcendental.
\end{theorem}

\begin{theorem}[Lindemann--Weierstrass; {\cite{Lindemann1882}}]\label{thm:LW}
If $\alpha_1,\dots,\alpha_n\in\Qbar$ are linearly independent over $\Q$, then $\ee^{\alpha_1},\dots,\ee^{\alpha_n}$ are algebraically independent over $\Q$.
\end{theorem}

\begin{theorem}[Gelfond--Schneider; {\cite{GS1934}}]\label{thm:GS}
If $\alpha\in\Qbar\setminus\{0,1\}$ and $\beta\in\Qbar\setminus\Q$, then every value of $\alpha^\beta$ is transcendental.
\end{theorem}

\begin{theorem}[Schneider, elliptic case; {\cite{Duke}}]\label{thm:schneider}
Let $\wp$ be a Weierstrass function with algebraic invariants $g_2,g_3$. If $z\in\Qbar\setminus\{0\}$, then $\wp(z)$ is transcendental.
\end{theorem}

\begin{definition}[Weierstrass function]
For a lattice $\Omega\subset\C$, $\wp_\Omega(z)=z^{-2}+\sum_{\omega\in\Omega\setminus\{0\}}\bigl((z-\omega)^{-2}-\omega^{-2}\bigr)$. Its invariants are $g_2(\Omega)=60\sum_{\omega\ne0}\omega^{-4}$ and $g_3(\Omega)=140\sum_{\omega\ne0}\omega^{-6}$. The lattice has \emph{complex multiplication} (CM) if $\operatorname{End}(\C/\Omega)\otimes\Q$ is an imaginary quadratic field.
\end{definition}

We also use Euler's identity $\ee^{i\pi}=-1$ and, in Section~\ref{sec:conditional}, the statement of Schanuel's conjecture.


\subsection{The disjunctive transcendence theorem}\label{sec:unconditional}

\begin{theorem}\label{thm:vieta}
At least one of $e+\pi$ and $e\pi$ is transcendental.
\end{theorem}

\begin{proof}
Set $S=e+\pi$ and $P=e\pi$. By Vieta's formulas, $e$ and $\pi$ are precisely the two roots of
\begin{equation}\label{eq:quadratic}
x^2-Sx+P=0.
\end{equation}
Suppose, for contradiction, that both $S$ and $P$ are algebraic. The field $\Qbar$ of algebraic numbers is closed under addition, subtraction, multiplication, division, and extraction of square roots. Hence the quadratic formula applied to \eqref{eq:quadratic},
\[
x=\frac{S\pm\sqrt{S^2-4P}}{2},
\]
shows that both roots are algebraic. Thus $e$ and $\pi$ would both be algebraic, contradicting Theorem~\ref{thm:HL}. Therefore $S$ and $P$ cannot both be algebraic; at least one is transcendental.
\end{proof}

\begin{remark}
The argument identifies neither which of the two is transcendental nor whether either could be rational. Any proof of the irrationality of $e+\pi$ must supply information beyond this elementary argument.
\end{remark}

\subsection{At most one rational among four combinations}

\begin{theorem}\label{thm:four}
Among the four numbers $e+\pi$, $e-\pi$, $e\pi$, $e/\pi$, at most one is rational.
\end{theorem}

\begin{proof}
There are six unordered pairs; we exclude each. Throughout, we use that $\Q$ is closed under the four arithmetic operations, and that $e,\pi>0$ with $e\ne\pi$.

\emph{Pair 1: $e+\pi$ and $e-\pi$.} If both are rational, their sum $2e$ and difference $2\pi$ are rational, so $e$ and $\pi$ are rational, contradicting Theorem~\ref{thm:HL}.

\emph{Pair 2: $e+\pi$ and $e\pi$.} If both are rational, then by \eqref{eq:quadratic}, $e$ and $\pi$ are roots of a quadratic polynomial with rational coefficients, hence algebraic, contradicting Theorem~\ref{thm:HL}.

\emph{Pair 3: $e-\pi$ and $e\pi$.} If both are rational, then $e$ and $-\pi$ are the roots of $x^2-(e-\pi)x-e\pi=0$, a quadratic with rational coefficients; hence $e$ and $\pi$ are algebraic, contradiction.

\emph{Pair 4: $e+\pi$ and $e/\pi$.} Suppose $e+\pi=r\in\Q$ and $e/\pi=s\in\Q$. Then $e=s\pi$, so $(s+1)\pi=r$. Since $e,\pi>0$, we have $s>0$, hence $s+1\ne0$, and $\pi=r/(s+1)\in\Q$, contradiction.

\emph{Pair 5: $e-\pi$ and $e/\pi$.} Suppose $e-\pi=r\in\Q$ and $e/\pi=s\in\Q$. Then $(s-1)\pi=r$. Since $e\ne\pi$, $s\ne1$, so $\pi=r/(s-1)\in\Q$, contradiction.

\emph{Pair 6: $e\pi$ and $e/\pi$.} If $e\pi=r\in\Q$ and $e/\pi=s\in\Q$, then multiplying gives $e^2=rs\in\Q$, so $e$ is a root of $x^2-rs\in\Q[x]$, hence algebraic, contradiction.

All six pairs are excluded, so at most one of the four numbers is rational.
\end{proof}

\subsection{Transcendence of \texorpdfstring{$e^\pi$}{e\^{}pi}}

\begin{theorem}\label{thm:epi}
$e^\pi$ is transcendental.
\end{theorem}

\begin{proof}
By Euler's identity, $\ee^{i\pi}=-1$. Taking the principal branch,
\[
\ee^\pi=(\ee^{i\pi})^{-i}=(-1)^{-i}.
\]
The base $-1$ is algebraic and distinct from $0$ and $1$. The exponent $-i$ is algebraic (it satisfies $x^2+1=0$) and irrational. By the Gelfond--Schneider theorem (Theorem~\ref{thm:GS}), $(-1)^{-i}$ is transcendental. Hence $\ee^\pi$ is transcendental.
\end{proof}

\begin{remark}
This proof exemplifies the only known paradigm for such results: transform the target constant into the standard input of a proved theorem. The transformation here exists solely because of Euler's identity. No analogous transformation is known for $e+\pi$.
\end{remark}

\subsection{Technical baseline: two classical irrationality proofs}

\begin{theorem}[Fourier]\label{thm:fourier}
$\ee$ is irrational.
\end{theorem}

\begin{theorem}[Niven]\label{thm:niven}
$\pi$ is irrational.
\end{theorem}

Both statements are weaker than the Hermite--Lindemann theorem
(Theorem~\ref{thm:HL}) and are recorded only as a technical
baseline. Fourier's proof for $\ee$ multiplies the truncated
exponential series by $n!$: the main term clears denominators and
is an integer, while the tail
$I_n=n!\bigl(\ee-\sum_{k\le n}1/k!\bigr)$ satisfies $0<I_n<1/n$.
Niven's proof for $\pi$ integrates $f(x)=x^n(a-bx)^n/n!$ against
$\sin x$ over $[0,\pi]$; every even derivative at the endpoints is
integral, and the integral is a positive integer tending to $0$,
hence eventually lies strictly in $(0,1)$. Both proofs construct a
quantity that is an integer under the rationality hypothesis and
independently lies strictly in $(0,1)$. Three ingredients are
required: (i) a rapidly convergent series or integral
representation; (ii) a ``denominator-clearing'' mechanism making
the main term integral; (iii) separation of main term and tail
with controlled error. The obstruction analysis of
\S\ref{sec:obstructions} shows that each attempted route to
$e+\pi$ fails at ingredient (ii) or (iii), at an identifiable
structural point. Expository accounts of both proofs are given in
\cite{Conrad}.
\subsection{Stronger known results (statements only)}\label{sub:stronger-known}

For context we record the following published results, used later only as stated: Nesterenko's theorem that $\pi$, $\ee^\pi$, $\Gamma(1/4)$ are algebraically independent over $\Q$ \cite{Nesterenko1996} (which already gives algebraic independence of $\pi$ and $\ee^\pi$); Tijdeman's work on algebraic independence of exponential values \cite{Tijdemann1971} (the precise ``at least two of $\ee,\pi,\ee^\pi$'' corollary is not verified from the original and is not relied upon); the conditional result that if $\ee^{\pi^2}$ is algebraic then $\ee$ and $\pi$ are algebraically independent, due to Brownawell/Waldschmidt \cite{Brownawell1974}; and the six exponentials theorem of Siegel--Lang--Ramachandra \cite{Lang1966,Ramachandra1968}. A systematic treatment of algebraic independence, including Nesterenko's method, is given in the monograph \cite{Zhu2024}.
\section{Window periods and the converse to denominator locking}\label{sec:windows}

This chapter establishes the periodic window structure that
rationality forces on the continued fraction of $e+x$.
Theorem~\ref{thm:explicit-period} gives an explicit necessary
period under rationality. Theorem~\ref{thm:digit-shift} gives the
complete digit change and forces eventual periodicity modulo four
in ordinary digit indices. Corollary~\ref{cor:digit-statistics}
gives tests in terms of the growth of ordinary convergent
denominators and the multiplicities of large digit values.
Theorem~\ref{thm:gcd-converse} identifies unboundedness of either
denominator quotient for $\pi$ with rational linear independence
of $1,e,\pi$. Corollary~\ref{cor:recover-slope} isolates the slope
needed for the sum. Whether the actual $\pi$ violates the slope
condition remains unresolved.
\subsection{Denominator restrictions at the Dirichlet scale}\label{sec:current-results}
Proposition~\ref{prop:farey-gap} excludes every denominator
$b<3635695100$ ($b\ge1$) and every adjacent match for those denominators
at $j\ge8$. Both conclusions were kernel-checked. For unrestricted $b$,
Proposition~\ref{prop:nest-transfer} transfers a hypothetical rational
affine relation to algebraic independence of $e,\exp(e),\Gamma(1/4)$;
this paper deduction gives a necessary condition. Proposition~\ref{prop:rational-span}
identifies P1 with full rational linear independence of $1,e,\pi$.
The fixed slope of the sum allows a more targeted divisibility test.

\begin{proposition}[A denominator restriction at the Dirichlet scale]
\label{prop:dirichlet-divisor}
Write $q=Q_j$, $P=P_j$, and $H=\lfloor q\sqrt j\rfloor$.
If $B(e+\pi)=A$ with $A\in\Z$ and $B$ a positive integer,
then for every $j>(3B/2)^2$ every pair of integers $p,d$ satisfying
\[
 1\le d\le H,\qquad |d\pi-p|\le 1/(H+1)
\]
satisfies $q/\gcd(q,B)\mid d$. Such pairs exist for every $j\ge1$.
\end{proposition}
\begin{proof}
Among the $H+1$ fractional parts of $0,\pi,\ldots,H\pi$ on the
unit circle, some circular gap is at most $1/(H+1)$. The difference
of its indices, with sign chosen positive, gives an admissible $d$.
Now the integer $Aqd-B(Pd+pq)$ has absolute value at most
\[
 B\bigl(d|qe-P|+q|d\pi-p|\bigr)
 < B\left(\frac{H}{2jq}+\frac{q}{H+1}\right)
 \le \frac{3B}{2\sqrt j}<1,
\]
using $|qe-P|<1/(2jq)$ and $H\le q\sqrt j<H+1$.
It vanishes. Reduction modulo $q$ gives $q\mid BPd$; since $P/q$
is a reduced convergent, $q\mid Bd$. Cancelling the gcd proves the claim.
\end{proof}
In particular, witnesses with $q/\gcd(q,d)$ unbounded would prove
irrationality, because under rationality these integers all divide
the fixed $B$ at sufficiently large indices. Dirichlet's theorem
provides witnesses but gives no such gcd estimate for $\pi$.
The substitute $x=A/B-e$ satisfies the same restriction; a proof
must use arithmetic information specific to $\pi$. This proposition was independently kernel-checked in \texttt{DVerify.prop\_8\_1}; witness existence for every $j\ge1$ is \texttt{DVerify.dirichlet\_witness}. All ten audited declarations use only Lean base axioms. We make no priority claim for this elementary test.
\subsection{Two primitive witnesses, signed forms, and a series obstruction}
\paragraph{Unconditional information on the sum.}
The finite separation of Proposition~\ref{prop:farey-gap} applies
uniformly to its stated denominator range. No argument here extends
that range to all denominators. The following results concern the
structure of potential proofs.

\begin{proposition}[Two primitive witnesses]\label{prop:two-window}
For any irrational $x$ and integer $H\ge1$, at most two distinct reduced
fractions $p/d$ with $1\le d\le H$ satisfy $|dx-p|\le1/(H+1)$.
If two occur their determinant has absolute value one. In the window of
Proposition~\ref{prop:dirichlet-divisor}, two such fractions for $\pi$
imply $Q_j\mid B$ under its rationality and threshold hypotheses.
\end{proposition}
\begin{proof}
Two distinct solutions on the same side have determinant of absolute
value strictly less than $H/(H+1)$: it is a difference of two
nonnegative products, each at most $H/(H+1)$. On opposite sides the
bound is $2H/(H+1)<2$. Integrality proves both assertions about the
fractions. Their denominators $d,f$ are coprime. Hence
$\operatorname{lcm}(q/\gcd(q,d),q/\gcd(q,f))=q$ by prime valuations.
Each quotient divides $B$ by Proposition~\ref{prop:dirichlet-divisor}.
\end{proof}
Exact rational enclosures (Machin's formula and the exponential series)
reproduce 17 double windows in $1\le j\le60$, and zero for $6-e$.
The count uses the last convergent $a/b$ below or at the cutoff and
$(r+ta)/(s+tb)$, where $r/s$ is its predecessor (initially $1/0$)
and $t=\lfloor(H-s)/b\rfloor$. These are Farey neighbours bracketing
$x$ with denominator sum exceeding $H$. Their positive errors in
absolute value $E,F$ satisfy $(s+tb)E+bF=1$. In their unimodular
basis a positive combination of both violates the denominator cutoff;
opposite coefficients give error at least $E+F\ge1/H$. This exhausts
all possible primitive solutions. For $6-e$, $u=6q-P_j$ is coprime to
$q$, and $|q(6-e)-u|<1/(2jq)\le1/(H+1)$.
For $j\ge3$ a second solution must have
$|ud-qp|<3/(2\sqrt j)<1$ and thus the same slope. For $j=1$ the
denominator is one; for $j=2$ use $q=7,H=9, |7e-19|<0.033$.

\begin{proposition}[Signed forms and a parity test]\label{prop:signed-round}
Put $q=Q_j$, $P=P_j$, $H=\lfloor q\sqrt j\rfloor$, $j\ge2$.
For a convergent $p_k/d_k$ of $\pi$ with $k\equiv j\bmod 2$ and
$d=d_k\le H$, set $s=d_{k+1}$, $g=\gcd(q,d)$,
$U=qd/g$, $V=(dP+qp_k)/g$, $h=\gcd(U,V)$.
Then $L=(U(e+\pi)-V)/h$ has sign $(-1)^j$ and along any sequence
of such indices tending to infinity,
\[
 L\to0\quad\Longleftrightarrow\quad ghs/q\to\infty .
\]
In particular, if the last convergent with denominator at most $H$
has index of the same parity as $j$ infinitely often, $e+\pi$ is irrational.
\end{proposition}
\begin{proof}
The two errors $\delta=qe-P$ and $\eta=d\pi-p_k$ have the stated
common sign by the Hermite identity and the convergent ordering
(starting $p_0/d_0=3/1$). Thus $gh|L|=d|\delta|+q|\eta|$.
Using $|\delta|<1/(2jq)$ and $1/(2s)<|\eta|<1/s$ bounds this sum
below by $q/(2s)$ and above by $1/(2\sqrt j)+q/s$.
For the last convergent $s>H$, hence $s>q\sqrt j$.
Under rationality with denominator $B$, every nonzero $L$ has
$B|L|\ge1$. This contradicts convergence to zero.
\end{proof}
The sign argument settles nonvanishing. Infinite parity agreement for
the specific $\pi$ remains open in this work. As a further conditional
consequence choose the last same-parity $k$ below
$\lfloor jq/B\rfloor$ for sufficiently large $j$. Then
$1\le B|L|<(1/2+Bq/s)/(gh)$, so $s<2Bq$.
Since $a_{k+2}s+d=d_{k+2}>jq/B$ and $d\le s$,
$a_{k+2}>j/(2B^2)-1$. This uses a different cutoff from the parity test.

\begin{proposition}[Obstruction for synchronous Machin truncations]
Let $S_{n,m}$ be $n!$ times the sum of the first $m$ terms in each
series of $16\arctan(1/5)-4\arctan(1/239)$. Along any unbounded
sequence $n\to\infty$ with $|n!\pi-S_{n,m}|=O(1/n)$, the reduced
denominator of $S_{n,m}$ eventually contains a prime $p>n$ exactly once.
\end{proposition}
\begin{proof}
For $v=2m+1$ the signed arctangent remainder at $1/a$ has magnitude
$\int_0^{1/a}t^{v-1}/(1+t^2)\,dt$. Bounding the integrand gives the
combined absolute error strictly between $15n!/(v5^v)$ and
$16n!/(v5^v)$, since $16(25/26)-4(5/239)^v>15$.
Thus $D=2m-1>2n$ eventually (otherwise the lower bound diverges).
By Bertrand's postulate take $D/2<p\le D$; eventually $p>317$.
Only exponent $p$ contributes a denominator divisible by $p$.
Its combined numerator apart from sign and $n!$ is
$16\cdot239^p-4\cdot5^p\equiv3804\not\equiv0\pmod p$.
All other terms have denominators prime to $p$, proving the claim.
\end{proof}
This rules out integrality of the complete truncated sum at the
required precision. It gives no lower bound for its distance to a
prescribed residue. The unrestricted factorial digit problem remains
open. The companion's dimension results concern other points as well
as $\pi$; their paper proofs and literature attribution are separate
from the arithmetic consequences above.



\subsection{An explicit period forced by rationality}
\begin{theorem}\label{thm:explicit-period}
For irrational $x$, let $k(j)$ be the last index of a regular convergent of
$x$ with denominator at most $H_j=\lfloor Q_j\sqrt j\rfloor$, indexing the
integer part by zero. If $e+x=A/B$, with $A\in\mathbb Z$ and $B$ a positive
integer, then for every $j>9B^2/4$ the count $r_j=k(j+1)-k(j)$ is odd, $r_j\ge3$, and
$r_{j+2B^2}=r_j$.
\end{theorem}
\begin{proof}
Write $v_i=(p_i,d_i)$ for the convergent vectors and $m=k(j)$.
Since $d_{m+1}\ge H_j+1$, the convergent error is less than $1/(H_j+1)$.
The integer estimate of Proposition~\ref{prop:dirichlet-divisor}, with $x$ in place of $\pi$,
forces $A Q_j d_m=B(P_j d_m+Q_jp_m)$ whenever $j>(3B/2)^2$.
Thus $v_m=w_j/M_j$, where $w_j=(AQ_j-BP_j,BQ_j)$ and $M_j$ is its
positive content. Since $\gcd(P_j,Q_j)=1$, we have
$M_j\mid B\gcd(AQ_j-BP_j,Q_j)=B\gcd(B,Q_j)\mid B^2$.
The sign of $d_m x-p_m$ is opposite to that of $Q_j e-P_j$.
Alternation gives $m+j$ odd. At the stated threshold $Q_j\ge j>B$ by induction, so $d_m\ge Q_j/B>1$ and a predecessor exists. The ratio of successive locked denominators is $(Q_{j+1}/Q_j)M_j/M_{j+1}>(4j+2)/B^2>1$, so $r=r_j$ is positive and odd.

Put $T=2B^2$. The Hermite determinant has absolute value two, whence
$D=|\det(v_m,v_{m+r})|=T/(M_jM_{j+1})$.
When $A/B$ is reduced, adjacent contents are coprime. Indeed, a prime
dividing both divides $B$ because each content divides $B^2$.
Reduction of $AQ_i-BP_i$ modulo that prime then forces it to divide
both $Q_j$ and $Q_{j+1}$, contradicting adjacent coprimality.
Thus $M_jM_{j+1}\mid B^2$ and $D$ is even. In fact, each reduced
content is odd: for odd $B$ its second coordinate is odd, and
for even $B$ its first coordinate is odd.
The determinant $D$ is invariant under rescaling $(A,B)$ by a positive
integer. The same evenness therefore holds for a nonreduced representation.
Consecutive convergents have determinant of absolute value one, so
$r\ne1$. Since $r$ is positive and odd, $r\ge3$.
Expanding in the unimodular
basis gives $v_{m+r}=C v_m+D v_{m-1}$ with
$C/D=[a_{m+1};\ldots,a_{m+r}]$ by the convergent recurrence.
Primitivity of $v_m$ determines $C\bmod D$ from the two end vectors
modulo $D$. For $D>1$ the fractional part determines the digits after the
integer part up to the two terminal representations, whose lengths differ
by one. Odd length selects exactly one. Consequently $D$ and $C\bmod D$ determine $r$.

For completeness, put $b_{jt}=(2t)!\binom{j+t}{2t}/t!$ for $0\le t\le j$
and zero otherwise. The identities
\[
 P_j=\sum_t b_{jt},\qquad Q_j=(-1)^j\sum_t(-1)^t b_{jt}
\]
follow from the initial pairs $(P_0,Q_0)=(1,1)$, $(P_1,Q_1)=(3,1)$
and $b_{j+1,t}-b_{j-1,t}=(4j+2)b_{j,t-1}$ for $t\ge1$
by factorial cancellation (the constant terms cancel).
Vandermonde's identity expresses $b_{j+T,t}-b_{jt}$ as a sum with factors
$(2t)!\binom{T}{l}/t!$, $1\le l\le\min(2t,T)$.
Every such $l$ divides $(2t)!/t!$: for $l\le t$ the product contains
a multiple of $l$, and for $l>t$ it contains $l$. Since
$l\binom{T}{l}=T\binom{T-1}{l-1}$, each summand is divisible by $T$.
Thus $P_j,Q_j$ repeat modulo $T$ after $T$ steps, as $T$ is even.
Now $M_j$ is the gcd of both coordinates of $w_j$ and $B^2$.
Further, $M_j D$ and $M_{j+1}D$ divide $T$, so the two reduced vectors
modulo $D$ can be recovered before division from $w_j,w_{j+1}\bmod T$.
This proves the asserted period.
\end{proof}
For $x=17/3-e$, the residue calculation in the proof, starting with
the class $j=1\pmod {18}$, gives the repeating word
\[
3,7,3,3,5,3,7,5,7,7,3,7,5,3,5,5,7,5.
\]
Comparison with shifts by the proper divisors of $18$ gives least period
$18$ for this example. For $B=1$ the determinant $D=2$ forces the
fractional part of $C/D$ to be $1/2$, and its odd-length representation
has the two terminal digits $(1,1)$. Thus $r_j=3$ throughout the
locked range for integer translates of $-e$. For the actual $\pi$, it would suffice to prove
that for every $B\ge1$ there exists at least one $j>9B^2/4$ with
$r_{j+2B^2}\ne r_j$. The witness may depend on $B$. Infinitely many windows with $r_j\le2$ would also suffice.
The companion article proves that Lebesgue-almost every variable $x$
has at least logarithmically many windows with $r_j(x)\le1$. The same
article constructs a full-dimensional null set with $r_j(x)\to\infty$,
even when the eventual parity of $k(j)$ is prescribed at every window.
In particular, avoidance of the short-window criterion can hold on a
set of full Hausdorff dimension.
The required infinite-scale information about $\pi$ remains open. The cutoff existence and parity implication for the actual $\pi$ were checked separately in Lean (\texttt{RoundTwoD.irrational\_of\_parity\_agreements}). The explicit period is proved on paper.
\begin{theorem}[Partial quotients in a specified window]\label{thm:digit-shift}
Under the hypotheses of Theorem~\ref{thm:explicit-period}, let
$T=2B^2$, $M_j=\gcd(AQ_j-BP_j,BQ_j)$, and write
$W_j=(a_{k(j)+1},\ldots,a_{k(j+1)})=(b_j,S_j)$, where $S_j$
is the list after the first digit. For every $j$ with $j-1>9B^2/4$,
\[
                W_{j+T}=(b_j+4M_j^2,S_j).
\]
\end{theorem}
\begin{proof}
Put $V_i=v_{k(i)}$, $U_i=v_{k(i)-1}$. Write the change of basis as
\[
 (V_{i+1},U_{i+1})=(V_i,U_i)
 \begin{pmatrix} C_i&E_i\\D_i&F_i\end{pmatrix},
 \qquad D_i=T/(M_iM_{i+1}).
\]
The matrix is the ordered product over $a\in W_i$ of
$\left(\begin{smallmatrix}a&1\\1&0\end{smallmatrix}\right)$.
By Theorem~\ref{thm:explicit-period} it has determinant $-1$.
Its lower row depends only on $S_i$: the lower row after the first
factor is $(1,0)$. The fractional part of $C_i/D_i$ and the odd
length determine $S_i$, as in that theorem, so
$S_{i+T}=S_i$ and $F_{i+T}=F_i$ for $i>9B^2/4$.


Inverting the preceding matrix gives $V_{j-1}=-F_{j-1}V_j+D_{j-1}U_j$.
The unnormalized vectors $M_iV_i=(AQ_i-BP_i,BQ_i)$ obey the Hermite
recurrence. Substitution and comparison of the coefficient of $V_j$
therefore give
\[
 C_j=\frac{(4j+2)M_j-M_{j-1}F_{j-1}}{M_{j+1}} .
\]
The three contents repeat after $T$ steps. Subtracting and dividing by
$D_j=D_{j+T}$ yields
\[
 \frac{C_{j+T}}{D_{j+T}}-\frac{C_j}{D_j}
 =\frac{4T M_j}{M_{j+1}D_j}=4M_j^2 .
\]
These rationals have identical tails, so the difference equals the
difference of the first digits.
\end{proof}

Perron's classical theorem already gives the Hurwitzian structure under
rational linear fractional transformations and preserves the degrees
in the quasi-period (see Stambul~\cite[p.~141]{Stambul}). Here the
windows specify the positions, the onset and the individual increments.
In ordinary digit indices let $K=\sum_{i=j}^{j+T-1}r_i$, independent
of $j$ beyond the threshold by periodicity. Then $k(j+T)=k(j)+K$.
Every sufficiently late digit obeys $a_{n+K}\equiv a_n\pmod4$.
In particular, failure of eventual periodicity of the partial quotients
of $\pi$ modulo four would suffice for the original problem.
No proof of that premise is available here. An irrationality exponent
of two alone is compatible with every model $A/B-e$.

For an exact finite comparison, the Machin alternating-series bounds
give $W_{27}(\pi)=(22,1,2,2,1)$ and $W_{45}(\pi)=(1,2,1,15,1)$.
For $B=3$ the lag is $18$ and the two lengths agree.
The first digit decreases and the tails differ. This excludes that
denominator by the digit test at this pair, already covered by the
paper's earlier finite denominator bound. The endpoint check uses
800 terms for $\arctan(1/5)$ and 300 for $\arctan(1/239)$ and their
respective first omitted terms, stopping the common continued fraction
after its first denominator exceeding the final cutoff.
\begin{corollary}[Tests without a candidate period]\label{cor:digit-statistics}
Suppose $e+x=A/B$. Write $d_n$ for the ordinary convergent denominators
of $x$, and let $K$ be the total number of digits in $T=2B^2$
consecutive windows beyond the threshold of Theorem~\ref{thm:digit-shift}.
With $\alpha=T/K\in\mathbb Q\cap(0,1/3]$,
\[
 \log d_n=\alpha n\log n+
       \alpha(\log(4\alpha)-1)n+O_{A,B}(\log n).
\]
For each sufficiently large integer $d$, the total number $N(d)$ of
occurrences of the digit $d$ is at most $T$. Moreover,
$N(d+4B^4)=N(d)$ for every sufficiently large $d$.
\end{corollary}
\begin{proof}
The periodic lengths give $k(j)=j/\alpha+O_{A,B}(1)$.
The recurrence gives $Q_{j+1}/Q_j=4j+2+O(1/j)$, hence
$\log(Q_{j+1}/Q_j)=\log(4j)+1/(2j)+O(j^{-2})$.
Summing and applying Stirling's formula yields
$\log Q_j=j\log j+(\log4-1)j+O(1)$.
For $k(j)\le n<k(j+1)$ the bounding denominators are
$BQ_j/M_j$ and $BQ_{j+1}/M_{j+1}$. Their logarithms differ by
$O_{A,B}(\log j)$ and $j=\alpha n+O_{A,B}(1)$.
Substitution proves the formula.

Choose a complete period of windows after the threshold. Let $b_i$
and $M_i$ denote the first digit and content in each of these $T$
windows. Beyond the largest digit in the discarded prefix, the tails
of these windows and their first entries, the only occurrences come
from the $T$ progressions $b_i+4M_i^2t$, $t\ge0$. Thus
\[
 N(d)=\sum_{i=1}^T {\bf1}_{d\equiv b_i\pmod {4M_i^2}}.
\]
Each modulus divides $4B^4$, since $M_i\mid B^2$.
\end{proof}
In particular either of the following would suffice for $e+\pi$ to be
irrational: (i) along some subsequence $\log d_n/(n\log n)\to0$ for
$\pi$, or (ii) for every $D,L\ge1$ some digit $d>D$ occurs at least
$L$ times in its expansion. In (ii) each individual digit is allowed
to occur only finitely often. These hypotheses on $\pi$ are unproved
here. The companion paper also determines the dimension of the set
that passes the weaker eventual-periodicity test modulo four.
\subsection{The converse to denominator locking}
Write $p_j/d_j$ for the last reduced convergent of an irrational $x$
at the cutoff $H_j$ (the subscript here counts windows). Put
\[
 g_j=\gcd(Q_j,d_j),\quad u_j=Q_j/g_j,\quad v_j=d_j/g_j,
 \qquad Y_j=Q_jp_j/d_j .
\]
Thus $u_j/v_j=Q_j/d_j$ in lowest terms. Both ratios in the following
classification are observable from ordinary convergents.
\begin{theorem}[Classification by either denominator quotient]\label{thm:gcd-converse}
For irrational $x$, each of the following conditions is equivalent to
$x\in\mathbb Q+\mathbb Q e$:
\begin{enumerate}
\item $u_j$ is eventually bounded;
\item $v_j$ is eventually bounded;
\item $\displaystyle\limsup_{j\to\infty}\frac{\log u_j}{\log Q_j}<1/3$.
\end{enumerate}
In particular every $x$ outside that affine span has limsup at least
$1/3$. No optimality assertion about this threshold is needed.
\end{theorem}
\begin{proof}
First suppose $x=(A+Ce)/B$ with integers $B>0$ and $C\ne0$.
Consider $(AQ_j+CP_j)/(BQ_j)$ and call the positive common content
$M_j$. Since $\gcd(P_j,Q_j)=1$,
$M_j\mid B\gcd(C,Q_j)\mid B|C|$. Its reduced denominator
$d=BQ_j/M_j$ is comparable to $Q_j$, uniformly in $j$ for fixed $B,C$.
The error before reduction as an approximation to $x$ is
$|C|\,|Q_je-P_j|/(BQ_j)=O(1/(jQ_j^2))$. Legendre's criterion
therefore makes it a convergent for large $j$. Its successor has
denominator $d'>1/|dx-p|-d\gg jQ_j$ by the convergent error formula.
Eventually $d\le H_j<d'$, so this fraction is exactly $p_j/d_j$.
Now $Q_j/d_j=M_j/B$, giving both bounds. Bounded $u_j$ implies (3).

For the converse put $E_j=Q_j x-Y_j$. The successor at the cutoff
has denominator at least $H_j+1$, hence
\[
 |E_j|<\frac{u_j}{v_j(H_j+1)},\qquad
 v_j\le\sqrt j\,u_j,\qquad v_jY_j\in\mathbb Z .
\]
Let $R_j=Y_{j+1}-(4j+2)Y_j-Y_{j-1}$. The product
$V_j=v_{j-1}v_jv_{j+1}$ clears its denominators. Because $Q_j$ obeys
the Hermite recurrence, replacing each $Y_i$ by $-E_i$ gives the same
$R_j$. Under (3) choose $0<\theta<1/3$ with $u_i\le Q_i^\theta$
eventually. The displayed estimates yield
\[
 |V_j R_j|\ll j^2
 \frac{u_{j-1}u_ju_{j+1}}{H_{j-1}+1}
 \ll j^{2+\theta}Q_{j-1}^{\theta-1}Q_j^{2\theta}.
\]
Here $Q_{j+1}/Q_j\ll j$, and likewise $Q_j/Q_{j-1}\ll j$.
Thus the bound is a polynomial in $j$ times $Q_{j-1}^{3\theta-1}$,
which tends to zero since $Q_j$ grows faster than exponentially.
The integer $V_jR_j$ therefore vanishes eventually.

There is a separate argument under (2). Choose an integer $L$ divisible
by every eventual $v_j$. Then $d_j\mid LQ_j$ and $LY_j$ is integral.
The $d_j$ are nondecreasing and tend to infinity. Equality
$d_j=d_{j+1}$ would give $d_j\mid L$ since consecutive $Q$'s are
coprime. Consequently $d_j$ eventually increases at each step and
$d_j\ge j-O(1)$. Now $|E_j|<Q_j/(d_j(H_j+1))<1/(d_j\sqrt j)$,
so $LR_j$ is an integer tending to zero, again giving $R_j=0$ eventually.

In either converse, two successive rational values of $Y$ determine
rational coefficients $c,b$ with $Y_j=cP_j+bQ_j$ on the whole tail:
the two fundamental solutions have determinant of absolute value two.
Division by $Q_j$ and passage to the limit give $x=ce+b$.
\end{proof}
\begin{corollary}[Recovering the slope]\label{cor:recover-slope}
Under any of the three equivalent conditions, the entirely rational
quantity
\[
 s_j=\frac{p_{j+1}/d_{j+1}-p_j/d_j}
              {P_{j+1}/Q_{j+1}-P_j/Q_j}
\]
eventually equals the coefficient $c$ in $x=b+ce$. Consequently
$e+x\in\mathbb Q$ exactly when a denominator quotient is eventually
bounded and $s_j=-1$ for infinitely many $j$.
\end{corollary}
\begin{proof}
The locked fractions in the forward proof satisfy
$p_j/d_j=b+cP_j/Q_j$ eventually. Consecutive selected Hermite
fractions are distinct by their determinant. The slope is unique
since $e$ is irrational.
\end{proof}
For instance $x=6-2e$ has $d_j=Q_j$ eventually because $Q_j$ is odd
and $\gcd(6Q_j-2P_j,Q_j)=1$. Both quotients are one, while $s_j=-2$.
Thus a proof for $\pi$ that either quotient is unbounded would in
fact establish rational linear independence of $1,e,\pi$. Excluding
only slope $-1$ has the narrower conclusion sought for the sum.

\subsection{Window periods for rational affine images}
\begin{theorem}[Affine window period and digit increment]\label{thm:affine-window}
Let $x=(A+C e)/B$ with $A,C\in\Z$, $C\ne0$, and $B\in\N_{>0}$,
and set
\[
 T=2B|C|,\qquad L=\max(B,|C|)^2,\qquad
 G_j=\gcd(AQ_j+C P_j,BQ_j)>0 .
\]
For all $j>L$, the window length $r_j$ is positive and odd,
$r_{j+T}=r_j$, and
\[
 r_j=1\quad\Longleftrightarrow\quad G_jG_{j+1}=T.
\]
Write the window digit list as $W_j=(b_j,S_j)$, with first digit
$b_j$ and remaining list $S_j$. For $j-1>L$,
\[
 W_{j+T}=(b_j+4G_j^2,S_j).
\]
When $\gcd(A,B,C)=1$, an eventual one-digit window can occur only
if $B$ is even and both $A,C$ are odd.
\end{theorem}
\begin{proof}
Write $c=|C|$, $q=Q_j$, and reduce the fraction
$(Aq+CP_j)/(Bq)$ to $p/d$. Because $P_j$ and $Q_j$ are coprime,
$G_j\mid B\gcd(C,Q_j)$ and $G_j\mid Bc$. The Hermite error bound gives
\[
 d|dx-p|=\frac{Bc Q_j}{G_j^2}|Q_j e-P_j|
            <\frac{Bc}{2jG_j^2}.
\]
For $j>L$ we have $j>Bc$ and Legendre's criterion applies.
Also $d\le Bq<q\sqrt j$, hence $d\le H_j$. If $d'$ is the next
convergent denominator, the error formula gives
$d'>1/|dx-p|-d>(2j/c-B)q>q\sqrt j$. Thus $p/d$ is exactly the
terminal convergent at the cutoff. Its denominator is at least
$Q_j/c>1$. The terminal denominators are strictly increasing since
their successive ratio exceeds $(4j+2)/(Bc)>1$.

Set $V_j=(p,d)$ at that cutoff. The sign of its error equals that of
$C(Q_j e-P_j)$. Alternation therefore gives a fixed parity offset
between $k(j)$ and $j$, and $r_j$ is odd. The determinant is
\[
 D_j=|\det(V_j,V_{j+1})|=\frac{T}{G_jG_{j+1}}.
\]
In the basis consisting of $V_j$ and its immediate predecessor, the
coefficient ratio for $V_{j+1}$ is a finite continued fraction
$C_j/D_j$ with digits $W_j$. If $D_j=1$, coprimality and the odd length
give the one-digit expansion. If $D_j>1$, the fractional part of
$C_j/D_j$ gives the suffix up to the two terminal representations,
and odd length selects one. Thus $r_j=1$ precisely when $D_j=1$.

The binomial congruence in Theorem~\ref{thm:explicit-period} holds
for every even modulus $T$ with shift $T$. In particular
$P_{j+T}\equiv P_j$ and $Q_{j+T}\equiv Q_j\pmod T$.
Since $G_j\mid Bc$, it is recovered by taking the gcd of both
unnormalized coordinates and $Bc$. Hence $G_j$ and $D_j$ repeat
after $T$ steps. Both $G_jD_j$ and $G_{j+1}D_j$ divide $T$, so the
normalized vectors modulo $D_j$ are recovered from those coordinates
modulo $T$. Primitivity of $V_j$ determines $C_j\bmod D_j$.
Therefore both $S_j$ and $r_j$ are periodic, including the empty
suffix when $D_j=1$.

For the first-digit increment, let $F_{j-1}$ denote the lower-right
entry of the preceding window matrix, as in the proof of
Theorem~\ref{thm:digit-shift}. It is determined by the preceding
suffix and equals zero for an empty suffix. The unnormalized
vectors $G_iV_i=(AQ_i+CP_i,BQ_i)$ obey the Hermite recurrence.
Inverting the preceding window matrix, whose determinant is $-1$,
therefore gives
\[
 C_j=\frac{(4j+2)G_j-G_{j-1}F_{j-1}}{G_{j+1}}.
\]
The shift by $T$ increases $C_j/D_j$ by
$4T G_j/(G_{j+1}D_j)=4G_j^2$. Its fractional part stays fixed, so
the same integer is the first-digit increment.

Finally, if $B$ is odd then each $G_j$ is odd since $Q_j$ is odd.
If $B$ is even and $A,C$ have opposite parity then the first
unnormalized coordinate is odd, with the same conclusion.
For a primitive triple these cases cover everything except even $B$
and odd $A,C$. If the contents are odd then $D_j$ is even,
ruling out a one-digit window.
\end{proof}

\begin{remark}[Lean assembly]\label{rem:affine-assembly}
The algebra of the two proofs above has been assembled in
\texttt{EpiFormal/AffineWindowAssembly.lean} (namespace \texttt{AffineWindow}):
content arithmetic, evenness of $D_j$, oddness and $r_j\ge3$,
period transfer from $G$-periodicity, and the first-digit increment
identity $4T G_j/(G_{j+1}D_j)=4G_j^2$ (seven lemmas, each with explicit
interface hypotheses for the locked denominators, the Hermite recurrence,
the mod-$T$ congruence, and the digit-determination rule).
The file elaborates with zero \texttt{sorry}; the kernel audit is recorded in
\texttt{verification/\allowbreak audit\_\allowbreak final\_\allowbreak 20261122.log}
(937 audited declarations project-wide across 134 Lean files, \texttt{sorryAx}=0, zero project-added axioms;
the seven \texttt{AffineWindow} declarations use only Lean base axioms, three of them strictly weaker including one axiom-free).
This remark does not affect the paper proofs, which are self-contained.
\end{remark}

For example $x=(1+e)/2$ has $B=2$, $C=1$, $G_j=2$ and $D_j=1$.
Its sufficiently late windows each contain one digit. By contrast,
for a primitive representation of $x=A/B-e$, odd $B$ gives odd
second coordinate, and even $B$ gives odd first coordinate. This
explains the stronger lower bound of Theorem~\ref{thm:explicit-period}
within the general affine family. Regarding prior literature,
earlier work of Davis~\cite{Davis1945} (as cross-referenced in
Davis~\cite{Davis1978} and Walters~\cite{Walters1968}) obtained regular
continued fractions and convergents for $e^{1/t}$ and $e^{2/t}$ via
diagonal Pad\'e approximants; for the homogeneous case $A=0$ (namely
$(m/n)e$), Matthews and Walters~\cite[Theorem~1(b)]{MatthewsWalters1970}
already established an explicit period $d=mn$, an explicit threshold,
and a squared-$\gcd$ arithmetic increment $p_t^2(2x-1)$ (increment
$4p_t^2$ per $2mn$ steps). Meanwhile, Perron~\cite[Band~I, \S33, beginning on p.~119]{Perron1954}
(as discussed in Hetyei~\cite[\S1.2, p.~4]{Hetyei2014}), Stambul~\cite{Stambul}, and
Hetyei~\cite{Hetyei2014} treat qualitative Hurwitzian closure and
repeated-constant Hurwitzian patterns without tracking how rational
M\"obius transformations affect the quasi-period length. In the modern
literature, Komatsu~\cite{Komatsu1999JTNB,Komatsu2003Acta,Komatsu2003Tsukuba,Komatsu2007Czech,Komatsu2008Sarajevo}
investigates ${}_0F_1$ hypergeometric closed forms, leaping convergents,
and inhomogeneous constants $M(\theta,\phi)$ for Hurwitz and Tasoev
continued fractions. Komatsu~\cite[Theorems~1--4, pp.~49--50]{Komatsu2005Monatsh}
gives series expressions for prescribed period-three $e$-type families and
their Tasoev analogues. Komatsu~\cite{Komatsu2009JNT}, by its abstract,
reduces periods to~$1$ using non-simple $N$-continued fractions.
Kalinin and Komatsu~\cite{KalininKomatsu2026} analyze factorial-series
error terms for $e^{1/s}$ and $\sqrt{v/u}\tanh(1/\sqrt{uv})$.
Havens, Barbero, Cerruti, and Murru~\cite{HBCM2020} resolve $|\det\sigma|=2$ linear
fractional transformations while explicitly noting that $|\det\sigma|>2$
requires different methods. Within a negative search of CrossRef,
OpenAlex, and arXiv records, no prior precedent was found for the
unified treatment of the non-homogeneous affine shift $A\ne0$ (notably
$A/B-e$), the explicit Hermite-cutoff window length criterion
$r_j=1\iff G_jG_{j+1}=T$, and the sharp five-window characterization
below (this reports a negative search, not a proof of non-existence).


\begin{corollary}[Runs of single-digit affine windows]\label{cor:affine-single}
Assume $\gcd(A,B,C)=1$ in Theorem~\ref{thm:affine-window}.
The following conditions are equivalent:
\begin{enumerate}
\item[(i)] $r_j=1$ eventually.
\item[(ii)] Five consecutive windows starting at some $j>L$ have $r_j=1$.
\item[(iii)] $|C|=1$, $B\in\{2,4\}$, and $A$ is odd.
\end{enumerate}
The five-window length in (ii) is sharp.
\end{corollary}
\begin{proof}
Condition (i) implies (ii). Suppose (ii) holds. For the six
terminal indices involved, every consecutive pair of contents
has product $2B|C|$ by Theorem~\ref{thm:affine-window}.
Suppose an odd prime $\ell$ divides $B|C|$. If $\ell\mid C$,
then $\ell\mid G_i$ is equivalent to $\ell\mid Q_i$:
when $\ell\nmid B$ this follows from the second coordinate, and
when $\ell\mid B$ it follows from $\ell\nmid A$ and the first.
Two consecutive $Q_i$ cannot both be divisible by $\ell$.
If $\ell\nmid C$, then $\ell\mid B$, and
$\ell\mid G_i$ is equivalent to $\ell\mid U_i$ for
$U_i=AQ_i+CP_i$. Consecutive $U_i$ cannot both vanish modulo $\ell$
because $U_{i+1}Q_i-U_iQ_{i+1}=\pm2C$.
Thus the six terms of the corresponding sequence $Q_i$ or $U_i$
alternate between zero and nonzero residues modulo $\ell$.
Among the four interior indices there are two nonzero terms
at indices $i$ and $i+2$, each between zero residues.
The Hermite recurrence forces $\ell$ to divide both $4i+2$
and $4(i+2)+2$, hence to divide eight, a contradiction.

Therefore $B|C|$ is a power of two. The presence of a one-digit
window gives odd $A,C$ and even $B$ by
Theorem~\ref{thm:affine-window}. Write $C=\pm1$, $B=2^b$, $b\ge1$.
Since $Q_i$ is odd and $U_{i+1}Q_i-U_iQ_{i+1}=\pm2C$, exactly one
of each pair of $U$'s is congruent to two modulo four, and the other
is divisible by four. Here $G_i=\gcd(U_i,2^b)$.
If $b\ge3$, the product $G_iG_{i+1}=2^{b+1}$ forces the term
divisible by four in each pair to vanish modulo eight.
The six terms then alternate between zero modulo eight and
two modulo four. At an interior index $i$ of the latter kind,
the recurrence forces $8\mid(4i+2)U_i$, whereas this product
is congruent to four modulo eight. This contradiction gives $b\le2$
and establishes (iii).

Conversely, for $b=1$ the two contents are both two, and for $b=2$
they are two and four in alternating order. Their product is
$2B|C|$ in both cases. Hence (iii) implies (i).

For sharpness, take $x=(1+e)/6$. Its content is $G_i=6$ when
$3\mid P_i+Q_i$, and $G_i=2$ otherwise. The recurrence gives
$(P_i+Q_i)\bmod3=(2,0,2,0,2)$ for $i=2,\ldots,6$.
The content has period $12$, so $G_iG_{i+1}=12$ for four
consecutive windows starting at $i\equiv2\pmod{12}$ beyond the
threshold. This model fails (iii).
\end{proof}

\paragraph{Lean verification of the recurrence steps.} In \texttt{AffineWindow\allowbreak{}Audit.lean}, \texttt{selected\allowbreak{}Numerator\_\allowbreak{}no\_\allowbreak{}five}
checks the odd-prime obstruction for $AQ_j+CP_j$, including $Q_j$ when
$(A,C)=(1,0)$. The declaration \texttt{selected\allowbreak{}Numerator\_\allowbreak{}no\_\allowbreak{}eight\_\allowbreak{}middle}
checks the incompatible middle residue modulo eight.
The declaration \texttt{sixthModel\_\allowbreak{}four\_\allowbreak{}products} checks the four
content products for the sharpness example. The first two checks depend
only on Lean's standard axioms. The third has an empty axiom report.
The lock, its determinant interpretation, and the passage from the
checked content products to the full periodic window statement are
proved above.

For $C=-B$, Theorem~\ref{thm:affine-window} improves the
threshold in Theorems~\ref{thm:explicit-period} and~\ref{thm:digit-shift}
to $j>B^2$ and $j-1>B^2$, respectively. After reducing $A/B$
the contents are odd by the argument in Theorem~\ref{thm:explicit-period},
so $r_j\ge3$ already at this smaller threshold.

\begin{proposition}[Prime-power reductions for rational translates]\label{prop:rational-gcd-valuations}
Assume $B\ge1$ and $\gcd(A,B)=1$. For a Hermite pair $(P_j,Q_j)$ of $e$ write
$G_j=\gcd(AQ_j-BP_j,BQ_j)$ and $g=\gcd(B,Q_j)$. Then
\[
G_j=g\gcd\bigl(A(Q_j/g)-(B/g)P_j,g\bigr).
\]
For $p^u\parallel B$ and $s=v_p(Q_j)$, the exponent of $p$ in $G_j$ is
\[
v_p(G_j)=
\begin{cases}
\min(u,s), & s\ne u,\\
u+\min\{u,v_p(A(Q_j/p^u)-(B/p^u)P_j)\}, & s=u .
\end{cases}
\]
Here $v_p(0)=\infty$. When $s=u$, an extra exponent of at least $k$,
$1\le k\le u$, occurs precisely for
\[
 A\equiv (B/p^u)P_j\,(Q_j/p^u)^{-1}\pmod{p^k}.
\]
Moreover $G_jG_{j+1}\mid B^2$.
\end{proposition}
\begin{proof}
Put $b=B/g$, $q=Q_j/g$ and $L=Aq-bP_j$. The assumptions
$\gcd(b,q)=\gcd(A,b)=\gcd(P_j,q)=1$ imply $\gcd(L,bq)=1$.
Consequently $G_j=g\gcd(L,g b q)=g\gcd(L,g)$. For a prime $p\mid B$,
if $s<u$ the first summand of $AQ_j-BP_j$ has valuation $s$ while the
second has valuation at least $u$. If $s>u$, then $p\nmid P_j$ and the
second has valuation $u$ while the first has valuation $s$. At equality
$s=u$ divide the two summands by $p^u$ and compare with the denominator,
whose valuation is $2u$. The resulting divisibility condition is the
stated congruence because $Q_j/p^u$ is a unit modulo $p^k$.
Finally $G_j\mid g^2$ by the first formula. All $Q_j$ are odd and
$|P_{j+1}Q_j-P_jQ_{j+1}|=2$, so $\gcd(Q_j,Q_{j+1})=1$.
Thus $\gcd(G_j,G_{j+1})=1$ and the product divides $B^2$.
\end{proof}
For $x=A/B-e$ the window determinant $D_j=2B^2/(G_jG_{j+1})$
can therefore be assembled prime by prime. The extra reduction at a fixed
$p$ is governed by the single residue class above. Conditions for distinct
primes combine by the Chinese remainder theorem.

\section{The bounded golden-orbit limit spectrum}\label{sec:golden}

This chapter evaluates, by a finite exact algorithm, the bounded
golden-ratio orbit limit spectrum for every reduced rational
affine image of $e$, and establishes the complete Case~1 and
Case~2 classifications of the exact constants $c_\varphi$, all
extremal families, and the second smallest limit points.
\subsection{Bounded digit blocks and the golden-ratio orbit}
Let $F_0=0$, $F_1=1$ and $F_{n+2}=F_{n+1}+F_n$.
\begin{corollary}[Two digit tests]\label{cor:bounded-digit-block}
Let $x=A/B-e$ with $B\ge1$ and $A/B$ reduced. Fix an integer $M\ge1$.
Choose an integer $J>B^2$ satisfying $2J+1\ge(M+2)B^2$.
Every run of $\ell$ consecutive partial quotients at most $M$ starting
after digit $k(J)$ satisfies $F_{\ell+1}\le2B^2$. Also, among any three
consecutive digits starting after $k(B^2+1)$, at most one exceeds $2B^2$.
\end{corollary}
\begin{proof}
Use the notation of Theorems~\ref{thm:explicit-period}
and~\ref{thm:affine-window}. For $j>B^2$ put $m=k(j)$ and $r=r_j$.
The coefficient of $v_{m-1}$ in $v_{m+r}=C_jv_m+D_jv_{m-1}$ is
the continuant of the suffix $S_j$, with $D_j=2B^2/(M_jM_{j+1})$.
Indeed, multiplying the digit matrices leaves the row $(1,0)$ after
the first digit and gives the suffix continuant as its final first entry.
For a word of $t$ positive digits the continuant is at least $F_{t+1}$
and at least every digit in the word. The first bound follows by
minimizing the positive recurrence at the all-one word. For the second,
adjoining a positive digit at either end never decreases the continuant.
In particular, $F_{r_j}\le D_j\le2B^2$.

The denominator coordinate of the vector identity gives
\[
 \frac{C_j}{D_j}
 =\frac{Q_{j+1}}{Q_j}\frac{M_j^2}{2B^2}-\frac{d_{m-1}}{d_m}
 >\frac{2j+1}{B^2}-1 .
\]
As $r_j\ge3$, the suffix has length at least two, and the first digit
is $\lfloor C_j/D_j\rfloor$. For $j\ge J$ it is therefore greater than $M$.
A run as in the statement lies inside one suffix, giving
$F_{\ell+1}\le F_{r_j}$. Only a first digit can exceed $2B^2$, and
successive first digits are at least three positions apart.
\end{proof}
The finite-word inequalities are checked in \texttt{DigitTail\allowbreak{}Bound.lean}.
Their application to the window determinant uses the paper proof above.
Thus it suffices either to find arbitrarily long bounded digit blocks
for one fixed $M$ at $\pi$, or to prove that
$\min(a_n(\pi),a_{n+h}(\pi))$ is unbounded for one $h\in\{1,2\}$.
The all-one test admits a formulation in terms of quadratic approximation.

Write $\varphi=(1+\sqrt5)/2$, let $E_\varphi$ be the modular orbit
of $\varphi$ and its conjugate as in \cite[(1.1)--(1.2)]{BugeaudQuad},
and put $h(\beta)=2/|\beta-\beta'|$. For $x\notin\Q\cup E_\varphi$
set $c_\varphi(x)=\liminf_{\beta\in E_\varphi,\,h(\beta)\to\infty}
h(\beta)|x-\beta|$.
\begin{corollary}[A lower bound for the golden-ratio orbit]\label{cor:golden-bound}
For $x=(A+Ce)/B$ as in Theorem~\ref{thm:affine-window}, put
$T=2B|C|$. Then
\[
 c_\varphi(x)\ge\frac{1}{10368T^2},
 \qquad c_\varphi(A/B-e)\ge\frac{1}{41472B^4}.
\]
\end{corollary}
\begin{proof}
The transcendence of $x$ ensures that it lies outside $\Q\cup E_\varphi$.
For each of the $T$ window residue classes, the head grows by at least
four per period by Theorem~\ref{thm:affine-window}. Eventually every
head exceeds one. A run of $\ell\ge1$ ones then lies inside a suffix
with continuant $D_j=T/(G_jG_{j+1})\le T$. The finite-word bound in
the preceding proof gives $F_{\ell+1}\le T$, also valid for $\ell=0$.
We use the estimates in \cite[Lemmas~2.1--2.2 and Corollary~2.3]{BugeaudQuad}.
It suffices to consider sequences of orbit points $\beta$ for which
$h(\beta)\to\infty$ and $h(\beta)|x-\beta|$ stays bounded. Then
$\beta\to x$, and the common prefix of their continued fractions grows
without bound, since $x$ lies in the interior of each of its cylinders.
Remove the common integer part when applying Lemma~2.2.

Write the expansion of $\beta$ with its shortest preperiod followed
by ones. If the first mismatch with $x$ occurs within the preperiod,
Corollary~2.3 of the cited paper with period alphabet bound $2$ gives
$h(\beta)|x-\beta|\ge1/(96\cdot2^3\cdot3^2)=1/6912$.
Otherwise $\beta=[a_0;\ldots,a_s,\overline{1}]$ with the entire
preperiod shared and $a_s\ne1$. Here $s$ tends to infinity on this
subsequence, because bounded $s$ allows only finitely many such
points. Let $\ell\ge0$ be the number of following ones shared with $x$
and put $m=s+\ell+1$, the first mismatch. For sufficiently large $s$
we have $F_{\ell+1}\le T$. Write $q_i$ for the denominators of $\beta$.
The recurrence on its all-one tail gives
\[
 q_m=F_{\ell+2}q_s+F_{\ell+1}q_{s-1}
     \le F_{\ell+3}q_s\le3F_{\ell+1}q_s\le3Tq_s .
\]
Lemma~2.1 gives $h(\beta)\ge q_s^2/32$ since $s\ge3$ eventually.
Lemma~2.2, equation (2.4), gives $|x-\beta|\ge 1/(36q_m^2)$
since the two digits of $\beta$ following the mismatch are ones.
Multiplication gives the stated bound. Taking $C=-B$ gives the
second inequality.
\end{proof}
For irrationals whose regular continued fraction eventually consists of
triples $(A_k,1,1)$ with $A_k\to\infty$, the golden-ratio orbit constant
$c_\varphi(x)$ and its second smallest limit point can be determined
exactly. In the quadratic Lagrange spectrum literature
\cite{ParkkonenPaulin2011,BugeaudQuad,Pejkovic2016,Lin2018},
Lin~\cite[Lemma~4.2]{Lin2018} (\texttt{arXiv:1703.04847v2}, p.~14)
treats the regime $\liminf_{r\to\infty} a_r=\infty$ and proves
$c_\varphi(\xi)=3/\sqrt5-1=K_\varphi$, whereas the $(A_k,1,1)$ regime
here has $\liminf_{r\to\infty} a_r=1$ and yields the smaller exact value
$c_\varphi(x)=1-2/\sqrt5$ together with the second limit point
$4/\sqrt5-1$, with no overlap. (A search of 2018--2026 literature
found no precedent for this $(A_k,1,1)$ evaluation; this is a negative
search and does not prove non-existence of prior work.) This upgrades
the earlier upper bound on $6-e$ to an exact evaluation across the
entire $\mathrm{GL}_2(\Z)$-orbits of $e^{1/m}$ ($m\ge1$). We first
record an algebraic rigidity lemma in $\Z[\varphi]$.

\begin{lemma}[{$\Z[\varphi]$ norm parity rigidity}]\label{lem:norm-parity-rigidity}
For every $\eta\in E_\varphi$, there exists $m\in\Z$ such that
\[
 h(\eta)\frac{|2-\eta|\,|\eta'|}{2}=\left|\frac{2m}{\sqrt5}-1\right|.
\]
Consequently, $h(\eta)|2-\eta|\,|\eta'|/2 = 1-2/\sqrt5$ if and only if
$\eta\in\{\varphi,\varphi+1\}$, and for every
$\eta\in E_\varphi\setminus\{\varphi,\varphi+1\}$ one has
\[
 h(\eta)\frac{|2-\eta|\,|\eta'|}{2}\ge \frac{4}{\sqrt5}-1.
\]
\end{lemma}
\begin{proof}
Write $\eta=(a\varphi+b)/(c\varphi+d)$ with $a,b,c,d\in\Z$ and
$\varepsilon:=ad-bc\in\{\pm1\}$. Set
$\alpha=(2c-a)\varphi+(2d-b)\in\Z[\varphi]$ and
$\gamma=a\varphi+b\in\Z[\varphi]$. Since
$\eta-\eta'=\varepsilon\sqrt5/((c\varphi+d)(c\varphi'+d))$, we have
$h(\eta)=2|(c\varphi+d)(c\varphi'+d)|/\sqrt5$, while
$2-\eta=\alpha/(c\varphi+d)$ and $\eta'=\gamma'/(c\varphi'+d)$, so the
denominators cancel and
\[
 h(\eta)\frac{|2-\eta|\,|\eta'|}{2}=\frac{|\alpha\gamma'|}{\sqrt5}.
\]
Expanding $\alpha\gamma'$ via $\varphi\varphi'=-1$ and
$\varphi=(1+\sqrt5)/2$, $\varphi'=(1-\sqrt5)/2$, the coefficient of
$\sqrt5$ is $((2c-a)b-(2d-b)a)/2=cb-ad=-\varepsilon$, and the rational
part is
$m_0=-(2c-a)a+(2d-b)b+((2c-a)b+(2d-b)a)/2=-2ac+a^2+2bd-b^2+bc+ad-ab\in\Z$.
Thus $\alpha\gamma'=m_0-\varepsilon\sqrt5$. For any $p\varphi+q\in\Z[\varphi]$,
its field norm $N(p\varphi+q)=-p^2+pq+q^2$ is even if and only if
$p\equiv q\equiv0\pmod2$. Since $|ad-bc|=1$, we have $\gcd(a,b)=1$, so
$(a,b)\not\equiv(0,0)\pmod2$ and hence
$(2c-a,2d-b)\equiv(a,b)\not\equiv(0,0)\pmod2$. Thus
$N(\alpha)\equiv N(\gamma)\equiv1\pmod2$, giving
\[
 m_0^2-5=N(\alpha\gamma')=N(\alpha)N(\gamma)\equiv1\pmod2,
\]
so $m_0=2k_0\in2\Z$. Setting $m=\varepsilon k_0\in\Z$ yields the displayed
identity.

As $m$ ranges over $\Z$, the values of $|2m/\sqrt5-1|$ are $1-2/\sqrt5$
at $m=1$, $4/\sqrt5-1$ at $m=2$, $1$ at $m=0$, and at least
$6/\sqrt5-1>1$ otherwise. Equality at $m=1$ forces $k_0=\varepsilon$, so
$\alpha\gamma'=\varepsilon(2-\sqrt5)$. Taking norms gives
$N(\alpha)N(\gamma)=-1$, so $N(\gamma)\in\{\pm1\}$. Since
$\alpha+\gamma=2(c\varphi+d)$, we obtain
$\alpha/\gamma=\alpha\gamma'/N(\gamma)=\sigma(2-\sqrt5)$ with
$\sigma=\varepsilon N(\gamma)\in\{\pm1\}$, whence
$\eta=2/(1+\alpha/\gamma)$ equals $\varphi+1$ (for $\sigma=+1$) or
$\varphi$ (for $\sigma=-1$). Direct evaluation confirms both values.
\end{proof}

\begin{theorem}[Exact golden-ratio orbit constant and second limit point]\label{thm:golden-exact-second}
Let $x=[a_0;a_1,a_2,\ldots]$ be an irrational number for which there
exists $s_0\ge1$ such that, with $s_k:=s_0+3k$ ($k\ge0$),
\[
 a_{s_k}=A_k\to\infty,\qquad a_{s_k+1}=1,\qquad a_{s_k+2}=1.
\]
Define the two families of orbit points
$\beta_k^{(1)}:=[a_0;\ldots,a_{s_k},\overline1]$ and
$\beta_k^{(2)}:=[a_0;\ldots,a_{s_k},2,\overline1]$. Then:
\begin{enumerate}
\item $c_\varphi(x)=1-2/\sqrt5$, attained with limit $1-2/\sqrt5$ along
      both families $\beta_k^{(1)}$ and $\beta_k^{(2)}$.
\item Every sequence $(\beta_n)\subset E_\varphi$ with
      $h(\beta_n)\to\infty$ that is eventually disjoint from
      $\{\beta_k^{(1)},\beta_k^{(2)}:k\ge0\}$ satisfies
      \[
       \liminf_{n\to\infty} h(\beta_n)|x-\beta_n|\ge \frac{4}{\sqrt5}-1.
      \]
\item The second smallest limit point $4/\sqrt5-1$ is sharply attained,
      for example (taking $k$ sufficiently large so that $A_k\ge2$) by each
      of the families
      $[a_0;\ldots,a_{s_k},2,2,\overline1]$,
      $[a_0;\ldots,a_{s_k-1},A_k-1,\overline1]$,
      $[a_0;\ldots,a_{s_k-1},A_k+1,2,\overline1]$, and
      $[a_0;\ldots,a_{s_k-1},2,\overline1]$.
\end{enumerate}
\end{theorem}
\begin{proof}
For any $m\ge3$, write $x=M_m(T_m)$ and $\beta=M_m(\eta)$ using the
convergent matrix $M_m=\bigl(\begin{smallmatrix}p_{m-1}&p_{m-2}\\
q_{m-1}&q_{m-2}\end{smallmatrix}\bigr)$ of $x$, with
$T_m=[a_m;a_{m+1},\ldots]$ and $z_{m-1}=q_{m-2}/q_{m-1}\in(0,1)$.
By the M\"obius error identity \cite[(3.4)]{BugeaudQuad},
\begin{equation}\label{eq:bugeaud-34}
 h(\beta)|x-\beta|
 =h(\eta)\frac{|T_m-\eta|\,|\eta'+z_{m-1}|}{T_m+z_{m-1}},
 \qquad
 h(\beta)=h(\eta)\,q_{m-1}^2\,|\eta+z_{m-1}|\,|\eta'+z_{m-1}|.
\end{equation}
At $m=s_k+1$, we have $T_m=[1;1,A_{k+1},1,1,\ldots]\to2$ and
$z_{m-1}=[0;A_k,\ldots]\to0$. Taking $\eta=\varphi$ (for $\beta_k^{(1)}$)
and $\eta=\varphi+1$ (for $\beta_k^{(2)}$) gives $|\eta\eta'|=1$, so
$h(\beta_k^{(i)})\sim (2/\sqrt5)q_{s_k}^2\to\infty$ and
$h(\beta_k^{(i)})|x-\beta_k^{(i)}|\to h(\eta)|2-\eta|\,|\eta'|/2=1-2/\sqrt5$.
Likewise, $[a_0;\ldots,a_{s_k},2,2,\overline1]$ corresponds to
$\eta=2+1/(\varphi+1)=4-\varphi$ at $m=s_k+1$ (with $h(\eta)=2/\sqrt5$),
giving limit $(2-\varphi)(4-\varphi')/\sqrt5=4/\sqrt5-1$; the other three
families in (3), evaluated at $m=s_k$ where $T_{s_k}=A_k+1/2+o(1)$ and
$z_{s_k-1}=1/2+o(1)$, also converge to $(2/\sqrt5)(2-\sqrt5/2)=4/\sqrt5-1$.

To prove (2), suppose for contradiction that there exists a sequence
$(\beta_n)\subset E_\varphi\setminus\{\beta_k^{(1)},\beta_k^{(2)}\}$
with $h(\beta_n)\to\infty$ and
$E_n:=h(\beta_n)|x-\beta_n|\to L<4/\sqrt5-1$. Then $\beta_n\to x$, and
since $x$ is irrational, the length $m_n$ of the maximal common prefix of
the regular continued fractions of $x$ and $\beta_n$ tends to $\infty$.
Eventually $m_n\ge s_0+3$, and we consider two cases.

\emph{Case 1: $m_n\in\{s_k+1,s_k+2\}$.}
Since $\beta_n$ agrees with $x$ through $a_{s_k}$, we apply
\eqref{eq:bugeaud-34} at $m=s_k+1$, writing $\beta_n=M_{s_k+1}(\eta_n)$
with $\eta_n=[b_{s_k+1};b_{s_k+2},\ldots]\in E_\varphi\cap(1,\infty)$,
$T_m=2-\delta_k$ ($\delta_k\to0^+$), and $z_{m-1}\to0^+$. Because
$m_n\le s_k+2$, we cannot have $b_{s_k+1}=b_{s_k+2}=1$, so
$\eta_n\ne\varphi$; because $\beta_n\ne\beta_k^{(2)}$, we also have
$\eta_n\ne\varphi+1$. By Galois's theorem on purely periodic continued
fractions together with Serret's theorem (every element of $E_\varphi$
eventually has an all-ones tail), $\varphi=[\overline1]$ is the unique
element of $E_\varphi$ satisfying $\eta>1$ and $-1<\eta'<0$. Since
$\eta_n\ne\varphi$, either
$\eta_n'>0$ or $\eta_n'<-1$, and in both cases
$|\eta_n'+z_{m-1}|\ge(1-z_{m-1})|\eta_n'|$. Moreover, if $m_n=s_k+1$ then
$b_{s_k+1}\ge2$, so $\eta_n>2>T_m$ and $|T_m-\eta_n|>|2-\eta_n|$; if
$m_n=s_k+2$ then $b_{s_k+1}=1$ and $b_{s_k+2}\ge2$, so $\eta_n<3/2$ and
$|T_m-\eta_n|=(2-\eta_n)-\delta_k\ge(1-2\delta_k)|2-\eta_n|$. Applying
Lemma~\ref{lem:norm-parity-rigidity} to $\eta_n\notin\{\varphi,\varphi+1\}$
yields
\[
 E_n
 \ge \frac{(1-2\delta_k)(1-z_{m-1})}{1+z_{m-1}/2}\,
     h(\eta_n)\frac{|2-\eta_n|\,|\eta_n'|}{2}
 \ge (1-o(1))\left(\frac{4}{\sqrt5}-1\right),
\]
a contradiction.

\emph{Case 2: $m_n=s_k$.}
Apply \eqref{eq:bugeaud-34} at $m=s_k$, where
$T_{s_k}=A_k+[0;1,1,A_{k+1},\ldots]>A_k+1/2$ (with
$T_{s_k}=A_k+1/2+o(1)$) and
$z_{s_k-1}=[0;1,1,A_{k-1},\ldots]>1/2$ (with $z_{s_k-1}=1/2+o(1)$); in
particular, for $2\le Q_n\le A_k-1$ below, the positive remainders
ensure that the numerator product strictly exceeds the corresponding
concave quadratic uniformly in $Q_n$. Write $\eta_n=Q_n+\theta_n$ with
$Q_n=b_{s_k}\in\Z_{\ge1}\setminus\{A_k\}$ and
$\theta_n\in(0,1)\cap E_\varphi$. Then $h(\eta_n)=h(\theta_n)$. Writing
$\theta_n=(a\varphi+b)/(c\varphi+d)$ with $|ad-bc|=1$, we have
$h(\theta_n)=2|N(c\varphi+d)|/\sqrt5$. Because $2$ and $3$ are inert in
$\Z[\varphi]$ and $\gcd(c,d)=1$, $|N(c\varphi+d)|$ cannot lie in
$\{2,3,4\}$. Hence either $|N(c\varphi+d)|=1$, which gives
$h(\theta_n)=2/\sqrt5$ and restricts $\theta_n\in(0,1)$ to
$\{\varphi-1,\,2-\varphi\}$, or $|N(c\varphi+d)|\ge5$, giving
$h(\theta_n)\ge2\sqrt5$.
\begin{itemize}
\item If $\theta_n=\varphi-1=[0;\overline1]$, then $h(\eta_n)=2/\sqrt5$
      and $\theta_n'=-\varphi$. The choice $Q_n=1$ gives
      $\beta_n=[a_0;\ldots,a_{s_{k-1}},1,1,1,\overline1]=\beta_{k-1}^{(1)}$,
      excluded. If $Q_n\ge A_k+1$, then
      $(Q_n-\varphi+z_{s_k-1})/(T_{s_k}+z_{s_k-1})\ge1-o(1)$ and
      $|T_{s_k}-\eta_n|\ge\sqrt5/2+o(1)$, giving $E_n\ge1-o(1)$. If
      $2\le Q_n\le A_k-1$, the concave quadratic
      $(A_k-Q+1-\sqrt5/2)(Q-\sqrt5/2)$ attains its minimum on
      $[2,A_k-1]$ at both endpoints $Q=2$ and $Q=A_k-1$, yielding
      $E_n\ge(2/\sqrt5)(2-\sqrt5/2)-o(1)=4/\sqrt5-1-o(1)$.
\item If $\theta_n=2-\varphi=[0;2,\overline1]$, then $h(\eta_n)=2/\sqrt5$
      and $\theta_n'=1+\varphi$. For $Q_n=1$,
      $E_n\to(2/\sqrt5)(5/2+\varphi)=6/\sqrt5+1$. For $2\le Q_n\le A_k-1$,
      concavity of $(A_k-Q+\sqrt5/2-1)(Q+2+\sqrt5/2)$ at $Q=2$ and
      $Q=A_k-1$ yields limits $8/\sqrt5+1$ and $1$. For $Q_n\ge A_k+1$,
      $(\eta_n'+z_{s_k-1})/(T_{s_k}+z_{s_k-1})\ge1-o(1)$ and
      $|T_{s_k}-\eta_n|\ge2-\sqrt5/2+o(1)$, giving
      $E_n\ge4/\sqrt5-1-o(1)$.
\item If $\theta_n\notin\{\varphi-1,2-\varphi\}$, then
      $h(\eta_n)\ge2\sqrt5$. Applying Galois's theorem to
      $\tau_n=1/\theta_n\in E_\varphi\cap(1,\infty)\setminus\{\varphi\}$
      gives $\tau_n'\notin(-1,0)$, so $\theta_n'=1/\tau_n'>-1$ and
      $\eta_n'+z_{s_k-1}>Q_n-1/2+o(1)$. Meanwhile
      $|T_{s_k}-\eta_n|\ge|A_k-Q_n|-1/2+o(1)$. Checking $Q_n=1$,
      $Q_n\ge A_k+1$, and (by concavity of $(A_k-Q-1/2)(Q-1/2)$)
      $2\le Q_n\le A_k-1$ gives in every case
      $E_n\ge 2\sqrt5\cdot(1/2)-o(1)=\sqrt5-o(1)>4/\sqrt5-1$.
\end{itemize}
Every subcase contradicts $L<4/\sqrt5-1$, completing the proof.
\end{proof}

\begin{corollary}[Exact values for $6-e$, $e$, and $M(e^{1/m})$]\label{cor:golden-exact-e}
For every integer $m\ge1$ and every $M\in\mathrm{GL}_2(\Z)$,
\[
 c_\varphi(6-e)=c_\varphi(e)=c_\varphi\bigl(M(e^{1/m})\bigr)=1-\frac{2}{\sqrt5},
\]
and the second smallest limit point of $h(\beta)|x-\beta|$ for each such
$x$ is $4/\sqrt5-1$.
\end{corollary}
\begin{proof}
We first note that $c_\varphi$ (and the full set of finite limit points
of $h(\beta)|x-\beta|$ as $h(\beta)\to\infty$) is $\mathrm{GL}_2(\Z)$-invariant:
for $M=\bigl(\begin{smallmatrix}a&b\\c&d\end{smallmatrix}\bigr)\in\mathrm{GL}_2(\Z)$,
$y=M(x)$, and $\gamma=M(\beta)$, direct calculation gives
\[
 h(\gamma)|y-\gamma|=h(\beta)|x-\beta|\cdot\frac{|c\beta'+d|}{|cx+d|},
 \qquad h(\gamma)=h(\beta)|c\beta+d|\,|c\beta'+d|.
\]
Whenever $h(\beta_n)\to\infty$ with $h(\beta_n)|x-\beta_n|$ bounded, we
have $\beta_n\to x$ and $|\beta_n-\beta_n'|=2/h(\beta_n)\to0$, hence
$\beta_n'\to x$; since $x\notin\Q$, $|cx+d|>0$, so $h(\gamma_n)\to\infty$
and $h(\gamma_n)|y-\gamma_n|\sim h(\beta_n)|x-\beta_n|$. Symmetry via
$M^{-1}\in\mathrm{GL}_2(\Z)$ gives invariance.

By the classical Hurwitz continued fractions,
$e=[2;1,2,\overline{1,1,2k+2}]_{k\ge1}$ (for $m=1$) and
$e^{1/m}=[1;\overline{(2k+1)m-1,1,1}]_{k\ge0}$ (for $m\ge2$) both satisfy
Theorem~\ref{thm:golden-exact-second}, and $6-e=\bigl(\begin{smallmatrix}-1&6\\0&1\end{smallmatrix}\bigr)(e)$
lies in the $\mathrm{GL}_2(\Z)$-orbit of $e$.
\end{proof}

For $B>1$ the constant $c_\varphi(A/B-e)$ can drop well below
$1-2/\sqrt5$, and the norm parity mechanism of
Lemma~\ref{lem:norm-parity-rigidity} generalizes to arbitrary rational
complete-quotient endpoints $(u,v)=(p_1/q_1,p_2/q_2)$. We establish this
general rational-endpoint rigidity together with the exact evaluation of
$c_\varphi(11/2-e)$ and its second smallest limit point for the $B=2$
control model $x=11/2-e$.

\begin{lemma}[Rational-endpoint and Pell norm-congruence rigidity]\label{lem:pell-b2-rigidity}
\begin{enumerate}
\item For any rational numbers $u=p_1/q_1$ and $v=p_2/q_2$ with
      $q_1,q_2>0$ and $D:=p_1q_2+p_2q_1>0$, and any
      $\eta=(a\varphi+b)/\allowbreak(c\varphi+d)\in E_\varphi$ with
      $\varepsilon:=ad-bc\in\{\pm1\}$, set
      $\xi_1=(p_1c-q_1a)\varphi+(p_1d-q_1b)\in\Z[\varphi]$ and
      $\xi_2=(q_2a+p_2c)\varphi+(q_2b+p_2d)\in\Z[\varphi]$. Then
      \[
       \Phi_{u,v}(\eta):=h(\eta)\frac{|u-\eta|\,|\eta'+v|}{u+v}
       =\frac{|2\xi_1\xi_2'|}{D\sqrt5}
       =\left|\frac{\varepsilon M_{u,v}}{D\sqrt5}-1\right|,
      \]
      where $M_{u,v}\in\Z$ satisfies
      $M_{u,v}^2-5D^2=4N(\xi_1)N(\xi_2)$, and the $\Z[\varphi]$ field norm
      satisfies the mod-$5$ ramification identity
      $N(p\varphi+q)=-(p+2q)^2+5(p+q)q$ (in particular,
      $5\mid N(p\varphi+q)\iff 5\mid(p+2q)$).
\item For $q_1\in\{3,5\}$ and every $\eta\in E_\varphi$, there exists an
      \emph{odd} integer $m\in 2\Z+1$ such that
      \[
       \Phi_{8/q_1,0}(\eta)=h(\eta)\frac{|8/q_1-\eta|\,|\eta'|}{8/q_1}
       =\left|\frac{m}{4\sqrt5}-1\right|.
      \]
      Every integer $m\ne9$ satisfies
      $|m/(4\sqrt5)-1|\ge 1-2/\sqrt5$, and every \emph{odd} integer
      $m\ne9$ satisfies $|m/(4\sqrt5)-1|\ge 1-7/(4\sqrt5)$.
      Consequently:
      \begin{itemize}
      \item $\Phi_{8/5,0}(\eta)=9/(4\sqrt5)-1$ if and only if
            $\eta=\varphi$, and $\Phi_{8/5,0}(\eta)\ge 1-7/(4\sqrt5)$
            for all $\eta\in E_\varphi\setminus\{\varphi\}$ (with
            equality at $\eta=\varphi+1$);
      \item $\Phi_{8/3,0}(\eta)=9/(4\sqrt5)-1$ if and only if
            $\eta=\varphi+1$, and $\Phi_{8/3,0}(\eta)\ge 1-7/(4\sqrt5)$
            for all $\eta\in E_\varphi\setminus\{\varphi+1\}$ (with
            equality at $\eta=\varphi$).
      \end{itemize}
\end{enumerate}
\end{lemma}
\begin{proof}
For (1), writing $u-\eta=\xi_1/(q_1(c\varphi+d))$,
$\eta'+v=\xi_2'/(q_2(c\varphi'+d))$,
$h(\eta)=2|(c\varphi+d)(c\varphi'+d)|/\sqrt5$, and
$u+v=D/(q_1q_2)$, the denominators cancel to give
$\Phi_{u,v}(\eta)=2|\xi_1\xi_2'|/(D\sqrt5)$. Expanding $2\xi_1\xi_2'$ via
$\varphi\varphi'=-1$ and $\varphi+\varphi'=1$, the coefficient of
$\sqrt5$ is
$((p_1c-q_1a)(q_2b+p_2d)-(p_1d-q_1b)(q_2a+p_2c))=-\varepsilon D$, and
the rational part is an integer $M_{u,v}\in\Z$. Taking norms in
$\Q(\sqrt5)$ yields $M_{u,v}^2-5D^2=4N(\xi_1)N(\xi_2)$, while expanding
$-(p+2q)^2+5(p+q)q=-p^2+pq+q^2=N(p\varphi+q)$ gives the mod-$5$
ramification equivalence.

For (2), take $p_1=8$, $p_2=0$, $q_2=1$, so $D=8$, $\xi_2=\gamma:=a\varphi+b$,
and $\xi_1\gamma'=m_{8,q_1}-4\varepsilon\sqrt5$ with
\[
 m_{8,q_1}:=q_1(a^2-ab-b^2)-8ac+4ad+4bc+8bd\in\Z,
 \qquad m_{8,q_1}^2-80=N(\xi_1)N(\gamma).
\]
Direct expansion gives the congruence decomposition
\begin{equation}\label{eq:mod4-rigidity}
 m_{8,q_1}+q_1N(\gamma)=4(-2ac+2bd+ad+bc).
\end{equation}
Since $|ad-bc|=1$, $N(\gamma)$ is odd (Lemma~\ref{lem:norm-parity-rigidity});
because $q_1\in\{3,5\}$ is odd, \eqref{eq:mod4-rigidity} shows that
$m_{8,q_1}$ and $m:=\varepsilon m_{8,q_1}$ are odd, and
$\Phi_{8/q_1,0}(\eta)=|m/(4\sqrt5)-1|$. Since $9^2-5\cdot4^2=1$, we have
$8<44/5<4\sqrt5=\sqrt{80}<9$, so the closest integer to $4\sqrt5$ is $9$
(giving $9/(4\sqrt5)-1$), the next closest integer is $8$ (giving
$1-8/(4\sqrt5)=1-2/\sqrt5<10/(4\sqrt5)-1$), and because $4\sqrt5<9=(7+11)/2$,
the second closest \emph{odd} integer to $4\sqrt5$ is $7$ (giving
$1-7/(4\sqrt5)<11/(4\sqrt5)-1$).

Finally, equality $m=9$ gives $m_{8,q_1}=9\varepsilon$ and
$N(\xi_1)N(\gamma)=9^2-80=1$, so $N(\gamma)\in\{\pm1\}$ and
$\sigma:=\varepsilon N(\gamma)\in\{\pm1\}$. Multiplying
\eqref{eq:mod4-rigidity} by $\varepsilon$ gives
$9+q_1\sigma=4\varepsilon(-2ac+2bd+ad+bc)\equiv0\pmod4$. For $q_1=5$
this forces $\sigma=-1$, whereas for $q_1=3$ it forces $\sigma=+1$. Since
$\xi_1+q_1\gamma=8(c\varphi+d)$ and
$\xi_1/\gamma=\xi_1\gamma'/N(\gamma)=\sigma(9-4\sqrt5)$, we obtain
$8(c\varphi+d)=(5-(9-4\sqrt5))\gamma=4(\sqrt5-1)\gamma=(8/\varphi)\gamma$
for $q_1=5$ (hence $\eta=\varphi$), and
$8(c\varphi+d)=(3+(9-4\sqrt5))\gamma=4(3-\sqrt5)\gamma=(8/(\varphi+1))\gamma$
for $q_1=3$ (hence $\eta=\varphi+1$). Direct evaluation confirms the
values at $\varphi$ and $\varphi+1$.
\end{proof}

\begin{proposition}[Conjugate positivity on $E_\varphi\cap(1,\infty)\setminus\{\varphi\}$]\label{prop:conjugate-positivity}
Every $\eta\in E_\varphi\cap(1,\infty)\setminus\{\varphi\}$ satisfies
\[
 \eta'\ge\varphi^{-2}=\frac{3-\sqrt5}{2}>0,
\]
with equality at $\eta=\varphi+1=[2;\overline1]$. Consequently, for every
$z>0$, $|\eta'+z|=\eta'+z>|\eta'|$.
\end{proposition}
\begin{proof}
Write the shortest pre-periodic continued fraction of $\eta$ as
$\eta=[a_0;a_1,\ldots,a_k,\overline1]$ with $a_0\ge1$ and (when $k\ge1$)
$a_k\ge2$. If $k=0$, then $\eta\ne\varphi=[1;\overline1]$ forces $a_0\ge2$,
so $\eta=a_0-1+\varphi$ and $\eta'=a_0-1+\varphi'\ge 1+\varphi'=\varphi^{-2}>0$.
If $k\ge1$, then $x_k:=a_k+\varphi'\ge 2-\varphi^{-1}=\frac{5-\sqrt5}{2}>1$;
inducting backward via $x_{i-1}=a_{i-1}+1/x_i>1$ ($a_{i-1}\ge1$) gives
$\eta'=x_0>a_0\ge1>\varphi^{-2}$.
\end{proof}

\begin{theorem}[Exact golden-ratio orbit constant and second limit point for $11/2-e$]\label{thm:golden-b2-exact}
The regular continued fraction of $x:=11/2-e$ is
\[
 \frac{11}{2}-e
 =\bigl[2;1,3,1,3,1,1,1,\overline{2,k,2,1,1,1,k,1,1,1}_{k=1}^\infty\bigr],
\]
with unbounded partial quotients $a_{s_k}=a_{t_k}=k$ at
$s_k=10k-1$ and $t_k=10k+4$ ($k\ge1$). Define the two extremal orbit
families $\beta_k^{(1)}:=[a_0;\ldots,a_{s_k},2,\overline1]$ and
$\beta_k^{(2)}:=[a_0;\ldots,a_{t_k},\overline1]$. Then:
\begin{enumerate}
\item $c_\varphi(11/2-e)=9/(4\sqrt5)-1\approx 0.006230589875$, attained
      with limit $9/(4\sqrt5)-1$ along both $\beta_k^{(1)}$ and
      $\beta_k^{(2)}$.
\item Every sequence $(\beta_n)\subset E_\varphi$ with
      $h(\beta_n)\to\infty$ that is eventually disjoint from
      $\{\beta_k^{(1)},\beta_k^{(2)}:k\ge1\}$ satisfies
      \[
       \liminf_{n\to\infty} h(\beta_n)\left|\Bigl(\frac{11}{2}-e\Bigr)-\beta_n\right|
       \ge 1-\frac{7}{4\sqrt5}\approx 0.217376207875.
      \]
\item The second smallest limit point $1-7/(4\sqrt5)$ is sharply
      attained along each of the four families
      $[a_0;\ldots,a_{s_k},\overline1]$,
      $[a_0;\ldots,a_{t_k},2,\overline1]$,
      $[a_0;\ldots,a_{s_k},2,1,2,\overline1]$, and
      $[a_0;\ldots,a_{t_k},1,1,1,2,\overline1]$.
\end{enumerate}
\end{theorem}
\begin{proof}
Applying Theorem~\ref{thm:affine-window} to $x=(11-2e)/2$ ($A=11,B=2,C=-2$)
gives $L=4$, $T=8$, and $G_j=\gcd(11Q_j-2P_j,2Q_j)=1$ for all $j\ge1$
(since $Q_j$ is odd and $\gcd(P_j,Q_j)=1$), hence
$W_{j+8}=(b_j+4,S_j)$ for all $j\ge6$. Matching the exact Hermite
convergents $V_j=(11Q_j-2P_j,2Q_j)$ over the base period $W_6,\ldots,W_{13}$
($W_{2m+4}=[m+1,1,1,1,2]$ and $W_{2m+5}=[m+2,2,1,1,1]$ for $m=1,2,3,4$)
establishes the displayed period-$10$ expansion for all $k\ge1$.

At $m=s_k+1$, we have
$T_{s_k+1}=[2;1,1,1,k,\ldots]=(8T_{t_k}+5)/(3T_{t_k}+2)=8/3-\delta_k^{(1)}$
with $0<\delta_k^{(1)}=1/(3(3T_{t_k}+2))\to0^+$ and $z_{s_k}\to0^+$.
At $m=t_k+1$, we have
$T_{t_k+1}=[1;1,1,2,k+1,\ldots]=(8T_{s_{k+1}}+3)/(5T_{s_{k+1}}+2)=8/5-\delta_k^{(2)}$
with $0<\delta_k^{(2)}=1/(5(5T_{s_{k+1}}+2))\to0^+$ and $z_{t_k}\to0^+$.
Evaluating \eqref{eq:bugeaud-34} with $\eta=\varphi+1=[2;\overline1]$ at
$m=s_k+1$ (for $\beta_k^{(1)}$) and $\eta=\varphi=[\overline1]$ at
$m=t_k+1$ (for $\beta_k^{(2)}$) yields limit $9/(4\sqrt5)-1$ by
Lemma~\ref{lem:pell-b2-rigidity}(2). Likewise, taking $\eta=\varphi$ and
$\eta=[2;1,2,\overline1]=(11\varphi+8)/(4\varphi+3)$ at $m=s_k+1$, or
$\eta=\varphi+1$ and $\eta=[1;1,1,2,\overline1]=(11\varphi+8)/(7\varphi+5)$
at $m=t_k+1$, gives $m=7$ in Lemma~\ref{lem:pell-b2-rigidity}(2) and
hence limit $1-7/(4\sqrt5)$, proving (3) and the upper bound in (1).

To prove (2), let $(\beta_n)\subset E_\varphi\setminus\{\beta_k^{(1)},\beta_k^{(2)}\}$
have $h(\beta_n)\to\infty$ and $E_n:=h(\beta_n)|x-\beta_n|\to L<1-7/(4\sqrt5)$.
Then the index $m_n$ of the first continued-fraction mismatch between
$x$ and $\beta_n$ tends to $\infty$.

\emph{Case 1: $m_n\in\{s_k+1,s_k+2,s_k+3,s_k+4\}$ or
$m_n\in\{t_k+1,t_k+2,t_k+3,t_k+4\}$.}
In the first subcase, apply \eqref{eq:bugeaud-34} at $m=s_k+1$ with
$\beta_n=M_{s_k+1}(\eta_n)$, $\eta_n=[b_1;b_2,b_3,b_4,\ldots]\in E_\varphi\cap(1,\infty)$.
Since $m_n\le s_k+4$, $(b_1,b_2,b_3,b_4)\ne(2,1,1,1)$, so
$\eta_n\ne\varphi+1$; moreover, checking the four prefix possibilities
($b_1\ne2$; $(2,b_2\ge2)$; $(2,1,b_3\ge2)$; and $(2,1,1,b_4\ge2)$) shows
that either $\eta_n>8/3>T_{s_k+1}$ or $\eta_n<[2;1,1,2]=13/5$, so
$|T_{s_k+1}-\eta_n|\ge(1-15\delta_k^{(1)})|8/3-\eta_n|$. In the second
subcase, apply \eqref{eq:bugeaud-34} at $m=t_k+1$ with
$\beta_n=M_{t_k+1}(\eta_n)$. Since $\beta_n\ne\beta_k^{(2)}$, we have
$\eta_n\ne\varphi$. Because $(b_1,b_2,b_3,b_4)\ne(1,1,1,2)$, either
$b_1\ge2$, $(1,b_2\ge2)$, $(1,1,b_3\ge2)$, $(1,1,1,1)$ (where
$\eta_n=[1;1,1,1,t]=(5t+3)/(3t+2)>8/5>T_{t_k+1}$ for all $t>1$), or
$(1,1,1,b_4\ge3)$ (where $\eta_n<[1;1,1,3]=11/7$), so in all cases
$|T_{t_k+1}-\eta_n|\ge(1-35\delta_k^{(2)})|8/5-\eta_n|$. By
Proposition~\ref{prop:conjugate-positivity}, since
$\eta_n\in E_\varphi\cap(1,\infty)\setminus\{\varphi\}$ in both subcases
($b_{s_k+1}=2$ in the first subcase and $\eta_n\ne\varphi$ in the second),
we have $\eta_n'\ge\varphi^{-2}>0$ and hence
$|\eta_n'+z_{m-1}|=\eta_n'+z_{m-1}>|\eta_n'|$.
Lemma~\ref{lem:pell-b2-rigidity}(2)
therefore gives $E_n\ge(1-o(1))(1-7/(4\sqrt5))$ in both subcases.

\emph{Case 2: $m_n=s_k$ or $m_n=t_k$.}
At $m=s_k$ (with $r=3/8$) or $m=t_k$ (with $r=5/8$), we have
$T_m=k+r+o(1)$ and $z_{m-1}=r+o(1)$. Write $\eta_n=Q_n+\theta_n$ with
$Q_n\in\Z_{\ge1}\setminus\{k\}$ and $\theta_n\in E_\varphi\cap(0,1)$. By
the height dichotomy in the proof of Theorem~\ref{thm:golden-exact-second},
either $\theta_n\in\{\varphi-1,2-\varphi\}$ with $h(\eta_n)=2/\sqrt5$,
or $h(\eta_n)\ge2\sqrt5$ with $\theta_n'>-1$:
\begin{itemize}
\item If $\theta_n=\varphi-1$ ($\theta_n'=-\varphi$), then at $m=t_k$
      ($r=5/8$) the choice $Q_n=1$ gives
      $[a_0;\ldots,a_{s_k},2,1,1,1,1,\overline1]=\beta_k^{(1)}$, excluded,
      while at $m=s_k$ ($r=3/8$) the choice $Q_n=1$ gives limit
      $(2/\sqrt5)|1-\varphi+3/8|=1-7/(4\sqrt5)$
      ($=[a_0;\ldots,a_{t_{k-1}},1,1,1,2,\overline1]$). For
      $2\le Q_n\le k-1$, concavity of
      $(k+r-Q_n-(\varphi-1))(Q_n-\varphi+r)$ on $[2,k-1]$ bounds $E_n$
      below by its endpoint limit
      $(2/\sqrt5)(2-\varphi+r)\ge 15/(4\sqrt5)-1>1-7/(4\sqrt5)$, and for
      $Q_n\ge k+1$ we have $E_n\ge(2/\sqrt5)(\varphi-r)-o(1)\ge 1-1/(4\sqrt5)-o(1)$.
\item If $\theta_n=2-\varphi$ ($\theta_n'=1+\varphi$), checking $Q_n=1$,
      $2\le Q_n\le k-1$, and $Q_n\ge k+1$ gives
      $E_n\ge(2/\sqrt5)(3-\varphi-r)-o(1)\ge 15/(4\sqrt5)-1-o(1)$.
\item If $\theta_n\notin\{\varphi-1,2-\varphi\}$, then $h(\eta_n)\ge2\sqrt5$,
      $\eta_n'+z_{m-1}\ge Q_n-1+r-o(1)\ge 3/8-o(1)$, and
      $|T_m-\eta_n|\ge|k-Q_n|-(1-r)-o(1)\ge 3/8-o(1)$, yielding
      $E_n\ge 2\sqrt5\cdot(3/8)-o(1)=3\sqrt5/4-o(1)>1-7/(4\sqrt5)$.
\end{itemize}
Every case contradicts $L<1-7/(4\sqrt5)$.
\end{proof}

To extend exact golden-ratio orbit evaluation beyond $B=1,2$ to all
coprime models $A/3-e$ and $A/4-e$, we upgrade
Lemma~\ref{lem:pell-b2-rigidity} to general rational endpoints and
establish a uniform separation lemma for even-length suffixes.

\begin{lemma}[General rational-endpoint arithmetic rigidity and algebraic inversion]\label{lem:general-endpoint-rigidity}
Let $u=p_1/q_1$ and $v=p_2/q_2$ be reduced rationals with
$\gcd(p_1,q_1)=\gcd(p_2,q_2)=1$, $q_1,q_2>0$, and $D:=p_1q_2+p_2q_1>0$.
For $\eta=(a\varphi+b)/(c\varphi+d)\in E_\varphi$ ($\varepsilon:=ad-bc\in\{\pm1\}$),
keep the notation of Lemma~\ref{lem:pell-b2-rigidity}(1):
$\xi_1=(p_1c-q_1a)\varphi+(p_1d-q_1b)$,
$\xi_2=(q_2a+p_2c)\varphi+(q_2b+p_2d)$, and
$2\xi_1\xi_2'=M_{u,v}-\varepsilon D\sqrt5$, so that
$\Phi_{u,v}(\eta)=|\varepsilon M_{u,v}/(D\sqrt5)-1|$ and
$M_{u,v}^2-5D^2=4N(\xi_1)N(\xi_2)$. Then:
\begin{enumerate}
\item For every prime $\ell$ inert in $\Z[\varphi]$ (i.e., $\ell=2$ or
      $\ell\equiv\pm2\pmod5$), $\ell\nmid N(\xi_1)$ and $\ell\nmid N(\xi_2)$.
\item If $D=2d$ is even, then $M_{u,v}=2m_0$ is even and $m:=\varepsilon m_0\in\Z$
      satisfies
      \[
       \Phi_{u,v}(\eta)=\left|\frac{m}{d\sqrt5}-1\right|,\qquad
       m^2-5d^2=N(\xi_1)N(\xi_2),\qquad m\equiv d+1\pmod2.
      \]
      Moreover, if an inert prime $\ell$ divides $d$, then $\ell\nmid m$.
\item We have the congruences
      \[
       q_2(M_{u,v}+\varepsilon D)+2q_1N(\xi_2)\equiv0\pmod{2D},\qquad
       p_2(M_{u,v}+\varepsilon D)-2p_1N(\xi_2)\equiv0\pmod{2D}.
      \]
      When $D=2d$, these reduce to
      $q_2(\varepsilon m+d)+q_1N(\xi_2)\equiv0\pmod{2d}$ and
      $p_2(\varepsilon m+d)-p_1N(\xi_2)\equiv0\pmod{2d}$; in particular,
      at $(u,v)=(2d/q_1,0)$ ($p_2=0,q_2=1$, $\gamma:=a\varphi+b$) one has
      $m\equiv d-\varepsilon q_1N(\gamma)\pmod{2d}$.
\item If $\eta\in E_\varphi$ attains $M_{u,v}=\sigma_M M$ ($\sigma_M\in\{\pm1\}$)
      at $(u,v)$, set
      $R:=\xi_1/\xi_2=(\xi_1\xi_2')/N(\xi_2)=(\sigma_M M-\varepsilon D\sqrt5)/(2N(\xi_2))$.
      Then
      \[
       \eta=\frac{p_1-p_2R}{q_1+q_2R},
      \]
      where $N(\xi_2)$ divides the nonzero integer $(M^2-5D^2)/4$. In
      particular, for fixed $(u,v)$ and $M$, all solutions $\eta\in E_\varphi$
      lie among at most $4\tau(|M^2-5D^2|/4)$ explicit algebraic candidates.
\end{enumerate}
\end{lemma}
\begin{proof}
For (1), let $\ell$ be any prime. The matrix $\bigl(\begin{smallmatrix}-a&c\\-b&d\end{smallmatrix}\bigr)$
has determinant $-\varepsilon\in\{\pm1\}$ and is invertible over $\mathbb{F}_\ell$.
Since $\gcd(p_1,q_1)=1$, $(q_1,p_1)^T\not\equiv(0,0)\pmod\ell$, so the
$\{1,\varphi\}$-coefficients of $\xi_1$ are not both divisible by $\ell$.
Similarly $\gcd(p_2,q_2)=1$ gives $\ell\nmid\xi_2$ in $\Z[\varphi]$. Since
an inert prime $\ell$ remains prime in $\Z[\varphi]$,
$\ell\mid N(\xi)\iff\ell\mid\xi$, whence $\ell\nmid N(\xi_1)$ and
$\ell\nmid N(\xi_2)$.

For (2), $D=2d$ gives $M_{u,v}^2=5(2d)^2+4N(\xi_1)N(\xi_2)\equiv0\pmod4$,
so $M_{u,v}=2m_0$ is even and $m^2-5d^2=N(\xi_1)N(\xi_2)$. By (1), $2\nmid N(\xi_1)N(\xi_2)$,
so $m^2-5d^2\equiv1\pmod2$, i.e., $m\equiv d+1\pmod2$. If an inert prime
$\ell$ divides $d$, then (1) implies $\ell\nmid(m^2-5d^2)$, hence $\ell\nmid m$.

For (3), direct linear elimination gives
$q_2\xi_1+q_1\xi_2=D(c\varphi+d)$ and $-p_2\xi_1+p_1\xi_2=D(a\varphi+b)$.
Multiplying the first relation by $2\xi_2'$ and substituting $\sqrt5=2\varphi-1$
into $2\xi_1\xi_2'=(M_{u,v}+\varepsilon D)-2\varepsilon D\varphi$ yields
\[
 q_2(M_{u,v}+\varepsilon D)-2q_2\varepsilon D\varphi+2q_1N(\xi_2)
 =2D(c\varphi+d)\xi_2'\in 2D\Z[\varphi].
\]
Comparing the constant coefficient in the $\Z$-basis $\{1,\varphi\}$ gives
$q_2(M_{u,v}+\varepsilon D)+2q_1N(\xi_2)\equiv0\pmod{2D}$; doing the same
with $-p_2\xi_1+p_1\xi_2=D(a\varphi+b)$ yields
$p_2(M_{u,v}+\varepsilon D)-2p_1N(\xi_2)\equiv0\pmod{2D}$. When $D=2d$,
$M_{u,v}=2m_0=2\varepsilon m$ (from $m=\varepsilon m_0$); dividing by~$2$
and noting $-d\equiv d\pmod{2d}$ yields the reduced congruences and the
specialization at $(p_2,q_2)=(0,1)$.

For (4), dividing $D(a\varphi+b)=p_1\xi_2-p_2\xi_1$ by
$D(c\varphi+d)=q_1\xi_2+q_2\xi_1$ (noting $D>0$ and $\xi_2\ne0$ since
$v\in\Q$ while $\eta'\notin\Q$) gives $\eta=(p_1-p_2R)/(q_1+q_2R)$.
Since $N(\xi_1)N(\xi_2)=(M^2-5D^2)/4\ne0$ (because $5$ is not a square)
and $\sigma_M,\varepsilon\in\{\pm1\}$ with signed divisors $N(\xi_2)$ give
at most $2\cdot 2\tau(|M^2-5D^2|/4)$ choices for $R$, the bound follows.
\end{proof}

\begin{lemma}[Uniform separation for even-length suffixes]\label{lem:even-suffix-sep}
Let $S=(s_1,\ldots,s_{2t})\in\Z_{\ge1}^{2t}$ be a finite sequence of even
length, write $u=[s_1;s_2,\ldots,s_{2t}]$, and let
$T=[s_1;\ldots,s_{2t},b,\ldots]=u-\delta$ ($b\ge1$, $\delta>0$). Then for
every irrational $\eta=[b_1;b_2,\ldots]\in(1,\infty)$ whose first $2t$
partial quotients satisfy $(b_1,\ldots,b_{2t})\ne(s_1,\ldots,s_{2t})$,
either $\eta>u>T$ (so $|T-\eta|>|u-\eta|$) or $\eta\le u-\Delta(S)$, where
$\Delta(S)>0$ depends only on $S$. In particular, as $\delta\to0^+$,
\[
 |T-\eta|\ge\bigl(1-O_S(\delta)\bigr)|u-\eta|
\]
uniformly over all irrationals $\eta>1$ with $(b_1,\ldots,b_{2t})\ne S$.
\end{lemma}
\begin{proof}
Let $i\in\{1,\ldots,2t\}$ be the least $1$-based index with $b_i\ne s_i$.
By standard monotonicity of finite continued fractions:
\begin{itemize}
\item if $i=2r-1$ and $b_i>s_i$, then $\eta>[s_1;\ldots,s_{2r-2},s_{2r-1}+1]>u$;
\item if $i=2r$ and $b_i<s_i$, irrationality of $\eta$ gives
      $b_{2r}+1/(\cdots)<b_{2r}+1\le s_{2r}$, so $\eta>u$;
\item if $i=2r-1$ and $b_i<s_i$, then
      $\eta<[s_1;\ldots,s_{2r-2},b_i+1]\le[s_1;\ldots,s_{2r-1}]<u$;
\item if $i=2r$ and $b_i>s_i$, then $\eta<[s_1;\ldots,s_{2r-1},s_{2r}+1]<u$.
\end{itemize}
Setting
$\Delta(S):=u-\max\bigl(\{[s_1;\ldots,s_{2r-1}]:1\le r\le t\}\cup\{[s_1;\ldots,s_{2r-1},s_{2r}+1]:1\le r\le t\}\bigr)>0$
completes the proof.
\end{proof}

\begin{theorem}[Exact golden-ratio orbit constant and second limit point for $B=3$]\label{thm:golden-b3-exact}
For every integer $A$ with $3\nmid A$, write $x:=A/3-e=[a_0;a_1,a_2,\ldots]$. Then:
\begin{enumerate}
\item For all $j\ge0$, $3\nmid Q_j$ (in fact $Q_j\bmod3$ has period~$6$ given
      by $(1,1,1,2,2,2)$), so $G_j=1$ and $D_j=18$ for $j>9$. Up to
      $\mathrm{GL}_2(\Z)$-equivalence, $x$ reduces to $10/3-e$ ($A\equiv1\pmod3$)
      or $11/3-e$ ($A\equiv2\pmod3$). In both classes, the windows
      $W_j=(b_j,S_j)$ repeat with period $T=18$ and increment
      $b_{j+18}=b_j+4$, and exhibit the same six suffixes $S^{(q_0)}$ of
      value $18/q_0$ for $q_0\in\{1,5,7,11,13,17\}$:
      \[
       S^{(1)}=[17,1],\quad S^{(5)}=[3,1,1,2],\quad S^{(7)}=[2,1,1,3],
      \]
      \[
       S^{(11)}=[1,1,1,1,2,1],\quad S^{(13)}=[1,2,1,1,1,1],\quad S^{(17)}=[1,17].
      \]
\item We have
      \[
       c_\varphi(A/3-e)=1-\frac{20}{9\sqrt5}\approx 0.0061920100000935,
      \]
      attained with limit $1-20/(9\sqrt5)$ along—and only along—the four
      pairwise-disjoint extremal orbit families indexed by $j$ with $b_j=a_{k(j)+1}$:
      $[a_0;\ldots,b_j,3,\overline1]$ (at $S^{(5)}$),
      $[a_0;\ldots,b_j,2,\overline1]$ (at $S^{(7)}$),
      $[a_0;\ldots,b_j,\overline1]$ (at $S^{(11)}$), and
      $[a_0;\ldots,b_j,1,2,\overline1]$ (at $S^{(13)}$).
\item Every sequence $(\beta_n)\subset E_\varphi$ with $h(\beta_n)\to\infty$
      eventually disjoint from these four extremal families satisfies
      \[
       \liminf_{n\to\infty} h(\beta_n)|x-\beta_n|\ge\frac{22}{9\sqrt5}-1\approx 0.0931887889998972.
      \]
\item The second smallest limit point $c_{\mathrm{sec}}^{(3)}:=\frac{22}{9\sqrt5}-1$
      is sharply attained along each of the four families
      $[a_0;\ldots,b_j,2,\overline1]$ (at $S^{(5)}$),
      $[a_0;\ldots,b_j,2,1,1,2,\overline1]$ (at $S^{(7)}$),
      $[a_0;\ldots,b_j,1,1,1,1,2,\overline1]$ (at $S^{(11)}$), and
      $[a_0;\ldots,b_j,\overline1]$ (at $S^{(13)}$).
\end{enumerate}
\end{theorem}
\begin{proof}
Since $Q_0=Q_1=1$ and $Q_{j+1}=(4j+2)Q_j+Q_{j-1}\equiv(j+2)Q_j+Q_{j-1}\pmod3$,
six steps verify $(Q_0,\ldots,Q_7)\bmod3=(1,1,1,2,2,2,1,1)$, so $3\nmid Q_j$
and $2\nmid Q_j$ for all $j\ge0$. Thus $G_j=\gcd(AQ_j-3P_j,3Q_j)=1$ and
$D_j=18$ for $j>9$ by Theorem~\ref{thm:affine-window}. Integer shifts preserve
$\mathrm{GL}_2(\Z)$-equivalence (Corollary~\ref{cor:golden-exact-e}), reducing
to $10/3-e$ and $11/3-e$, whose period-$18$ window suffixes $W_{10},\ldots,W_{27}$
are directly verified to be the six listed even-length sequences.

At $D=18$ ($d=9$), Lemma~\ref{lem:general-endpoint-rigidity}(2) forces $m$ to
be even with $3\nmid m$ (since $3\mid 9$ is inert in $\Z[\varphi]$). Because
$9\sqrt5=\sqrt{405}\in(20,21)$, the closest even integer to $9\sqrt5$ is $m=20$,
and since $3\mid 18$, the second closest admissible even integer is $m=22$
(as $22-9\sqrt5<9\sqrt5-16$). Applying Lemma~\ref{lem:general-endpoint-rigidity}(4)
with $m=20$ ($20^2-5\cdot9^2=-5$) to the six endpoints $(18/q_0,0)$ yields
unique solutions in $E_\varphi$: none for $q_0\in\{1,17\}$, $[3;\overline1]$
for $q_0=5$, $[2;\overline1]$ for $q_0=7$, $\varphi=[\overline1]$ for $q_0=11$,
and $[1;2,\overline1]$ for $q_0=13$ (noting $[a_0;\ldots,b_j,1,2,1,1,1,1,1,\overline1]=[a_0;\ldots,b_j,1,2,\overline1]$
when $b_{j+1}$ is replaced by~$1$). Likewise $m=22$ ($22^2-5\cdot9^2=79$)
gives $[2;\overline1]$, $[2;1,1,2,\overline1]$, $[1;1,1,1,2,\overline1]$, and
$\varphi=[\overline1]$ at $q_0=5,7,11,13$, respectively. Evaluating~\eqref{eq:bugeaud-34}
at $m=k(j)+2$ proves (4) and the upper bound in (2).

For (3), let $(\beta_n)\subset E_\varphi$ be eventually disjoint from the four
extremal families with $h(\beta_n)\to\infty$ and $E_n:=h(\beta_n)|x-\beta_n|\to L<c_{\mathrm{sec}}^{(3)}$.
The first mismatch index $m_n$ tends to $\infty$.

\emph{Case 1: $m_n$ lies inside a bounded suffix $S_j$.}
Apply~\eqref{eq:bugeaud-34} at $m=k(j)+2$, where $T_m=18/q_0-\delta_j$
($\delta_j\to0^+$) and $z_{m-1}\to0^+$, with $\beta_n=M_m(\eta_n)$ and
$\eta_n\in E_\varphi\cap(1,\infty)$. Since $m_n$ lies inside $S_j$, $\eta_n$
differs from the unique $m=20$ solution at $18/q_0$ (for $q_0=13$, any
$\eta_n$ starting with $S^{(13)}=(1,2,1,1,1,1)$ would match past $S_j$).
By Lemma~\ref{lem:even-suffix-sep}, $|T_m-\eta_n|\ge(1-O(\delta_j))|18/q_0-\eta_n|$.
If $\eta_n\ne\varphi$, Proposition~\ref{prop:conjugate-positivity} gives
$|\eta_n'+z_{m-1}|=\eta_n'+z_{m-1}>|\eta_n'|$; if $\eta_n=\varphi$ (so
$q_0\ne11$ by exclusion of $\beta_j^{(11)}$), direct evaluation of
$\Phi_{18/q_0,0}(\varphi)$ yields $m\in\{10,14,16,22,26\}$ and
$|\varphi'+z_{m-1}|\ge(1-2z_{m-1})|\varphi'|$. In every case
Lemma~\ref{lem:general-endpoint-rigidity}(2) forces $E_n\ge(1-o(1))c_{\mathrm{sec}}^{(3)}$.

\emph{Case 2: $m_n=k(j)+1$ occurs at an unbounded digit $b_j$.}
Here $T_m=b_j+r_1+o(1)$ and $z_{m-1}=r_2+o(1)$ with
$r_1,r_2\in\{1/18,5/18,7/18,11/18,13/18,17/18\}$ (reverse suffixes give
$r_2=[0;\mathrm{rev}(S_{j-1})]=q_1'/18$). Write $\eta_n=Q_n+\theta_n$ with
$Q_n\in\Z_{\ge1}\setminus\{b_j\}$ and $\theta_n\in E_\varphi\cap(0,1)$:
\begin{itemize}
\item If $\theta_n=\varphi-1$ ($\theta_n'=-\varphi$, $h(\eta_n)=2/\sqrt5$),
      then $Q_n=1$ yields $E_n\to|(18+2q_1')/(18\sqrt5)-1|$. Here $q_1'=11$
      corresponds to $S_{j-1}=S^{(13)}$ (producing $\beta_{j-1}^{(13)}$, excluded),
      $q_1'=13$ gives $c_{\mathrm{sec}}^{(3)}$, and $q_1'\in\{1,5,7,17\}$ give
      at least $1-16/(9\sqrt5)>0.204$. For $2\le Q_n\le b_j-1$, concavity on
      $[2,b_j-1]$ bounds $E_n$ below by the endpoint limits
      $(2/\sqrt5)(2-\varphi+r_2)$ and $(2/\sqrt5)(2-\varphi+r_1)$, both at
      least $(2/\sqrt5)(2-\varphi+1/18)\approx 0.39133$; for $Q_n\ge b_j+1$,
      $E_n\ge(2/\sqrt5)(\varphi-17/18)-o(1)\approx 0.60248$.
\item If $\theta_n=2-\varphi$ ($\theta_n'=1+\varphi$, $h(\eta_n)=2/\sqrt5$),
      then $Q_n=1$, $2\le Q_n\le b_j-1$, and $Q_n\ge b_j+1$ give lower bounds
      $\ge(2/\sqrt5)(2+\varphi)>3.23$, $\ge(2/\sqrt5)(\varphi-1+1/18)\approx0.60248$,
      and $\ge(2/\sqrt5)(3-\varphi-17/18)\approx0.39133$, respectively.
\item If $\theta_n\notin\{\varphi-1,2-\varphi\}$, then $h(\eta_n)\ge2\sqrt5$,
      and $\tau_n=1/\theta_n\in E_\varphi\cap(1,\infty)\setminus\{\varphi\}$
      satisfies $\theta_n'=1/\tau_n'\in(0,\varphi^2)$ by
      Proposition~\ref{prop:conjugate-positivity}. Checking $Q_n=1$
      ($\ge 2\sqrt5\varphi^{-2}\approx1.7082$), $2\le Q_n\le b_j-1$
      ($\ge 2\sqrt5 r_1\ge\sqrt5/9\approx0.24845$), and $Q_n\ge b_j+1$
      ($\ge 2\sqrt5(1-r_1)\ge\sqrt5/9\approx0.24845$) excludes $L<c_{\mathrm{sec}}^{(3)}$.
\end{itemize}
All branches contradict $L<c_{\mathrm{sec}}^{(3)}$, proving (2) and (3).
\end{proof}

\begin{remark}[An explicit extremal window for $17/3-e$]\label{rem:golden-b3-window}
For $x=17/3-e$, the adjacent locked vectors at $j=12,13$ are the
reductions of $(17Q_j-3P_j,3Q_j)$. The Euclidean algorithm, with the
finite expansion chosen by the parity condition in
Theorem~\ref{thm:affine-window}, gives
$W_{12}=(2,1,2,1,1,1,1)$ and $C_{12}/D_{12}=49/18$.
Thus $j\equiv12\pmod{18}$ has suffix $S^{(13)}$. In~\eqref{eq:bugeaud-34}
cut the common prefix just after the first $1$ in that suffix.
Then $T_m\to13/5$ and $z_{m-1}\to1$ along this progression.
Choosing $\eta=\varphi+1$ gives an orbit point with height tending
to infinity and normalized error tending to
$1-4\sqrt5/9=1-20/(9\sqrt5)$. This is the $S^{(13)}$ extremal
family of Theorem~\ref{thm:golden-b3-exact}. The exact enumeration of
Theorem~\ref{thm:golden-finite-shift} independently recovers the first
two distinct values for this input, using the residue class of $17$
modulo $3$ and invariance under integer translation.
\end{remark}

\begin{theorem}[Exact golden-ratio orbit constant, mod-$16$ exclusion, and second limit point for $B=4$]\label{thm:golden-b4-exact}
For every odd integer $A$, write $x:=A/4-e=[a_0;a_1,a_2,\ldots]$. Then:
\begin{enumerate}
\item Since $Q_j$ is always odd, $G_j=1$ and $D_j=32$ for $j>16$, and
      $\mathrm{GL}_2(\Z)$-equivalence reduces to the two classes
      $13/4-e$ ($A\equiv1\pmod4$) and $15/4-e$ ($A\equiv3\pmod4$). For
      both classes, starting at $j=17$, both the suffix $S_j$ and the
      window length $r_j=1+|S_j|$ repeat with minimal period~$8$, while
      $b_{j+8}=b_j+1$ (so $b_{j+32}=b_j+4$ at the affine period $T=32$).
      The suffixes have values $32/q_0$ with:
      \begin{itemize}
      \item for $A\equiv1\pmod4$: $q_0\in\{7,15,23,25,31\}$
            ($[4,1,1,3]$, $[2,7,1,1]$, $[1,2,1,1,3,1]$, $[1,3,1,1,2,1]$, $[1,31]$);
      \item for $A\equiv3\pmod4$: $q_0\in\{7,9,15,23,31\}$
            (replacing $S^{(25)}$ by $S^{(9)}=[3,1,1,4]$).
      \end{itemize}
      Every occurring $q_0$ satisfies $q_0\equiv\pm1\pmod8$, i.e.,
      $q_0\bmod16\in\{\pm1,\pm7\}$.
\item For every odd integer $A$,
      \[
       c_\varphi(A/4-e)=1-\frac{35}{16\sqrt5}\approx 0.0217202598438420.
      \]
      This minimum is attained along—and only along—three pairwise-disjoint
      families in each class: $[a_0;\ldots,b_j,4,\overline1]$ (at $S^{(7)}$),
      $[a_0;\ldots,b_j,1,2,\overline1]$ (at $S^{(23)}$), and
      $[a_0;\ldots,b_j,1,3,\overline1]$ (at $S^{(25)}$) for $A\equiv1\pmod4$;
      and $[a_0;\ldots,b_j,4,\overline1]$ (at $S^{(7)}$),
      $[a_0;\ldots,b_j,3,\overline1]$ (at $S^{(9)}$), and
      $[a_0;\ldots,b_j,1,2,\overline1]$ (at $S^{(23)}$) for $A\equiv3\pmod4$.
\item \emph{Mod-$16$ congruence exclusion of $m=37$.} At $D=32$ ($d=16$),
      $m$ must be odd. Although $m=37$ (giving $37/(16\sqrt5)-1\approx0.03418144$)
      is closer to $16\sqrt5\approx35.777088$ than $m=33$ (giving
      $1-33/(16\sqrt5)\approx0.07762196$) and $37^2-5\cdot16^2=89$ has no
      inert prime factors, every divisor $N(\gamma)\in\{\pm1,\pm89\}$ satisfies
      $N(\gamma)\bmod16\in\{\pm1,\pm7\}$. By Lemma~\ref{lem:general-endpoint-rigidity}(3),
      $m\equiv-\varepsilon q_0N(\gamma)\pmod{16}$ is confined to the order-$4$
      multiplicative subgroup $\{\pm1,\pm7\}\pmod{16}$, which excludes
      $m=37\equiv5\pmod{16}$ completely.
\item Every sequence $(\beta_n)\subset E_\varphi$ with $h(\beta_n)\to\infty$
      eventually disjoint from the three extremal families satisfies
      \[
       \liminf_{n\to\infty} h(\beta_n)|x-\beta_n|
       \ge c_{\mathrm{sec}}^{(4)}:=1-\frac{33}{16\sqrt5}\approx 0.0776219592813368,
      \]
      and this second smallest limit point is sharply attained in both
      classes at $S^{(15)}=[2,7,1,1]$ by both $[a_0;\ldots,b_j,2,\overline1]$
      and $[a_0;\ldots,b_j,2,7,\overline1]$.
\end{enumerate}
\end{theorem}
\begin{proof}
Because $P_j,Q_j\pmod{16}$ has minimal period~$8$ (and period~$16$ modulo~$32$),
direct computation of $W_{17},\ldots,W_{24}$ together with
Theorem~\ref{thm:affine-window} establishes $G_j=1$, $D_j=32$, minimal
period~$8$ for $S_j$ and $r_j$ with $b_{j+8}=b_j+1$, and the stated five-suffix
sets for $13/4-e$ and $15/4-e$.

At $D=32$ ($d=16$), Lemma~\ref{lem:general-endpoint-rigidity}(2) forces $m$ to
be odd. Since $16\sqrt5\approx35.777088$, the closest odd integer is $m=35$,
followed by $m=37$ and then $m=33$. Part~(3) rules out $m=37$ across all
$q_0\in\{7,9,15,23,25,31\}$ since $q_0\bmod16\in\{\pm1,\pm7\}$ and
$89\equiv9\equiv-7\pmod{16}$, whereas $37\equiv5\pmod{16}$. Applying
Lemma~\ref{lem:general-endpoint-rigidity}(4) at $m=35$ ($35^2-1280=-55$)
yields unique solutions $[4;\overline1]$ at $q_0=7$, $[3;\overline1]$ at
$q_0=9$, $[1;2,\overline1]$ at $q_0=23$, $[1;3,\overline1]$ at $q_0=25$, and
none at $q_0\in\{15,31\}$. At $m=33$ ($33^2-1280=-191$), solutions occur
exclusively at $q_0=15$, namely $[2;\overline1]$ and $[2;7,\overline1]$.
This proves the upper bound in (2) and the sharp attainment in (4).

For the lower bound in (4), apply~\eqref{eq:bugeaud-34} as in
Theorem~\ref{thm:golden-b3-exact}. In \emph{Case 1} ($m_n$ inside $S_j$),
Lemma~\ref{lem:even-suffix-sep}, Proposition~\ref{prop:conjugate-positivity}
(noting $\Phi_{32/q_0,0}(\varphi)$ yields $m\in\{23,25,31,39,41,47\}$), and
the exclusion of $m=35,37$ force $E_n\ge(1-o(1))c_{\mathrm{sec}}^{(4)}$.
In \emph{Case 2} ($m_n=k(j)+1$ at $b_j$), we have $T_m=b_j+r_1+o(1)$ and
$z_{m-1}=r_2+o(1)$ with $r_1\in[7/32,31/32]$ and
$r_2=[0;\mathrm{rev}(S_{j-1})]\in\{1/32,7/32,9/32,17/32,23/32,25/32\}$
(note $S^{(31)}=[1,31]$ gives $r_2=1/32$). Writing $\eta_n=Q_n+\theta_n$:
\begin{itemize}
\item For $\theta_n=\varphi-1$: $Q_n=1$ gives $|(32+2q_1')/(32\sqrt5)-1|$,
      which equals $c_{\mathrm{sec}}^{(4)}$ at $q_1'=17$ (when $S_{j-1}=S^{(15)}$),
      equals $39/(16\sqrt5)-1\approx0.090083$ at $q_1'=23$, and exceeds $0.145$
      otherwise. For $2\le Q_n\le b_j-1$, concavity gives the minimum of the
      endpoint limits $(2/\sqrt5)(2-\varphi+r_2)$ (at $Q_n=2$, bounded below by
      $(2/\sqrt5)(2-\varphi+1/32)\approx0.36959$) and $(2/\sqrt5)(2-\varphi+r_1)$
      (at $Q_n=b_j-1$, $\ge0.53729$); for $Q_n\ge b_j+1$, we get
      $\ge(2/\sqrt5)(\varphi-31/32)-o(1)\approx0.58073$.
\item For $\theta_n=2-\varphi$: the three ranges $Q_n=1$, $2\le Q_n\le b_j-1$,
      and $Q_n\ge b_j+1$ yield lower bounds $\ge(2/\sqrt5)(2+\varphi)>3.23$,
      $\ge(2/\sqrt5)(\varphi-1+7/32)\approx0.74844$, and
      $\ge(2/\sqrt5)(3-\varphi-31/32)\approx0.36959$.
\item For $\theta_n\notin\{\varphi-1,2-\varphi\}$: $h(\eta_n)\ge2\sqrt5$ and
      $\theta_n'>0$ (Proposition~\ref{prop:conjugate-positivity}) give lower
      bounds $\ge2\sqrt5\varphi^{-2}\approx1.7082$ ($Q_n=1$),
      $\ge2\sqrt5(7/32)\approx0.97828$ ($2\le Q_n\le b_j-1$), and
      $\ge2\sqrt5(1-31/32)=\sqrt5/16\approx0.139754>c_{\mathrm{sec}}^{(4)}$
      ($Q_n\ge b_j+1$).
\end{itemize}
Every branch exceeds $c_{\mathrm{sec}}^{(4)}$, completing the proof.
\end{proof}

\begin{theorem}[Effective golden-orbit spectrum of rational affine images of $e$]
\label{thm:golden-finite-shift}
Let $x=(A+Ce)/B$, where $B\ge1$, $C\ne0$, and $\gcd(A,B,C)=1$.
Let $\mathcal L_\varphi(x)$ denote the finite limits of
$h(\beta_n)|x-\beta_n|$ along sequences $\beta_n\in E_\varphi$ with
$h(\beta_n)\to\infty$. There is an effectively computable positive integer
$D$ such that
\[
 \mathcal L_\varphi(x)\subseteq
 \left\{\left|1-\frac{m}{D\sqrt5}\right|:m\in\Z\right\}.
\]
For each rational $H>0$, the intersection $\mathcal L_\varphi(x)\cap[0,H]$
is effectively computable by the three families in the proof.
The spectrum has positive minimum and uniform separation between
distinct values. Its first two values are found with $H=3$.
\end{theorem}

\begin{proof}
Use Theorem~\ref{thm:affine-window}. On each residue class, the head
$b_j$ tends to infinity and the suffix $S_j$ is constant with even
length, possibly zero. At a head the complete quotient and reversed
denominator ratio satisfy
\[
 T=b_j+r_j^*,\qquad r_j^*\longrightarrow r=[0;S_j]\in[0,1),
 \qquad z\longrightarrow v=[0;\operatorname{rev}(S_{j-1})]\in[0,1),
\]
where the empty word has value $[0;\varnothing]=0$. These follow by
cutting at the next or previous growing head. For a nonempty suffix,
the limits immediately after its head are $u=[S_j]>1$ and $0$. We may pass to a
subsequence with fixed residue class. If a normalized error has a finite
limit as the height grows, then $\beta_n\to x$, so the first mismatch
index tends to infinity.

\emph{Mismatch inside a suffix.} Cut immediately after the growing
head. The orbit tail $\eta$ differs from $S_j$ within its first
$|S_j|$ digits. At this cut $T\to u$ from below. By
Lemma~\ref{lem:even-suffix-sep}, $|T-\eta|\ge(1-o(1))|u-\eta|$
uniformly over these tails. The conjugate is bounded away from zero
in absolute value by Proposition~\ref{prop:conjugate-positivity}
and the separate value $\varphi'$. Hence \eqref{eq:bugeaud-34}
bounds the normalized error below by $(1-o(1))\Phi_{u,0}(\eta)$.
For a bounded error, Lemma~\ref{lem:general-endpoint-rigidity}
allows only finitely many integers in the endpoint expression and
finitely many corresponding $\eta$ by algebraic inversion.
After taking a constant-tail subsequence, the limit is
$\Phi_{u,0}(\eta)$. Conversely, each such tail with the required suffix
mismatch realizes that limit and has transported height tending to
infinity by \eqref{eq:bugeaud-34}.

\emph{Mismatch at a head.} Write $\eta=q+\theta$, where $q\ge1$ is an
integer different from $b_j$ and $\theta\in E_\varphi\cap(0,1)$. We have
$h(\eta)=h(\theta)\ge2/\sqrt5$ and $\theta'\ge-\varphi>-2$.
Indeed, the primitive quadratic polynomial has discriminant $5$,
and applying Proposition~\ref{prop:conjugate-positivity} to
$1/\theta$ gives $\theta'>0$ except at $\theta=\varphi-1$,
where $\theta'=-\varphi$. If $q\le b_j/2$, then for $b_j\ge8$
\[
 \frac{|T-q-\theta|}{T+z}\ge\tfrac16.
\]
Bounded normalized error forces $q$ to be bounded, since
$|\eta'+z|\ge q-2$ and $h(\theta)\ge2/\sqrt5$. For
$\theta\ne\varphi-1$ the stronger $|\eta'+z|\ge q$ also bounds
$h(\theta)$. For the exceptional tail the height is fixed.
If $q>b_j/2$, then $|\eta'+z|/(T+z)\ge1/6$ for $b_j\ge8$.
The factor $|T-\eta|$ is at least $|q-b_j|-1$, so bounded error
forces $d_0=q-b_j$ to be bounded. In that case the quotient
$|\eta'+z|/(T+z)$ tends to one uniformly in $\theta$, because
$|\theta'-\theta|\le\sqrt5/a$, where $a\ge1$ is the leading coefficient
of the primitive polynomial of $\theta$. For $r>0$ the mismatch
gives $|T-\eta|\ge\frac12\min(r,1-r)$ eventually, hence bounded
height. For $r=0$ this works with an eventual lower bound $1/2$
whenever $d_0\ne-1$ (and $d_0\ne0$). For $r=0$ and $d_0=-1$,
\[
 |T-\eta|=r_j^*+1-\theta\ \ge\ 1-\theta.
\]
The polynomial $F(t)=a(t-\theta)(t-\theta')\in\Z[t]$ has $F(1)\ne0$.
Writing $M=a(1-\theta)>0$ gives
\[
 1\le |F(1)|\le \frac{M(M+\sqrt5)}{a}.
\]
Thus bounded normalized error also bounds $a$ in this boundary case.
All branches now admit a constant-tail subsequence.

An orbit tail $\theta\in(0,1)$ of bounded height has only finitely
many possibilities. If its primitive polynomial is
$a t^2+ct+d$ with $a>0$, then $a=h(\theta)\sqrt5/2$ and
$|c|\le2a+\sqrt5$. Passing to constant data in the two branches
gives the respective limiting expressions
\[
 h(\theta)|q+\theta'+v|,\qquad
 h(\theta)|r-d_0-\theta|,\quad d_0=q-b_j\in\Z\setminus\{0\}.
\]
Both are attained for the displayed fixed choices by
\eqref{eq:bugeaud-34}. In the first branch $q+\theta'+v\ne0$
because $\theta'$ is irrational and $v$ rational, so the transported
height tends to infinity. In the second branch both complete-quotient
factors in the height formula grow without bound.

For clarity, the following are explicit search bounds for limits
at most $H$. In the suffix branch write $u=p/s$ in lowest terms and
let $m$ be the signed integer from the endpoint formula
$\Phi_{u,0}(\eta)=|m/(p\sqrt5)-1|$. Retain
$(1-H)p\sqrt5\le m\le(1+H)p\sqrt5$ and apply
Lemma~\ref{lem:general-endpoint-rigidity} to its signed norm divisors.
Check the resulting orbit tails and the suffix mismatch.
In the fixed-$q$ head branch, $|q+\theta'+v|\ge q-2$ and
$h(\theta)>1/2$ give $q\le 2H+2$. For
$\theta\ne\varphi-1$ the stronger bound $|q+\theta'+v|\ge q$
gives $a\le H\sqrt5/2$. Include $\varphi-1$ separately.
In the moving-offset branch,
$|r-d_0-\theta|\ge |d_0|-1$ gives $|d_0|\le2H+1$.
If $r>0$, then $a\le H\sqrt5/(2\min(r,1-r))$. If $r=0$
and $d_0\ne-1$, the distance is at least $1$, giving $a\le H\sqrt5/2$.
If $r=0$ and $d_0=-1$, put $M=a(1-\theta)\le H\sqrt5/2$ in the
norm bound to obtain $a\le5H(H+2)/4$.

The window phases are obtained from one period of locked rational
convergents, using the expansion with the required parity.
For each head bound enumerate the discriminant-$5$ polynomials
with the stated bounds on $a$ and $c$, and retain the roots in
$(0,1)$ in the orbit. Orbit membership and prefixes are decidable by
exact continued-fraction iteration in $\Q(\sqrt5)$, stopping at
$\varphi$ or a repeated complete quotient. All comparisons and
values are in $\Q(\sqrt5)$. To see the lattice assertion, take $D$ to
be the least common multiple of the numerators of the nonempty
suffix endpoints $u$ and the denominators of the head endpoints $r,v$.
The suffix formula already has the required form. In the two head
formulas use $2a\theta=-c\pm\sqrt5$, $2a\theta'=-c\mp\sqrt5$, and
$h(\theta)=2a/\sqrt5$. This expresses each value as
$|m/(D\sqrt5)-1|$ with integral $m$, changing its sign if necessary.
All lattice values are at least
$\|D\sqrt5\|/(D\sqrt5)>0$, where $\|\cdot\|$ is distance to the
nearest integer. The difference of distinct values is at least
$\min(1,\|2D\sqrt5\|)/(D\sqrt5)>0$ by considering whether their
integers lie on the same side of $D\sqrt5$. The two moving offsets $d_0=1,2$
for $\theta=\varphi-1$ give two distinct values below $3$,
so the first two distinct limits can always be found using $H=3$.
\end{proof}

\begin{corollary}[Full spectrum when the digits diverge]\label{cor:golden-pure-heads}
For every irrational $x$ whose partial quotients $a_n(x)$ tend to
infinity,
\[
 \mathcal L_\varphi(x)=
 \left\{\left|1-\frac{2j+1}{\sqrt5}\right|:j\in\Z\right\}.
\]
The first two values are $3/\sqrt5-1$ and $1-1/\sqrt5$.
In the primitive affine family, digit divergence occurs exactly for
$|C|=1$, $B\in\{2,4\}$ and odd $A$.
\end{corollary}

\begin{proof}
Every late digit is a growing head and both limiting endpoints
are zero. The head proof of Theorem~\ref{thm:golden-finite-shift}
uses only that successive heads tend to infinity in this case.
The coefficient $c$ is odd, since $c^2-4ad=5$, so both head
formulas give values of the displayed form with odd numerator.
Every positive odd numerator $2j+1$ ($j\ge0$) occurs in the fixed
branch by taking $\eta=\varphi+j$. Every negative odd numerator
$-(2j+1)$ occurs in the moving branch by taking
$\theta=\varphi-1$ and $d_0=j+1$. Ordering the odd integers by
their distance to $\sqrt5$ gives the stated first two values.
The affine characterization follows from Corollary~\ref{cor:affine-single}:
a nonempty suffix recurring in a phase contains a fixed digit.
The minimum agrees with Lin~\cite[Lemma~4.2]{Lin2018}; the formula
here describes all finite limits.
\end{proof}

\begin{corollary}[A mixed one-digit and three-digit model]\label{cor:golden-mixed}
For $x=(1+e)/6$, the slice $\mathcal L_\varphi(x)\cap[0,1]$ is
\[
 \left\{\left|1-\frac{m}{3\sqrt5}\right|:
          m\in\{1,3,5,7,9,11,13\}\right\}.
\]
The first two limits are $7/(3\sqrt5)-1$ and $1-\sqrt5/3$.
\end{corollary}
\begin{proof}
Run the exact enumeration of Theorem~\ref{thm:golden-finite-shift}
with $(A,B,C)=(1,6,1)$ and $H=1$. The starting data are the
locked convergents from Theorem~\ref{thm:affine-window}; the
certificate checks one phase period and its translate. It evaluates
the suffix candidates and both head families in $\Q(\sqrt5)$.
The two smallest values are obtained at suffix $(1,2)$ with
$\eta=\varphi$ and $\eta=\varphi+1$, respectively.
\end{proof}

\begin{corollary}[Exact constants beyond denominator four]
\label{cor:golden-finite-shift-examples}
The two smallest distinct golden-orbit limits are as follows.
\[
\begin{array}{c|cc}
x & c_\varphi(x) & \text{second smallest limit}\\ \hline
A/5-e,\ (A,5)=1 & 56/(25\sqrt5)-1 & 1-54/(25\sqrt5)\\
A/6-e,\ (A,6)=1 & 1-79/(36\sqrt5) & 1-77/(36\sqrt5)\\
22/7-e & 110/(49\sqrt5)-1 & 1-108/(49\sqrt5)
\end{array}
\]
\end{corollary}

\begin{proof}
Apply Theorem~\ref{thm:golden-finite-shift} with $H=1$. The accompanying
\texttt{golden\_finite\_spectrum.py} obtains the window phases from the
locked rational convergents and performs both enumerations in
$\Q(\sqrt5)$ using rational arithmetic. Integer translation reduces
the denominator-five and denominator-six cases to the four and two
coprime classes, respectively. The script also checks the first two
values at denominators one through four against the preceding results.
For example, at $22/7-e$ the smallest value is attained by tail
$(5-\sqrt5)/2$ immediately after a growing head with suffix
$(1,2,1,1,1,2,2,1)$; the second value uses tail
$3/2+\sqrt5/10$ at suffix $(1,1,2,1,1,3,1,1)$.
\end{proof}

\paragraph{Verification scope.}
The three inequalities for the growing-head cutoff are checked in
\texttt{GoldenHeadCutoff.lean} with only the base axioms. The finite
spectrum theorem has a paper proof and its tabulated instances use the
exact enumeration certificate. The real inequality used in the empty-suffix
norm cutoff is kernel-checked in \texttt{GoldenAffineEndpoint.lean}, with
the norm and conjugate-separation inequalities as inputs.
The corresponding golden-orbit
approximation property of $\pi$ remains open.

\section{General quadratic orbits, ideal classes, and uniform interval counting}\label{sec:quadratic}

This chapter extends the finite-spectrum analysis from the golden
ratio to arbitrary real quadratic orbits, factors the trace test
through invertible ideal classes, and counts the finite limits
uniformly over admissible intervals.
\subsection{Arbitrary quadratic orbits and counting the finite limits}
Let $\tau$ be a real quadratic irrational with primitive discriminant
$\Delta>0$. Write $E_\tau$ for the orbit of $\tau$ and $\tau'$ under
integer unimodular fractional-linear maps, and write
$\mathcal L_\tau(x)$ for the finite limits of
$h(\beta)|x-\beta|$ along $\beta\in E_\tau$ with $h(\beta)\to\infty$,
where $h(\beta)=2/|\beta-\beta'|$.
Suppose the continued fraction of $x$ is eventually a concatenation of
words $(b_j,S_j)$, where $b_j\to\infty$ and the finite suffix words $S_j$
belong to a finite list. A phase is a pair $(S_{j-1},S_j)$ occurring
infinitely often. We allow empty words.

\begin{theorem}[Finite-phase spectra for arbitrary quadratic orbits]\label{thm:general-quadratic-phase}
Given the recurring phases and $\tau$, the set
$\mathcal L_\tau(x)\cap[0,H]$ is effectively computable for every positive
rational $H$. An explicit positive integer $D$ obtained from the phases
satisfies
\[
 \mathcal L_\tau(x)\subseteq
 \{\,|1-m/(D\sqrt\Delta)|:m\in\Z\,\}.
\]
Consequently the spectrum has positive minimum and a positive bound on
the distance between distinct values. For any $\theta\in E_\tau\cap(0,1)$
and any phase with forward endpoint $r=[0;S_j]$, the spectrum contains
the entire progression
\[
 h(\theta)(d+\theta-r),\qquad d=1,2,\ldots .
\]
In particular, for $N_\tau(x,H)=|\mathcal L_\tau(x)\cap[0,H]|$,
\[
 \max\!\left(0,\left\lfloor H/h(\theta)+r-\theta\right\rfloor\right)
 \le N_\tau(x,H)\le \lfloor 2D\sqrt\Delta\,H\rfloor+1.
\]
The hypotheses and the phase list are effective for every nonconstant
rational affine image of $e$.
\end{theorem}

\begin{proof}
For any orbit element $y$ let $F_y(t)=a(t-y)(t-y')$ be its primitive
integer polynomial with $a>0$. Primitive discriminant is invariant
under unimodular substitution, hence $a|y-y'|=\sqrt\Delta$ and
$h(y)=2a/\sqrt\Delta$. For $s=P/Q\in\Q$ in lowest terms and
$M=a|s-y|$, the integer $Q^2 F_y(s)$ is nonzero. Therefore
\begin{equation}\label{eq:quadratic-norm-cutoff}
  a\le Q^2 M(M+\sqrt\Delta).
\end{equation}
The same bound holds using $M=a|s-y'|$.

Fix a phase, write $v=[0;\operatorname{rev}(S_{j-1})]$, and use $0$
for the empty word. A bounded normalized error with height tending
to infinity implies $\beta\to x$, so the first difference in their
continued fractions goes to infinity. Pass to a fixed phase and to
one of the following three cases.

\emph{Mismatch within a nonempty suffix.}
Cut immediately after $b_j$. The forward quotient tends to
$u=[S_j]\ge1$, and the backward parameter tends to zero. At this cut the
tail $\eta$ of the orbit approximant is outside the interval with
prefix $S_j$. The finite continued-fraction map sends $(1,\infty)$ to
the open interval of this prefix and sends $\infty$ to $u$.
The other endpoint is at fixed positive distance from $u$.
Thus $|T-\eta|\ge(1-o(1))|u-\eta|$ uniformly for such tails, by considering
each side of the interval separately. If $a\to\infty$ then
$\eta'\ge\eta-\sqrt\Delta/a>1/2$ eventually. The identity
\eqref{eq:bugeaud-34}, which holds by the same fractional-linear
calculation for this orbit, bounds $a|u-\eta|$ and
\eqref{eq:quadratic-norm-cutoff} bounds $a$, a contradiction.
Once $a$ is bounded, \eqref{eq:bugeaud-34} also bounds $\eta$: for large $\eta$
both $\eta'+z$ and $|T-\eta|$ grow. A constant-tail subsequence
exists, and its value is $h(\eta)|u-\eta|\,|\eta'|/u$.

\emph{Mismatch at a growing head.}
Write $\eta=q+\theta$, where $q\ge1$ is an integer, $q\ne b_j$,
and $\theta\in E_\tau\cap(0,1)$. At this cut $T=b_j+r_j^*$ with
$r_j^*\to r$ and $0<r_j^*<1$, while $z\to v$.
If $q\le b_j/2$, then $|T-\eta|/(T+z)\ge1/6$ eventually.
Since $\theta'$ is bounded uniformly by the discriminant identity,
bounded error forces bounded $q$. If $a\to\infty$, then
$q+\theta'+z\ge1/2$ eventually, which contradicts bounded error.
A constant pair $(q,\theta)$ consequently gives the limit
$h(\theta)|q+\theta'+v|$.

If $q>b_j/2$, then $|\eta'+z|/(T+z)\ge1/6$ eventually, and
$|T-\eta|\ge|q-b_j|-1$ bounds the offset $d=q-b_j$.
For bounded $d$, that quotient tends uniformly to one. When $r>0$,
the mismatch bounds $|T-\eta|$ away from zero. When $r=0$ the same
holds for $d\ne-1,0$. In the boundary case $r=0,d=-1$,
\[
 |T-\eta|=r_j^*+1-\theta\ge 1-\theta.
\]
Bounded error bounds $a(1-\theta)$ and
\eqref{eq:quadratic-norm-cutoff} at $s=1$ bounds $a$.
Thus in every moving case a constant $(d,\theta)$ gives
$h(\theta)|r-d-\theta|$.

Conversely, each candidate in these three families is attained by
transporting its tail along the chosen phase. Append a tail differing
inside $S_j$ for the suffix family and take $q$ fixed or $q=b_j+d$ for
the head families. Transported heights tend to infinity by the
second identity of \eqref{eq:bugeaud-34}. In the first two families
the pertinent conjugate plus the limiting backward parameter is
nonzero, since the latter is rational. In the moving family both
complete-quotient factors grow.

For an explicit search with limit at most $H$, put $M_0=H\sqrt\Delta/2$.
In the suffix family, if $a\ge 2\sqrt\Delta$ then $\eta'>1/2$ and
$a|u-\eta|\le 2uM_0$, so \eqref{eq:quadratic-norm-cutoff} bounds $a$
in terms of $u,H,\Delta$. If $\eta\ge1+\sqrt\Delta$ then
$\eta'\ge1$ and $\eta\le u+uM_0$. This bounds its integer polynomial
coefficients. In the fixed head family $q\le M_0+\sqrt\Delta$.
If $a\ge2\sqrt\Delta$ then $q+\theta'+v\ge1/2$, hence $a\le2M_0$.
In the moving head family $|d|\le1+M_0$ and a single bound for every
$d$ is $a\le Q_r^2 M_0(M_0+\sqrt\Delta)$, by
\eqref{eq:quadratic-norm-cutoff} at $s=r-d$, where $Q_r$ is the
denominator of $r$. Integer polynomials are enumerated within these
bounds and checked exactly for discriminant, orbit membership, and
the respective tail condition. Membership is decidable by the
continued-fraction algorithm for quadratic irrationals, comparing the
eventual periods up to cyclic shift for $\tau$ and $\tau'$.

To obtain the lattice, take $D$ to be the common multiple of the
numerators of the nonempty suffix endpoints $u$ and the denominators
of all $r,v$, in lowest terms. In the suffix value use
\[
 2a(u-\eta)\eta'/u=2a\eta'-2c/u,\qquad c=a\eta\eta'\in\Z,
 \qquad 2a\eta'=-b\pm\sqrt\Delta .
\]
For each head value use $2a\theta=-b\pm\sqrt\Delta$ and the
corresponding conjugate formula. This gives the stated lattice
inclusion. The positive distance to zero and the positive separation
follow from $\|D\sqrt\Delta\|>0$ and $\|2D\sqrt\Delta\|>0$.
Counting the admissible integers $m$ gives the upper bound.
The choices $d\ge1$ in the moving family give the progression and
the lower count. The affine phase data follow from
Theorem~\ref{thm:affine-window}.
\end{proof}

\begin{corollary}[Closed spectra for constant-period quadratic tails]
\label{cor:metallic-spectrum}
Suppose the partial quotients of $x$ tend to infinity and let
$\tau_a=(a+\sqrt{a^2+4})/2$, $a\in\N_{>0}$. Then
\[
 \mathcal L_{\tau_a}(x)=
 \left\{\,\left|1-\frac{m}{\sqrt{a^2+4}}\right|:
                 m\equiv a\pmod{2}\,\right\},\qquad
 N_{\tau_a}(x,H)=\sqrt{a^2+4}\,H+O(1).
\]
\end{corollary}

\begin{proof}
Both endpoints are zero, and only the two head families occur.
All primitive polynomials of the orbit have discriminant $a^2+4$,
so their middle coefficient has the parity of $a$. This forces the
indicated parity in both head values. Conversely, the fixed family
with $\eta=\tau_a+j$, $j\ge0$, gives $m=a+2j$. The moving family with
$\theta=\tau_a-a$ and $d=j+1$ gives $m=a-2(j+1)$.
The two families cover all admissible $m$. Different integers give
different values because twice the square root is irrational.
Counting the integers of this parity in the appropriate interval
gives the claimed asymptotic.
\end{proof}

\begin{remark}[A nonconstant-period example]
For the orbit of $\tau=(1+\sqrt3)/2=[\overline{1,2}]$ and the
mixed target $x=(1+e)/6$, the exact enumerator gives
\[
 \mathcal L_\tau(x)\cap[0,1/2]
   =\{\,|1-m/(3\sqrt{12})|:m\in\{6,8,10,12,14\}\,\}.
\]
The phase list is the same as for Corollary~\ref{cor:golden-mixed}.
The certificate enumerates primitive discriminant-$12$ polynomials
and tests every admissible tail of the three families by exact
quadratic-field arithmetic.
\end{remark}

Bugeaud's formula \cite[(3.4)]{BugeaudQuad} gives the underlying error
identity. Lin \cite[Lemma~4.2]{Lin2018} computes the minimum for
diverging digits for a class that includes these constant periods.
Corollary~\ref{cor:metallic-spectrum} gives the complete finite-limit
spectrum and its counting law for each such orbit; the preceding
theorem allows any real quadratic orbit and recurring bounded
suffixes. These results have paper proofs.


\subsection{Arithmetic structure and quantitative stability of the orbit spectrum}

\begin{theorem}[Pure-head trace test]\label{thm:pure-trace-test}
Suppose the partial quotients of $x$ tend to infinity, and let $\tau$
have primitive discriminant $\Delta$. Define $T_\tau$ to consist of
the integers $m$ for which there is a positive divisor $A$ of
$|N_m|$, where $N_m=(m^2-\Delta)/4\in\Z$, such that
\[
 \gcd(A,m,N_m/A)=1,\qquad
 A X^2+mX+N_m/A\text{ has a root in }E_\tau.
\]
Then
\[
 \mathcal L_\tau(x)=
 \{\,|1-m/\sqrt\Delta|:m\in T_\tau\,\}.
\]
This is a divisor test for the entire finite-limit spectrum.
\end{theorem}

\begin{proof}
The pure-head cases of Theorem~\ref{thm:general-quadratic-phase}
give $h(\eta)|\eta'|$ for $\eta\in E_\tau$ with $\eta>1$,
and $h(\gamma)|\gamma|$ for $\gamma=d+\theta$ with $d\ne0$
and $0<\theta<1$. The latter class contains all negative elements
of $E_\tau$. The orbit is closed under sign change and conjugation,
and $h(-\gamma)=h(\gamma)$. The two lists together therefore give
exactly all $h(\gamma)|\gamma|$, $\gamma\in E_\tau$.

If the primitive polynomial of $\gamma$ is $A X^2+mX+C$, then
$AC=N_m$ and $2A\gamma=-m\pm\sqrt\Delta$. The two roots are both
in the orbit. Applying $X\mapsto-X$ shows $T_\tau=-T_\tau$,
so both signs are accounted for by the displayed formula.
Conversely, every polynomial in the test supplies these values.
Since $\Delta$ is nonsquare, $N_m\ne0$. This proves the effective
finite divisor test.
\end{proof}

\begin{corollary}[Monic orbits]\label{cor:principal-orbit-spectrum}
If $E_\tau$ contains a root of a monic integer quadratic, then for
every $x$ with diverging partial quotients,
\[
 \mathcal L_\tau(x)
  =\{\,|1-m/\sqrt\Delta|:m\in\Z,\ m\equiv\Delta\pmod2\,\},
 \qquad N_\tau(x,H)=\sqrt\Delta\,H+O(1).
\]
\end{corollary}

\begin{proof}
Fix a monic polynomial with middle coefficient $b$. For any
$m\equiv b\pmod2$, substitution $X\mapsto X+(m-b)/2$ gives a monic
orbit polynomial with middle coefficient $m$. The discriminant
condition gives necessity of $m\equiv\Delta\pmod2$. Apply
Theorem~\ref{thm:pure-trace-test}; counting one parity in the
interval of length $2\sqrt\Delta H$ gives the asymptotic.
\end{proof}

\begin{corollary}[A prime-norm obstruction]\label{cor:prime-norm-orbit}
If $N_m$ is an integer and $|N_m|$ is one or prime, then
$|1-m/\sqrt\Delta|$ belongs to the pure-head spectrum if and only
if $E_\tau$ contains a monic root.
For discriminant $60$, the orbits of $\sqrt{15}$ and
$\sqrt{15}/3$ have respective minima
\[
 8/\sqrt{60}-1,\qquad 1-6/\sqrt{60}.
\]
\end{corollary}

\begin{proof}
Since $\sqrt\Delta$ is irrational, the absolute-value expression
determines $m$ uniquely. In the divisor test either $A=1$ or
$|C|=1$. In the second case replace a root by its reciprocal and
multiply the polynomial by a sign if needed to make it monic.
The converse is Corollary~\ref{cor:principal-orbit-spectrum}.

The continued-fraction periods of $\sqrt{15}$ and $\sqrt{15}/3$
are cyclically $(1,6)$ and $(2,3)$, respectively. The conjugate
periods give the same respective cyclic classes. These period
computations, certified by exact continued fractions, distinguish
the orbits. Any monic root of discriminant $60$ is an integer
translate of $\sqrt{15}$ or its conjugate. Thus $m=8$, for which
$N_8=1$, is excluded in the second orbit. The coefficient $m=6$
is attained by an integer translation of $3X^2-5$.
Among the even integers the closest to $\sqrt{60}$ is $8$,
followed by $6$, giving the asserted minima.
\end{proof}

The prime-factor step in Corollary~\ref{cor:prime-norm-orbit}
is checked in \texttt{OrbitAndFlow.lean}. The certificate
\texttt{verify\_ten\_rounds.py} evaluates the two discriminant-$60$
orbits by both the trace test and the finite-phase enumerator.

\begin{proposition}[Unimodular transport]\label{prop:orbit-mobius}
For a unimodular fractional-linear map $g$ and an irrational target
$x$, one has $\mathcal L_\tau(gx)=\mathcal L_\tau(x)$.
In particular the finite-phase algorithm also computes every
bounded slice at $g((A+Ce)/B)$, $C\ne0$, using the original affine
parameters.
\end{proposition}

\begin{proof}
Write the denominator of $g$ as $cX+d$. A bounded normalized error
and $h(\beta)\to\infty$ imply $\beta,\beta'\to x$. The exact
transformation formula is
\[
 h(g\beta)|gx-g\beta|=
 h(\beta)|x-\beta|\frac{|c\beta'+d|}{|cx+d|},\qquad
 h(g\beta)=h(\beta)|c\beta+d|\,|c\beta'+d|.
\]
The multiplier of the error tends to one, and the height ratio
tends to $|cx+d|^2>0$. Apply the same facts to the inverse map.
\end{proof}

\begin{theorem}[Rate of approach to the finite-phase spectrum]
\label{thm:phase-convergence-rate}
In Theorem~\ref{thm:general-quadratic-phase}, fix a bound $H>0$.
Consider an orbit approximant with normalized error at most $H$
whose first disagreement with the target is in window $j$. For
all sufficiently large $j$, its distance from some member of
$\mathcal L_\tau(x)$ is
\[
 O_{H,x,\tau}\!\left(b_{j-1}^{-1}+b_j^{-1}+b_{j+1}^{-1}\right).
\]
Each limit value has a witnessing tail family with the same rate
on a recurring phase. For a rational affine image of $e$, this
rate is $O(j^{-1})$.
\end{theorem}

\begin{proof}
First the bounded-error families in the proof of
Theorem~\ref{thm:general-quadratic-phase} bound the candidate data
uniformly for all sufficiently large $j$. In detail, if the
positive leading coefficients in the suffix branch were unbounded
along a sequence, the conjugate would stay above $1/2$, and the
norm cutoff at $u$ would bound them. Then the error also bounds
$\eta$. In the first head branch the conjugate separation bounds
$q$ and prevents unbounded $A$. The moving branch bounds $d$;
an unbounded $A$ is excluded by the positive mismatch distance or
by the norm cutoff at $1$ in the exceptional empty-suffix case.
Finitely many phases therefore leave only finitely many actual
tail data for bounded errors. All such late data yield limit
values on their recurring phases.

A fixed finite continued-fraction word gives a fractional-linear
map of the reciprocal of the next complete quotient, regular at
zero. Thus a forward endpoint at a head has error
$O(b_{j+1}^{-1})$, and the backward endpoint has error
$O(b_{j-1}^{-1})$. At the cut after a head, the forward endpoint
has error $O(b_{j+1}^{-1})$ and the backward parameter is
$O(b_j^{-1})$. The error identity \eqref{eq:bugeaud-34}
then gives the asserted rate in the suffix case. In the fixed
head case the forward factor $|T-\eta|/(T+z)$ equals
$1+O(b_j^{-1})$; in the moving head case the conjugate factor
$|\eta'+z|/(T+z)$ equals $1+O(b_j^{-1})$. Combining these estimates
gives the same rate in each case, with uniform constants over the
finite candidate data. For any limit, the fixed witnessing data
give the rate along their phase by the same formulas.
Finally Theorem~\ref{thm:affine-window} gives $b_j\asymp j$
for affine images of $e$.
\end{proof}


\begin{lemma}[General Case~2 exclusion criterion for periodic affine windows]\label{lem:general-case2-criterion}
Let $x=(A+Ce)/B$ ($B\in\N_{>0}$, $C\in\Z\setminus\{0\}$) have eventual
windows $W_j=(b_j,S_j)$ for $j>L=\max(B,|C|)^2$, where $b_j\to\infty$
and every suffix $S_j$ has even length $2t_j\ge0$ (Theorem~\ref{thm:affine-window}),
drawn from a finite periodic set of pairs $(D_j,S_j)$. Set
$r_1(j):=0$ if $D_j=1$ ($S_j=\emptyset$) and $r_1(j):=1/[S_j]=q_{0,j}/D_j\in(0,1)$
if $D_j>1$; likewise set $r_2(j-1):=0$ if $D_{j-1}=1$ and
$r_2(j-1):=[0;\mathrm{rev}(S_{j-1})]=q'_{1,j-1}/D_{j-1}\in(0,1)$ if $D_{j-1}>1$.
Put $r_{1,\min}:=\min_j r_1(j)$, $r_{1,\max}:=\max_j r_1(j)$,
$r_{2,\min}:=\min_j r_2(j-1)$, and (when $\max_j D_j>1$)
$r_{1,\min}^+:=\min\{r_1(j):D_j>1\}>0$.

Suppose $(\beta_n)\subset E_\varphi$ has $h(\beta_n)\to\infty$ and $\beta_n\to x$,
and pass to a subsequence on which the first mismatch $m_n=k(j)+1$ occurs
at the growing head $b_j$ ($j\to\infty$). Write $\beta_n=M_{m_n}(\eta_n)$
with $\eta_n=Q_n+\theta_n$, $Q_n\in\Z_{\ge1}\setminus\{b_j\}$,
$\theta_n\in E_\varphi\cap(0,1)$, and $E_n:=h(\beta_n)|x-\beta_n|$. Then:
\begin{enumerate}
\item[(a)] If $\theta_n=\varphi-1$ ($\theta_n'=-\varphi$, $h(\eta_n)=2/\sqrt5$):
  \begin{itemize}
  \item $Q_n=1$ gives $\beta_n=[a_0;\ldots,a_{k(j)},\overline1]$ and
        $E_n=\bigl|(D_{j-1}+2q'_{1,j-1})/(D_{j-1}\sqrt5)-1\bigr|+o(1)$
        (reducing to $1-1/\sqrt5$ when $D_{j-1}=1$, and matching the Case~1
        value $\Phi_{D_{j-1}/q_{0,j-1},0}([S_{j-1};\overline1])$ when $D_{j-1}>1$);
  \item $2\le Q_n\le b_j-1$ gives
        $E_n\ge\frac{2}{\sqrt5}\bigl(2-\varphi+\min(r_{1,\min},r_{2,\min})\bigr)-o(1)$;
  \item $Q_n\ge b_j+1$ gives
        $E_n\ge\frac{2}{\sqrt5}(\varphi-r_{1,\max})-o(1)$.
  \end{itemize}
\item[(b)] If $\theta_n=2-\varphi$ ($\theta_n'=1+\varphi$, $h(\eta_n)=2/\sqrt5$):
  \begin{itemize}
  \item $Q_n=1$ gives $E_n\ge\frac{2}{\sqrt5}(2+\varphi+r_{2,\min})-o(1)\ge 1+\sqrt5-o(1)$;
  \item $2\le Q_n\le b_j-1$ gives
        $E_n\ge\frac{2}{\sqrt5}(\varphi-1+r_{1,\min})-o(1)\ge 1-1/\sqrt5-o(1)$;
  \item $Q_n\ge b_j+1$ gives
        $E_n\ge\frac{2}{\sqrt5}(3-\varphi-r_{1,\max})-o(1)$.
  \end{itemize}
\item[(c)] If $\theta_n\notin\{\varphi-1,2-\varphi\}$, then $h(\eta_n)\ge2\sqrt5$
      and $\theta_n'\in(0,\varphi^2)$:
  \begin{itemize}
  \item $Q_n=1$ gives $E_n\ge2\sqrt5(\varphi^{-2}+r_{2,\min})-o(1)\ge3\sqrt5-5-o(1)$;
  \item $Q_n\ge2$ gives $E_n\ge2\sqrt5\min(r_{1,\min}^+,1-r_{1,\max})-o(1)$
        whenever $D_j>1$, and $E_n\ge(\sqrt5-1)-o(1)$ whenever $D_j=1$.
  \end{itemize}
\end{enumerate}
\end{lemma}
\begin{proof}
Because $S_j$ and $S_{j-1}$ have even length ($\ge0$), cutting at $b_j$ gives
$T_{m_n}=b_j+r_1(j)+\varepsilon_{1,j}$ and
$z_{m_n-1}=r_2(j-1)+\varepsilon_{2,j}$ with $\varepsilon_{1,j},\varepsilon_{2,j}>0$
tending to~$0$ as $j\to\infty$ (monotonicity of even-length continued fractions).
By the discriminant-$5$ height dichotomy, either
$\theta_n\in\{\varphi-1,2-\varphi\}$ with $h(\eta_n)=2/\sqrt5$, or
$h(\eta_n)\ge2\sqrt5$ and Proposition~\ref{prop:conjugate-positivity} applied
to $1/\theta_n\in E_\varphi\cap(1,\infty)\setminus\{\varphi\}$ gives
$\theta_n'\in(0,\varphi^2)$. Applying~\eqref{eq:bugeaud-34} at $m=m_n$:

In~(a) ($\theta_n=\varphi-1$), $Q_n=1$ yields
$(2/\sqrt5)|1-\varphi+r_2(j-1)|+o(1)=|(D_{j-1}+2q'_{1,j-1})/(D_{j-1}\sqrt5)-1|+o(1)$,
and when $D_{j-1}>1$ the expansion $[a_0;\ldots,b_{j-1},S_{j-1},\overline1]$
equals $M_{k(j-1)+2}([S_{j-1};\overline1])$, matching
$\Phi_{D_{j-1}/q_{0,j-1},0}([S_{j-1};\overline1])$. For $2\le Q_n\le b_j-1$,
concavity of $(b_j+r_1(j)+\varepsilon_{1,j}-\varphi+1-Q_n)(Q_n-\varphi+r_2(j-1)+\varepsilon_{2,j})$
on $[2,b_j-1]$ bounds $E_n$ below by the smaller of its endpoint limits
$(2/\sqrt5)(2-\varphi+r_2(j-1))$ (at $Q_n=2$) and
$(2/\sqrt5)(2-\varphi+r_1(j))$ (at $Q_n=b_j-1$). For $Q_n\ge b_j+1$,
$|T_{m_n}-\eta_n|\ge\varphi-r_1(j)-o(1)$ and $(\eta_n'+z_{m_n-1})/(T_{m_n}+z_{m_n-1})\ge1-o(1)$.

In~(b) ($\theta_n=2-\varphi$), $Q_n=1$ gives $(2/\sqrt5)(2+\varphi+r_2(j-1))+o(1)$,
$2\le Q_n\le b_j-1$ gives by concavity the minimum of the endpoint limits
$(2/\sqrt5)(3+\varphi+r_2(j-1))$ and $(2/\sqrt5)(\varphi-1+r_1(j))$ (always the latter),
and $Q_n\ge b_j+1$ gives $|T_{m_n}-\eta_n|\ge3-\varphi-r_1(j)-o(1)$.

In~(c) ($\theta_n\notin\{\varphi-1,2-\varphi\}$), $Q_n=1$ places
$\eta_n=1+\theta_n\in E_\varphi\cap(1,\infty)\setminus\{\varphi\}$, so
Proposition~\ref{prop:conjugate-positivity} gives $\eta_n'\ge\varphi^{-2}$ and
$E_n\ge2\sqrt5(\varphi^{-2}+r_{2,\min})-o(1)$. For $Q_n\ge2$ and $D_j>1$:
at $Q_n=b_j-1$ we have $T_{m_n}-\eta_n>r_1(j)$ and
$(\eta_n'+z_{m_n-1})/(T_{m_n}+z_{m_n-1})\ge1-o(1)$; on $2\le Q_n\le b_j-2$,
concavity of $(b_j-Q_n-1+r_1(j))Q_n$ gives ratio $\ge(1+r_1(j))-o(1)>r_1(j)$;
and at $Q_n\ge b_j+1$, $\eta_n-T_{m_n}\ge1-r_1(j)-o(1)$ with ratio $\ge1-o(1)$.
Hence $E_n\ge2\sqrt5\min(r_1(j),1-r_1(j))-o(1)$. Finally, if $D_j=1$ ($r_1(j)=0$),
then $2\le Q_n\le b_j-2$ and $Q_n\ge b_j+1$ both give $|T_{m_n}-\eta_n|\ge1-o(1)$
and $E_n\ge2\sqrt5-o(1)$, while at $Q_n=b_j-1$ we have
$E_n\ge(1-o(1))h(\theta_n)(1-\theta_n)$. If $\theta_n\le1/2$ this is $\ge\sqrt5$;
if $\theta_n>1/2$, writing $\theta_n=1/(1+\psi)$ with
$\psi=(a_1\varphi+b_1)/(c_1\varphi+d_1)\in E_\varphi\cap(0,1)\setminus\{\varphi-1\}$
gives $h(\theta_n)(1-\theta_n)=h(\psi)\psi(1+\psi')=2|\xi_1\xi_2'|/\sqrt5=|m/\sqrt5-1|$
for $\xi_1=a_1\varphi+b_1$, $\xi_2=(a_1+c_1)\varphi+(b_1+d_1)$, and an odd
integer $m$ with $m^2-5=4N(\xi_1)N(\xi_2)$. Since $|m|\in\{1,3\}$ forces
$|N(\xi_1)|=|N(\xi_2)|=1$ and $\xi_1/\xi_2=1-\theta_n\in\{\varphi^{-1},\varphi^{-2}\}$
(i.e., $\theta_n\in\{2-\varphi,\varphi-1\}$, excluded), we must have $|m|\ge5$
and hence $h(\theta_n)(1-\theta_n)\ge\sqrt5-1$ (attained at $\theta_n=[0;1,2,\overline1]$).
\end{proof}

\begin{theorem}[Exact golden-ratio orbit constant, extremal classification, and second limit point for $B=5$]\label{thm:golden-b5-exact}
For every integer $A$ with $5\nmid A$, write $x:=A/5-e=[a_0;a_1,a_2,\ldots]$. Then:
\begin{enumerate}
\item For all $j\ge0$, $Q_j\bmod5$ has minimal period~$10$ with $Q_{j+5}\equiv-Q_j\pmod5$
      and $5\nmid Q_j$, while $2\nmid Q_j$; hence $G_j=1$ and $D_j=50$ for $j>25$.
      Up to $\mathrm{GL}_2(\Z)$-equivalence, $x$ reduces to $16/5-e$, $17/5-e$,
      $18/5-e$, or $19/5-e$. In all four classes, $(S_j,r_j)$ repeats with
      minimal period $T=50$ and $b_{j+50}=b_j+4$, exhibiting the exact same
      $20$-element suffix set $S^{(q_0)}$ ($[S^{(q_0)}]=50/q_0$) as $q_0$ ranges
      over $(\Z/50\Z)^\times$.
\item For every $A$ with $5\nmid A$,
      \[
       c_\varphi(A/5-e)=\frac{56}{25\sqrt5}-1\approx 0.0017584539199058,
      \]
      attained with limit $\frac{56}{25\sqrt5}-1$ along---and only along---four
      pairwise-disjoint families: $[a_0;\ldots,b_j,\overline1]$ (at $S^{(31)}=[1,1,1,1,1,2,1,1]$),
      $[a_0;\ldots,b_j,1,1,2,\overline1]$ (at $S^{(29)}=[1,1,2,1,1,1,1,1]$),
      $[a_0;\ldots,b_j,2,\overline1]$ (at $S^{(19)}=[2,1,1,1,2,2]$), and
      $[a_0;\ldots,b_j,2,2,\overline1]$ (at $S^{(21)}=[2,2,1,1,1,2]$).
\item Every sequence $(\beta_n)\subset E_\varphi$ with $h(\beta_n)\to\infty$
      eventually disjoint from these four extremal families satisfies
      \[
       \liminf_{n\to\infty}h(\beta_n)|x-\beta_n|\ge
       c_{\mathrm{sec}}^{(5)}:=1-\frac{54}{25\sqrt5}\approx 0.0340186337200909,
      \]
      and $c_{\mathrm{sec}}^{(5)}$ is sharply attained in every class along
      the eight families $[a_0;\ldots,b_j,1,1,1,1,1,2,\overline1]$ ($S^{(31)}$),
      $[a_0;\ldots,b_j,\overline1]$ ($S^{(29)}$),
      $[a_0;\ldots,b_j,1,3,\overline1]$ ($S^{(39)}=[1,3,1,1,4,1]$),
      $[a_0;\ldots,b_j,1,4,\overline1]$ ($S^{(41)}=[1,4,1,1,3,1]$),
      $[a_0;\ldots,b_j,2,1,1,1,2,\overline1]$ ($S^{(19)}$),
      $[a_0;\ldots,b_j,2,\overline1]$ ($S^{(21)}$),
      $[a_0;\ldots,b_j,4,\overline1]$ ($S^{(11)}=[4,1,1,5]$), and
      $[a_0;\ldots,b_j,5,\overline1]$ ($S^{(9)}=[5,1,1,4]$).
\end{enumerate}
\end{theorem}
\begin{proof}
Because $(Q_0,\ldots,Q_9)\bmod5=(1,1,2,1,1,4,4,3,4,4)$ and the recurrence
multiplier $4j+2$ is $5$-periodic, induction gives $Q_{j+5}\equiv-Q_j\pmod5$
and $5\nmid Q_j$ for all $j\ge0$. Since $2\nmid Q_j$, $G_j=1$ and $D_j=50$ for
$j>25$ by Theorem~\ref{thm:affine-window}. Integer shifts reduce to
$16/5-e,\ldots,19/5-e$, where computing one period of $50$ windows confirms
minimal period~$50$, $b_{j+50}=b_j+4$, and the shared $20$-element suffix set
over $(\Z/50\Z)^\times$.

At $D=50$ ($d=25$), Lemma~\ref{lem:general-endpoint-rigidity}(2) forces $m$ to
be even. Since $25\sqrt5\approx55.901699$, the closest even integer is $m=56$
and the second closest is $m=54$ (as $1-54/(25\sqrt5)<58/(25\sqrt5)-1$).
Applying Lemma~\ref{lem:general-endpoint-rigidity}(4) across all $20$ suffixes
at $m=56$ ($56^2-5\cdot25^2=11$) yields exclusively the four stated extremal
solutions, and at $m=54$ ($54^2-5\cdot25^2=-209$) yields exclusively the eight
stated second-limit-point solutions. Evaluating~\eqref{eq:bugeaud-34} proves
the upper bound in~(2) and the sharp attainment in~(3).

For the lower bound in~(3), let $(\beta_n)\subset E_\varphi$ be eventually
disjoint from the four extremal families with $E_n\to L<c_{\mathrm{sec}}^{(5)}$
and mismatch index $m_n\to\infty$.
In \emph{Case 1} ($m_n$ inside $S_j$), Lemma~\ref{lem:even-suffix-sep},
Proposition~\ref{prop:conjugate-positivity} (noting
$\Phi_{50/q_0,0}(\varphi)=|(25+q_0)/(25\sqrt5)-1|$ equals
$56/(25\sqrt5)-1$ only at $q_0=31$, excluded, equals $c_{\mathrm{sec}}^{(5)}$
at $q_0=29$, and exceeds $0.03753$ otherwise), and the exclusion of $m=56$
force $E_n\ge(1-o(1))c_{\mathrm{sec}}^{(5)}$.
In \emph{Case 2} ($m_n=k(j)+1$), apply Lemma~\ref{lem:general-case2-criterion}
with $r_{1,\min}=r_{2,\min}=1/50$ and $r_{1,\max}=49/50$:
in~(a) ($\theta_n=\varphi-1$), $Q_n=1$ gives $|(25+q_1')/(25\sqrt5)-1|$
($q_1'=31$ occurs iff $S_{j-1}=S^{(29)}$, producing the excluded extremal
family $[a_0;\ldots,b_{j-1},1,1,2,\overline1]$, while $q_1'=29$ gives
$c_{\mathrm{sec}}^{(5)}$ and all other $q_1'$ exceed $0.03753$),
$2\le Q_n\le b_j-1$ gives $\ge(2/\sqrt5)(2-\varphi+1/50)\approx0.359529$,
and $Q_n\ge b_j+1$ gives $\ge(2/\sqrt5)(\varphi-49/50)\approx0.570675$;
in~(b) ($\theta_n=2-\varphi$), the three ranges yield $\ge3.25395$, $\ge0.570675$,
and $\ge0.359529$; in~(c) ($\theta_n\notin\{\varphi-1,2-\varphi\}$), $Q_n=1$
gives $\ge2\sqrt5(\varphi^{-2}+1/50)\approx1.79765$ and $Q_n\ge2$ gives
$\ge2\sqrt5/50=\sqrt5/25\approx0.0894427>c_{\mathrm{sec}}^{(5)}$.
All branches contradict $L<c_{\mathrm{sec}}^{(5)}$.
\end{proof}

\begin{theorem}[Exact golden-ratio orbit constant, inert and mod-$8$ exclusions, and second limit point for $B=6$]\label{thm:golden-b6-exact}
For every integer $A$ with $\gcd(A,6)=1$, write $x:=A/6-e=[a_0;a_1,a_2,\ldots]$. Then:
\begin{enumerate}
\item Since $\gcd(Q_j,6)=1$ for all $j\ge0$, we have $G_j=1$ and $D_j=72$ for
      $j>36$. Up to $\mathrm{GL}_2(\Z)$-equivalence, $x$ reduces to $19/6-e$
      ($A\equiv1\pmod6$) or $23/6-e$ ($A\equiv5\pmod6$). In both classes,
      starting at $j=37$, $(S_j,r_j)$ repeats with minimal period~$18$ and
      $b_{j+18}=b_j+1$ (so $b_{j+72}=b_j+4$), exhibiting the exact same
      $12$-element suffix set $S^{(q_0)}$ ($[S^{(q_0)}]=72/q_0$) with
      \[
       q_0\in\{5,11,13,19,29,35,37,43,53,59,61,67\}.
      \]
      Every occurring $q_0$ satisfies $q_0\bmod24\in\{5,11,13,19\}$, and in
      particular $q_0\equiv\pm3\pmod8$.
\item For every $A$ with $\gcd(A,6)=1$,
      \[
       c_\varphi(A/6-e)=1-\frac{79}{36\sqrt5}\approx 0.0186146098750923,
      \]
      attained along---and only along---four pairwise-disjoint families:
      $[a_0;\ldots,b_j,\overline1]$ (at $S^{(43)}=[1,1,2,14]$),
      $[a_0;\ldots,b_j,1,13,2,\overline1]$ (at $S^{(67)}=[1,13,2,2]$),
      $[a_0;\ldots,b_j,2,\overline1]$ (at $S^{(29)}=[2,2,13,1]$), and
      $[a_0;\ldots,b_j,14,2,\overline1]$ (at $S^{(5)}=[14,2,1,1]$).
\item \emph{Inert-$3$ and mod-$8$ exclusions of $m=81,83$.} At $D=72$ ($d=36$),
      $m$ is odd, the inert prime $3\mid36$ excludes $m=81$ by
      Lemma~\ref{lem:general-endpoint-rigidity}(2), and since
      $83^2-5\cdot36^2=409\equiv1\pmod8$ is prime, every divisor
      $N(\gamma)\in\{\pm1,\pm409\}$ satisfies $N(\gamma)\equiv\pm1\pmod8$;
      together with $q_0\equiv\pm3\pmod8$, Lemma~\ref{lem:general-endpoint-rigidity}(3)
      forces $m\equiv36-\varepsilon q_0N(\gamma)\equiv4\pm3\equiv\pm1\pmod8$,
      excluding $m=83\equiv3\pmod8$ completely.
\item Every sequence $(\beta_n)\subset E_\varphi$ with $h(\beta_n)\to\infty$
      eventually disjoint from the four extremal families satisfies
      \[
       \liminf_{n\to\infty}h(\beta_n)|x-\beta_n|\ge
       c_{\mathrm{sec}}^{(6)}:=1-\frac{77}{36\sqrt5}\approx 0.0434598096250900,
      \]
      and $c_{\mathrm{sec}}^{(6)}$ is sharply attained in both classes along the
      four families $[a_0;\ldots,b_j,1,4,\overline1]$ ($S^{(59)}=[1,4,1,1,5,1]$),
      $[a_0;\ldots,b_j,1,5,\overline1]$ ($S^{(61)}=[1,5,1,1,4,1]$),
      $[a_0;\ldots,b_j,5,\overline1]$ ($S^{(13)}=[5,1,1,6]$), and
      $[a_0;\ldots,b_j,6,\overline1]$ ($S^{(11)}=[6,1,1,5]$).
\end{enumerate}
\end{theorem}
\begin{proof}
Since $2\nmid Q_j$ and $3\nmid Q_j$ (Theorem~\ref{thm:golden-b3-exact}(1)),
$\gcd(A,6)=1$ gives $G_j=1$ and $D_j=72$ ($j>36$). Direct verification of
$W_{36},\ldots,W_{54}$ shows $S_{54}=S_{36}$ (with $b_{54}=b_{36}+1$) and
the twelve stated suffixes over $W_{37},\ldots,W_{54}$; because $4\cdot18=72$,
the recurrence $C_j=(4j+2)-F_{j-1}$ from the proof of Theorem~\ref{thm:affine-window}
yields $C_{j+18}=C_j+72$ inductively for all $j\ge37$, giving minimal
period~$18$ for $(S_j,r_j)$ with $b_{j+18}=b_j+1$.

At $D=72$ ($d=36$), $36\sqrt5\approx80.498447$. By~(3), both $m=81$ and $m=83$
are impossible, leaving $m=79$ as the closest admissible odd integer and
$m=77$ as the second closest (since $1-77/(36\sqrt5)<85/(36\sqrt5)-1$).
Lemma~\ref{lem:general-endpoint-rigidity}(4) at $m=79$ ($79^2-5\cdot36^2=-239$)
and $m=77$ ($77^2-5\cdot36^2=-551$) produces precisely the four families in~(2)
and the four families in~(4), respectively.

For the lower bound in~(4), \emph{Case 1} ($m_n$ inside $S_j$) closes via
Lemma~\ref{lem:even-suffix-sep}, Proposition~\ref{prop:conjugate-positivity}
(noting $\Phi_{72/q_0,0}(\varphi)=|(36+q_0)/(36\sqrt5)-1|$ equals $1-79/(36\sqrt5)$
at $q_0=43$, excluded, and is $\ge1-73/(36\sqrt5)\approx0.0931502>c_{\mathrm{sec}}^{(6)}$
for all other $q_0$), and the exclusion of $m=81,83,79$.
In \emph{Case 2} ($m_n=k(j)+1$), Lemma~\ref{lem:general-case2-criterion} with
$r_{1,\min}=r_{2,\min}=5/72$ and $r_{1,\max}=67/72$ gives:
in~(a) ($\theta_n=\varphi-1$), $Q_n=1$ has $q_1'=43$ iff $S_{j-1}=S^{(5)}$
(producing the excluded extremal family $[a_0;\ldots,b_{j-1},14,2,\overline1]$)
and $\ge1-73/(36\sqrt5)\approx0.0931502$ for all other $q_1'$, while
$2\le Q_n\le b_j-1$ gives $\ge0.403754$ and $Q_n\ge b_j+1$ gives $\ge0.614899$;
in~(b) ($\theta_n=2-\varphi$), the three ranges give $\ge3.29818$, $\ge0.614899$,
and $\ge0.403754$; in~(c) ($\theta_n\notin\{\varphi-1,2-\varphi\}$), $Q_n=1$
gives $\ge2.01877$ and $Q_n\ge2$ gives $\ge5\sqrt5/36\approx0.310565>c_{\mathrm{sec}}^{(6)}$.
Every branch exceeds $c_{\mathrm{sec}}^{(6)}$.
\end{proof}

\begin{theorem}[Exact constants, extremal families, and second limit points for multi-scale affine targets]\label{thm:golden-multiscale-exact}
Complete Case~1 and Case~2 classifications hold for the following mixed-$D_j$
rational affine targets:
\begin{enumerate}
\item[(i)] \emph{Target $x=(1+e)/6$ ($T=12$, minimal period~$6$, single-digit phases):}
      for $j>36$, $(G_j,D_j)$ repeats with period~$6$, taking $D_j=3$
      ($S_j=[1,2]$, $r_j=3$) for $j\equiv0,1\pmod6$ and $D_j=1$ ($S_j=\emptyset$,
      $r_j=1$) for $j\equiv2,3,4,5\pmod6$, with $b_{j+6}=b_j+2G_j^2$
      ($b_{j+12}=b_j+4G_j^2$, $G_j\in\{2,6\}$). We have
      \[
       c_\varphi\bigl((1+e)/6\bigr)=\frac{7}{3\sqrt5}-1\approx 0.043498389500,
      \]
      uniquely attained at $j\equiv0,1\pmod6$ by $[a_0;\ldots,b_j,\overline1]$.
      Every sequence $(\beta_n)\subset E_\varphi$ eventually disjoint from this
      family satisfies $\liminf_{n\to\infty}h(\beta_n)|x-\beta_n|\ge 1-\frac{\sqrt5}{3}\approx 0.2546440075$,
      sharply attained at $j\equiv0,1\pmod6$ by $[a_0;\ldots,b_j,2,\overline1]$
      and $[a_0;\ldots,b_j,1,2,\overline1]$.
\item[(ii)] \emph{Target $x=(1+e)/3$ ($T=6$, $D_j\in\{2,6\}$):}
      for $j>9$, $(G_j,D_j)$ repeats with period~$6$ ($b_{j+6}=b_j+4G_j^2$,
      $G_j\in\{1,3\}$), giving $D_j=6$ ($S_j=[1,5]$) at $j\equiv0,1\pmod6$ and
      $D_j=2$ ($S_j=[1,1]$) at $j\equiv2,3,4,5\pmod6$. We have
      \[
       c_\varphi\bigl((1+e)/3\bigr)=1-\frac{2}{\sqrt5}\approx 0.105572809000,
      \]
      attained at $D_j=2$ by $[a_0;\ldots,b_j,\overline1]$ and
      $[a_0;\ldots,b_j,2,\overline1]$; because inert $3\mid3$ excludes $m=6$
      at $D=6$, the second smallest limit point is
      $\frac{8}{3\sqrt5}-1\approx 0.192569588000$, sharply attained at $D_j=6$
      by $[a_0;\ldots,b_j,\overline1]$ and $[a_0;\ldots,b_j,1,4,\overline1]$.
\item[(iii)] \emph{Target $x=(1+2e)/3$ ($T=12$, minimal period~$6$, $D_j\in\{4,12\}$):}
      for $j>9$, $(G_j,D_j,S_j)$ repeats with period~$6$ ($b_{j+6}=b_j+2G_j^2$,
      $G_j\in\{1,3\}$), with $D_j=12$ at $j\equiv3,4\pmod6$ ($S_j=[1,11]$ and $[2,2,1,1]$)
      and $D_j=4$ at $j\equiv0,2\pmod6$ ($S_j=[1,3]$) and $j\equiv1,5\pmod6$
      ($S_j=[3,1]$). We have
      \[
       c_\varphi\bigl((1+2e)/3\bigr)=1-\frac{13}{6\sqrt5}\approx 0.031037209750>0
      \]
      (since inert $3\mid6$ excludes $m=15$ at $D=12$), attained at
      $S_j=[2,2,1,1]$ ($j\equiv4\pmod6$) by $[a_0;\ldots,b_j,2,\overline1]$ and
      $[a_0;\ldots,b_j,2,2,\overline1]$, and second smallest limit point
      $\frac{\sqrt5}{2}-1\approx 0.118033988750$, sharply attained at $D_j=4$
      by $[a_0;\ldots,b_j,\overline1]$ and $[a_0;\ldots,b_j,1,2,\overline1]$
      ($S_j=[1,3]$) and by $[a_0;\ldots,b_j,2,\overline1]$ and
      $[a_0;\ldots,b_j,3,\overline1]$ ($S_j=[3,1]$).
\item[(iv)] \emph{Target $x=22/7-e$ ($T=98$, three-scale $D_j\in\{2,14,98\}$):}
      for $j>49$, $W_j$ has period~$98$ ($b_{j+98}=b_j+4G_j^2$, $G_j\in\{1,7,49\}$)
      with $D_j$-frequencies $\{2:8,\,14:48,\,98:42\}$ per period ($8$ windows
      of $S_j=[1,1]$ at $j\equiv59,60,85,86,120,121,122,123\pmod{98}$, $6$
      suffixes at $D_j=14$, and $14$ suffixes at $D_j=98$). Because $7\equiv2\pmod5$
      is inert, $7\nmid m$ whenever $D_j\in\{14,98\}$. We have
      \[
       c_\varphi(22/7-e)=\frac{110}{49\sqrt5}-1\approx 0.003948887857,
      \]
      attained at $D_j=98$ along the four families
      $[a_0;\ldots,b_j,1,2,\overline1]$ ($S^{(71)}=[1,2,1,1,1,2,2,1]$),
      $[a_0;\ldots,b_j,1,2,2,\overline1]$ ($S^{(69)}=[1,2,2,1,1,1,2,1]$),
      $[a_0;\ldots,b_j,3,\overline1]$ ($S^{(27)}=[3,1,1,1,2,3]$), and
      $[a_0;\ldots,b_j,3,2,\overline1]$ ($S^{(29)}=[3,2,1,1,1,3]$), and
      second smallest limit point
      $c_{\mathrm{sec}}(22/7-e)=1-\frac{108}{49\sqrt5}\approx 0.014304728286$
      (since $7\mid112$), sharply attained at $D_j=98$ along the four families
      $[a_0;\ldots,b_j,1,1,2,\overline1]$ ($S^{(57)}$),
      $[a_0;\ldots,b_j,1,1,3,\overline1]$ ($S^{(55)}$),
      $[a_0;\ldots,b_j,2,2,\overline1]$ ($S^{(41)}$), and
      $[a_0;\ldots,b_j,2,3,\overline1]$ ($S^{(43)}$).
\end{enumerate}
\end{theorem}
\begin{proof}
For each target, Theorem~\ref{thm:affine-window} and finite computation over one
period confirm the stated $(G_j,D_j,S_j)$ lists, suffix evenness ($|S_j|=0$
for $D_j=1$ in~(i), and $|S_j|\ge2$ even in all other windows), and unbounded
head growth $b_j\to\infty$. In Case~1 ($m_n$ inside a non-empty suffix $S_j$),
Lemma~\ref{lem:even-suffix-sep}, Proposition~\ref{prop:conjugate-positivity},
and Lemma~\ref{lem:general-endpoint-rigidity} close every scale:
\begin{itemize}
\item In~(i), $D_j=1$ has no Case~1 suffix; at $D_j=3$ ($S=[1,2]$, $u=3/2$),
      inert $3\mid3$ gives $3\nmid M$ for odd $M$, so $M=7$ yields uniquely
      $\eta=\varphi$, $M=9$ is excluded, and $M=5$ yields $[2;\overline1]$ and
      $[1;2,\overline1]$ (the latter matching $[1,2]$ and diverging at $b_{j+1}$,
      i.e., Case~2 $Q_n=1$ of window $j+1$).
\item In~(ii), at $D=2$ ($d=1$, $S=[1,1]$), $m=2$ gives $1-2/\sqrt5$ via
      $\varphi$ and $[2;\overline1]$ (with $m=4$ giving $4/\sqrt5-1>0.788$);
      at $D=6$ ($d=3$, $S=[1,5]$), inert $3\mid3$ forbids even $m=6$, and
      $m=8$ gives $8/(3\sqrt5)-1$ via $\varphi$ and $[1;4,\overline1]$.
\item In~(iii), at $D=12$ ($d=6$), odd $m$ with $3\nmid m$ forbids $m=15$ and
      makes $m=13$ closest (exclusively at $S=[2,2,1,1]$ via $[2;\overline1]$
      and $[2;2,\overline1]$), with next $D=12$ value at $m=11$
      ($1-11/(6\sqrt5)\approx0.180108$); at $D=4$ ($d=2$), $m=5$ yields the
      four stated $\sqrt5/2-1\approx0.118034$ families at $[1,3]$ and $[3,1]$.
\item In~(iv), at $D=2$ ($d=1$, $S=[1,1]$), $m$ is even (excluding $m=1,3,5$
      by parity) with smallest values at $m=2$ ($1-2/\sqrt5\approx0.105573$)
      and $m=4$ ($4/\sqrt5-1\approx0.788854$); at $D=14$ ($d=7$), inert $7\mid7$
      excludes $m=14$, leaving $m\ge16$ ($16/(7\sqrt5)-1\approx0.0222025>c_{\mathrm{sec}}$);
      at $D=98$ ($d=49$), $7\mid49$ excludes $m=112$, $m=110$ gives uniquely
      the four extremal families ($110^2-5\cdot49^2=95$), and $m=108$ gives
      uniquely the four $c_{\mathrm{sec}}$ families ($108^2-5\cdot49^2=-341$),
      while $\Phi_{D/q_0,0}(\varphi)\ge1-106/(49\sqrt5)\approx0.0325583$ at
      $D=98$.
\end{itemize}
In Case~2 ($m_n=k(j)+1$), Lemma~\ref{lem:general-case2-criterion} closes every
target:
\begin{itemize}
\item In~(i) ($r_{1,\min}=r_{2,\min}=0$, $r_{1,\min}^+=r_{1,\max}=2/3$),
      $Q_n=1,\theta_n=\varphi-1$ gives $1-\sqrt5/3=c_{\mathrm{sec}}$ at $D_{j-1}=3$
      ($r_2=1/3$) and $1-1/\sqrt5\approx0.552786$ at $D_{j-1}=1$ ($r_2=0$);
      all other branches in~(a)--(c) yield $\ge(2/\sqrt5)(2-\varphi)\approx0.341641$
      (including $D_j=1$, $Q_n\ge2$, $\theta_n\notin\{\varphi-1,2-\varphi\}$,
      which satisfies $E_n\ge(\sqrt5-1)-o(1)\approx1.236068$).
\item In~(ii) ($r_{1,\min}=1/2$, $r_{1,\max}=5/6$, $r_{2,\min}=1/6$),
      $Q_n=1,\theta_n=\varphi-1$ at $D_{j-1}=2$ produces the excluded extremal
      family $[a_0;\ldots,b_{j-1},\overline1]$, and at $D_{j-1}=6$ gives
      $1-4/(3\sqrt5)\approx0.403715$; all other branches yield $\ge0.490712>c_{\mathrm{sec}}$.
\item In~(iii) ($r_{1,\min}=1/4$, $r_{1,\max}=11/12$, $r_{2,\min}=1/12$),
      $Q_n=1,\theta_n=\varphi-1$ at $r_2=7/12$ produces the excluded extremal
      family $[a_0;\ldots,b_{j-1},2,2,\overline1]$, at $r_2=3/4$ gives
      $c_{\mathrm{sec}}=\sqrt5/2-1$, and at $r_2\in\{1/4,1/12\}$ gives $\ge0.329180$;
      all other branches yield $\ge\min(0.416176,\,\sqrt5/6\approx0.372678)>c_{\mathrm{sec}}$.
\item In~(iv) ($r_{1,\min}=r_{2,\min}=1/98$, $r_{1,\max}=97/98$),
      $Q_n=1,\theta_n=\varphi-1$ gives $\ge16/(7\sqrt5)-1\approx0.0222025$
      (at $D_{j-1}=14, q_1'=9$);~(a) and~(b) for $Q_n\ge2$ give $\ge0.350768$;
      and~(c) gives $\ge1.75384$ ($Q_n=1$) and $\ge2\sqrt5/98=\sqrt5/49\approx0.0456340>c_{\mathrm{sec}}$
      ($Q_n\ge2$).
\end{itemize}
Thus every non-extremal sequence has $\liminf E_n\ge c_{\mathrm{sec}}$ in all four cases.
\end{proof}

\begin{remark}[Spectral regimes and open parametric frontier for $B\ge7$]\label{rem:golden-general-regimes}
With Theorems~\ref{thm:golden-exact-second}, \ref{thm:golden-b2-exact},
\ref{thm:golden-b3-exact}, \ref{thm:golden-b4-exact}, \ref{thm:golden-b5-exact},
\ref{thm:golden-b6-exact}, and~\ref{thm:golden-multiscale-exact}, together with
Corollary~\ref{cor:affine-single} and Lemma~\ref{lem:general-case2-criterion},
complete Case~1 and Case~2 structural classifications (exact $c_\varphi$, all
extremal families, and exact second limit points) are now proved for every
coprime $A/B-e$ with $1\le B\le 6$, for pure unbounded degeneration ($|C|=1$,
$B\in\{2,4\}$, $A$ odd, where Lin~\cite[Lemma~4.2]{Lin2018} gives
$c_\varphi(x)=3/\sqrt5-1$), and for the small-period mixed-$D_j$ affine targets
$(1+e)/6$, $(1+e)/3$, $(1+2e)/3$, and $22/7-e$ (including single-digit
$r_j=1$ phases). For general parametric denominators $B\ge7$ and arbitrary
$(A,B,C)$, Theorem~\ref{thm:golden-finite-shift} guarantees finiteness and
effective computability up to a fixed bound for each fixed $(A+Ce)/B$ with $C\ne0$, whereas closed-form parametric
formulas for $c_\varphi(A/B-e)$ and its second limit point across arbitrary
$B\ge7$ remain open.
\end{remark}

\begin{remark}[Formalization and numerical verification scope for the golden-ratio orbit models]\label{rem:golden-b2-model}
High-precision verification (400 digits in
\texttt{verification/\allowbreak{}verify\_\allowbreak{}golden\_\allowbreak{}orbit\_\allowbreak{}model.py}
and \texttt{verification/\allowbreak{}verify\_\allowbreak{}golden\_\allowbreak{}b2\_\allowbreak{}exact.py},
and 2500 digits in
\texttt{verification/\allowbreak{}verify\_\allowbreak{}golden\_\allowbreak{}b3\_\allowbreak{}b4\_\allowbreak{}exact.py}
and \texttt{verification/\allowbreak{}verify\_\allowbreak{}golden\_\allowbreak{}b5\_\allowbreak{}b6\_\allowbreak{}multiscale.py})
confirms the norm parity and rational-endpoint rigidity identities across
all $744$ matrices in $\mathrm{GL}_2(\Z)\cap[-6,6]^4$, the conjugate lower
bound of Proposition~\ref{prop:conjugate-positivity} across $2920$ matrices in
$[-12,12]^4$, the single-digit Case~2 bound $h(\theta)(1-\theta)\ge\sqrt5-1$
across $688$ elements of $E_\varphi\cap(0,1)\setminus\{\varphi-1,2-\varphi\}$
in $[-12,12]^4$, the $\Z[\varphi]$ norm mod-$5$ ramification identity across
$10201$ pairs, the exact base periods, suffix sets, algebraic $4\tau$
inversions, modular exclusions (mod-$16$ at $B=4$ and mod-$8$ at $B=6$),
Case~2 branch bounds, extremal convergence, and non-extremal scans for
$B=2,\dots,6$ and the multi-scale targets $(1+e)/6$, $(1+e)/3$, $(1+2e)/3$,
and $22/7-e$. In Lean~4,
\texttt{EpiFormal/\allowbreak{}GoldenOrbit\allowbreak{}Rigidity.lean}
($15$ declarations audited in \texttt{Audit.lean}) and
\texttt{EpiFormal/\allowbreak{}GoldenOrbit\allowbreak{}B2Rigidity.lean}
($22$ declarations audited in \texttt{Audit.lean})
machine-verify the M\"obius error identity \eqref{eq:bugeaud-34},
Lemmas~\ref{lem:norm-parity-rigidity} and~\ref{lem:pell-b2-rigidity}, and
the core algebraic, uniform-separation, and concave-quadratic lower bounds of
Theorems~\ref{thm:golden-exact-second} and~\ref{thm:golden-b2-exact};
\texttt{EpiFormal/\allowbreak{}GeneralEndpoint\allowbreak{}Rigidity.lean}
($30$ declarations audited in \texttt{Audit.lean}, base axioms only)
machine-verifies Proposition~\ref{prop:conjugate-positivity}, the full
four-part arithmetic and algebraic-inversion rigidity of
Lemma~\ref{lem:general-endpoint-rigidity} (including general odd inert-prime
reciprocity and the mod-$2D$ congruence sieve), and the arithmetic sieve and
spectral-gap core inequalities for $B\in\{3,4,5,6\}$;
\texttt{GoldenHeadCutoff.lean} checks the three growing-head inequalities;
and the continued-fraction complete-quotient convergence, Galois/Serret purely
periodic uniqueness,
Lemmas~\ref{lem:even-suffix-sep} and~\ref{lem:general-case2-criterion},
and the final Case~1/Case~2 $\liminf$ assembly of
Theorems~\ref{thm:golden-b3-exact}--\ref{thm:golden-multiscale-exact} remain
paper proofs backed by 2500-digit numerical verification.
These results sharply characterize the $e$-side control models for
$B\in\{1,\dots,6\}$ and the four multi-scale targets, while leaving the
$\pi$-side condition $c_\varphi(\pi)=0$ below open.
\end{remark}

By \cite[Lemma~4.1]{BugeaudQuad}, $c_\varphi(\pi)=0$ is equivalent to
arbitrarily long runs of ones in the ordinary continued fraction of $\pi$.
For the forward construction from a run of $\ell$ ones starting after
digit $s$, choose a maximal run, so that $a_s\ne1$, and replace the
tail after $s$ by $\overline{1}$. The two numbers share a cylinder
of length at most $q_{s+\ell}^{-2}$, and
$q_{s+\ell}\ge F_{\ell+1}q_s$.
By \cite[Lemma~2.1]{BugeaudQuad}, the resulting height lies between
$q_s^2/32$ and $4q_s^2$ for $s\ge3$. Thus arbitrarily long runs with
$s\to\infty$ give heights tending to infinity and normalized errors
at most $4/F_{\ell+1}^2$.
Either statement excludes $\pi\in\Q+\Q e$ by the first bound in
Corollary~\ref{cor:golden-bound} and the irrationality of $\pi$.
Thus it would give rational linear independence of $1,e,\pi$ and,
in particular, the irrationality of $e+\pi$.
One possible arithmetic target is to produce orbit points
$\beta=(-v+\sigma\sqrt5)/(2u)$, where $u,v,w$ are integers,
$v^2-4uw=5$ and $\sigma\in\{-1,1\}$, for which $|u|\to\infty$ and
$|2u\pi+v-\sigma\sqrt5|\to0$. Indeed, the normalized error is exactly
$|2u\pi+v-\sigma\sqrt5|/\sqrt5$. The required sequence remains open.
Theorem~4.2 of \cite{BugeaudQuad} concerns Lebesgue measure and
Hausdorff dimension. Theorem~3 and Section~5 of \cite{Hildebrand2024}
give almost-everywhere digit frequencies and a finite experiment at $\pi$.

For a finite example, the first three consecutive ones at $\pi$ occur
at indices $5,6,7$ (with integer part $a_0$).
Replacing the tail after $[3;7,15,1,292]$ by $\overline{1}$ gives
$\beta=(103993\varphi+355)/(33102\varphi+113)$. The matrix has
determinant $-1$. Substitution in $\varphi^2-\varphi-1=0$ gives
$(u,v,w)=(-1091989109,6861169923,-10777500509)$ with $\sigma=-1$.
Exact intervals give
\[
  0.014695421944 < h(\beta)|\pi-\beta| < 0.014695421945 .
\]
This estimates the error at one orbit point.

\subsection{Ideal classes in the trace test}

\begin{theorem}[Ideal classes of coprime trace norms]\label{thm:trace-class-sieve}
Let $\tau$ have positive nonsquare primitive discriminant $\Delta$,
and set
$\mathcal O_\Delta=\Z[(\Delta+\sqrt\Delta)/2]$. If $A_\tau>0$ is the
leading coefficient of its primitive polynomial, let
$\mathfrak c_\tau=[(A_\tau,A_\tau\tau)]$ in the ordinary group of
invertible fractional ideal classes of $\mathcal O_\Delta$.
For an integer $m$ such that $N_m=(m^2-\Delta)/4\in\Z$ and
$\gcd(m,N_m)=1$, put $\omega=(-m+\sqrt\Delta)/2$. For each prime
$p\mid N_m$ set $\mathfrak p_p=(p,\omega)$ and $k_p=v_p(|N_m|)$.
For any $x$ with partial quotients tending to infinity,
\begin{equation}\label{eq:trace-class-sieve}
 |1-m/\sqrt\Delta|\in\mathcal L_\tau(x)
 \quad\Longleftrightarrow\quad
 \mathfrak c_\tau^{\pm1}\in
 \left\{\,\prod_{p\mid N_m}[\mathfrak p_p]^{i_p}:
                    0\le i_p\le k_p\,\right\},
\end{equation}
where membership on the right means that at least one of the two
classes occurs. In particular, when $|N_m|=p^k$ the only possible
orbit classes, up to inversion, are the consecutive powers
$1,[\mathfrak p_p],\ldots,[\mathfrak p_p]^k$;
$[\mathfrak p_p]^k=1$.
\end{theorem}
\begin{proof}
For a primitive polynomial $AX^2+bX+C$ with discriminant
$\Delta$, write $\eta=A\gamma$ for one of its roots multiplied by
$A$. Then $\mathcal O_\Delta=\Z[\eta]$, and the lattice
$I=\Z A+\Z\eta$ is an ideal by $\eta^2+b\eta+AC=0$.
The product $I\bar I/A$ is contained in $\mathcal O_\Delta$ and
contains $A,\eta,\bar\eta,C$. Since
$\eta+\bar\eta=-b$ and $\gcd(A,b,C)=1$, the product contains $1$.
Thus $I\bar I=A\mathcal O_\Delta$ and $I$ is invertible.
For two quadratic roots, homothety of the lattices
$\Z+\Z\gamma$ and $\Z+\Z\tau$ is equivalent to a unimodular
change of basis, hence to a $\operatorname{GL}_2(\Z)$ fractional
linear map between the roots. Conjugating inverts the ideal class
by the product identity. This is the ordinary, unoriented form of
the ideal correspondence.

For $a\mid |N_m|$, $a>0$, the polynomial
$aX^2+mX+N_m/a$ is automatically primitive, since any common prime
would divide $m$ and $N_m$. Its positive-square-root choice
has ideal $I_a=(a,\omega)$; the other choice has the inverse class.
For $1\le i\le k_p$,
\[
 (p,\omega)^i=(p^i,\omega).
\]
Indeed, in the induction step the product is generated by
$p^i,p^{i-1}\omega,p\omega,\omega^2$.
Since $p^i\mid N_m$ and $\omega^2=-m\omega-N_m$, its generators
give $m\omega$ and $p\omega$. B\'ezout, using $p\nmid m$, gives
$\omega$. Containment in the other direction follows from the same
identity. For coprime positive $a,b$ dividing $|N_m|$ whose product
also divides it,
$(a,\omega)(b,\omega)=(ab,\omega)$, using $a\omega,b\omega$
and $ab\mid N_m$. Hence $I_a$ is the product of the stated
prime-ideal powers, and $I_{|N_m|}=(\omega)$ by
$\omega\bar\omega=N_m$. Each $\mathfrak p_p$ is prime because
evaluation at $\omega=0$ gives quotient $\Z/p\Z$.
These identities, the homothety test and
Theorem~\ref{thm:pure-trace-test} prove the result.
\end{proof}

\paragraph{A nonprincipal contribution at a squared prime norm.}
For $\Delta=60$ and $m=16$, $N_m=49$. Here $\omega=\sqrt{15}-8$
and the prime above $7$ has order two: its square is principal by
Theorem~\ref{thm:trace-class-sieve}. If the prime were principal,
a generator in $\Z[\sqrt{15}]$ would have absolute norm $7$.
The equations $u^2-15v^2=7$ and $u^2-15v^2=-7$ are excluded modulo
$5$ and $3$, respectively. Its root $\omega/7$ belongs to the
orbit of $\tau=\sqrt{15}/3$, by
$\omega/7=(-2\tau+3)/(\tau-2)$ with determinant $1$.
Thus this trace contributes to both the monic orbit and this
nonprincipal orbit, with the divisors $1,49$ giving the principal
class and $7$ giving the nonprincipal class. This contrasts with
the prime case of Corollary~\ref{cor:prime-norm-orbit}.

The two norm congruences for this example are audited in
\texttt{ResearchAdvance.lean}. The ideal-class correspondence and
product formula are paper proofs.



\begin{theorem}[Trace classes for arbitrary nonzero norms]
\label{thm:trace-all-norms}
Keep $\Delta,\mathcal O_\Delta$ and $\mathfrak c_\tau$ as in
Theorem~\ref{thm:trace-class-sieve}. Let $N_m=(m^2-\Delta)/4$ be an
integer and $\omega=(-m+\sqrt\Delta)/2$. For each prime power
$p^{k_p}\parallel |N_m|$, define a set of ideal classes by
\[
 \mathcal F_p(m)=
 \begin{cases}
 \{[(p,\omega)]^i:0\le i\le k_p\},&p\nmid m,\\
 \{1,[(p^{k_p},\omega)]\},&p\mid m .
 \end{cases}
\]
All displayed classes are invertible classes, even when $p$
divides the conductor. For any $x$ with partial quotients tending
to infinity, $|1-m/\sqrt\Delta|\in\mathcal L_\tau(x)$ if and only
if $\mathfrak c_\tau$ or $\mathfrak c_\tau^{-1}$
belongs to $\prod_{p\mid N_m}\mathcal F_p(m)$.
\end{theorem}
\begin{proof}
For $a>0$ dividing $|N_m|$, the form with coefficients
$(a,m,N_m/a)$ is primitive precisely when, for every $p\mid m$,
the exponent of $p$ in $a$ is either zero or $k_p$.
Indeed a common prime occurs precisely when $p\mid m$ and both
outer coefficients contain $p$.
For coprime $a,b$ with $ab\mid N_m$,
$(a,\omega)(b,\omega)=(ab,\omega)$:
the generators $a\omega,b\omega$ generate $\omega$ over $\Z$,
and $\omega^2=-m\omega-N_m$ belongs to $(ab,\omega)$.
In the case $p\nmid m$, the prime-power product in the proof of
Theorem~\ref{thm:trace-class-sieve} applies. In the case $p\mid m$,
the whole factor $(p^{k_p},\omega)$ is invertible, since
$(p^{k_p},m,N_m/p^{k_p})$ is primitive. This uses the
product-with-conjugate identity of that theorem, which holds
in the quadratic order itself. Factoring $(a,\omega)$ over the
coprime prime-power divisors of $a$ now gives exactly the displayed
class sets. Apply the pure-head trace test and the lattice
correspondence from Theorem~\ref{thm:trace-class-sieve}.
\end{proof}

\begin{corollary}[Common prime-power traces]
\label{cor:common-prime-power}
If $|N_m|=p^k$ and $p\mid m$, then $|1-m/\sqrt\Delta|$ occurs
exactly in the monic orbit. For example, at $\Delta=20$ and $m=6$,
the conductor prime $2$ has exponent two. The admissible divisors
are $1,4$; $a=2$ would give the nonprimitive form $(2,6,2)$.
\end{corollary}
\begin{proof}
The only nontrivial factor in Theorem~\ref{thm:trace-all-norms}
is $(p^k,\omega)=(\omega)$, since $|N_m|=p^k$ and
$\omega^2+m\omega=-N_m$. The principal ideal class corresponds
to the monic orbit. In the example the order is
$\Z[\sqrt5]$, of index two in the maximal order of
$\Q(\sqrt5)$, and $N_m=4$.
\end{proof}

\begin{theorem}[Intersections of quadratic orbit spectra]
\label{thm:trace-common-progression}
Let $x$ have partial quotients tending to infinity.
\begin{enumerate}
\item For every fixed $\Delta$, the intersection of the spectra over
all orbits of primitive discriminant $\Delta$ contains all
$|1-m/\sqrt\Delta|$ from a nonempty arithmetic progression
of integers $m$. The progression can be chosen with
$\gcd(m,N_m)=1$ throughout.
\item More generally, for every finite family of quadratic orbits
in the same real quadratic field $\Q(\sqrt d)$, $d$ squarefree,
their spectra share all values $|1-r/\sqrt d|$ from a nonempty
arithmetic progression of integers $r$.
\item If two orbits lie in distinct real quadratic fields, their
spectra have at most the common value $1$.
\end{enumerate}
In each progression assertion the shared values have a positive
linear lower counting rate as the upper bound of the values grows.
\end{theorem}
\begin{proof}
We first record a choice of representatives. For a primitive
binary quadratic form $Q(U,V)=A U^2+bUV+CV^2$ with $A>0$
and a finite set of primes, select for each prime a vector at which
$Q$ is nonzero modulo that prime. One of $(1,0),(0,1),(1,1)$ works,
by primitivity. Lift these vectors by the Chinese remainder theorem.
Adding a sufficiently large multiple of the product of the primes
to the first coordinate makes the value positive. Divide the two
coordinates by their greatest common divisor: the divisor is a
unit modulo every selected prime, so the new value is still
coprime to them. Complete the primitive vector to a unimodular
basis and change variables. The new positive leading coefficient
is coprime to the selected primes. In a finite list we may thus
take the leading coefficients $A_i$ pairwise coprime and each
coprime to any fixed finite list of primes. For the resulting
form $(A_i,b_i,C_i)$, substituting $X\mapsto X+t$ into its
polynomial yields every middle coefficient
$m_i\equiv b_i\bmod 2A_i$ within the same orbit.

For (i), there are only finitely many relevant orbits.
A reduced complete quotient $\gamma$ from each orbit satisfies
$\gamma>1$ and $-1<\gamma'<0$. Its leading coefficient is
$A=\sqrt\Delta/(\gamma-\gamma')<\sqrt\Delta$ and
$|b|=A|\gamma+\gamma'|<\sqrt\Delta$, leaving only finitely
many primitive triples. Choose all $A_i$ coprime to $2\Delta$
as well as pairwise coprime. The $b_i$ have common parity
$\Delta\bmod2$. Writing $m$ with this parity, the requirements
$m\equiv b_i\bmod 2A_i$ become simultaneous congruences
modulo $A_i$ for $(m-(\Delta\bmod2))/2$.
Require also $m\equiv1\bmod p$ at odd primes $p\mid\Delta$.
When $\Delta$ is even, choose the parity of $m/2$ opposite
to that of $\Delta/4$, which makes $N_m$ odd.
All these conditions are compatible by the Chinese remainder
theorem. A common odd prime of $m,N_m$ would divide $\Delta$,
so the resulting progression is coprime as claimed.

For (ii), write $\Delta_i=d c_i^2$ with $c_i>0$ integers.
Choose $A_i$ pairwise coprime and coprime to every $2c_j$,
and choose $r\equiv d\bmod2$.
Then $c_i r$ and $b_i$ have the same parity. The congruences
$c_i r\equiv b_i\bmod 2A_i$ impose one residue on
$(r-(d\bmod2))/2$ modulo each $A_i$. The Chinese remainder theorem
again gives a progression. Since $m_i=c_i r$, all its traces
give the same values $|1-r/\sqrt d|$.
In both parts the asserted values follow from
Theorem~\ref{thm:pure-trace-test}. The trace-to-value map is
injective on the integers, since equality for distinct traces
would make their sum twice the irrational square root.
For a progression of common difference $M$, counting its
integers in the interval from the absolute-value bound $T$
gives $2\sqrt\Delta\,T/M+O(1)$ for (i), or
$2\sqrt d\,T/M+O(1)$ for (ii).

For (iii), equality of values would give
$m_1/\sqrt{\Delta_1}-m_2/\sqrt{\Delta_2}=0$
or $m_1/\sqrt{\Delta_1}+m_2/\sqrt{\Delta_2}=2$.
The intersection of distinct quadratic fields is $\Q$,
so either relation forces both individual terms to be rational,
hence zero. Only the first relation can then hold.
\end{proof}

For the two orbits of $\sqrt{15}$ and $\sqrt{15}/3$,
the corresponding discriminant-$60$ forms can be taken as
$(1,0,-15)$ and $(7,2,-2)$. The latter comes from
$3U^2-5V^2$ by $(U,V)=(2X+Y,X+Y)$.
Thus $m\equiv16\bmod 420$ gives a common progression
with $\gcd(m,N_m)=1$.
The primitive-factor restriction in
Theorem~\ref{thm:trace-all-norms} and the congruence behind
this progression are audited in \texttt{TraceMarkov.lean}.
The class and orbit statements use paper proofs. The
certificate \texttt{verify\_trace\_markov.py} checks the primitive
divisors and ideal products on a bounded range and independently
verifies two explicit progressions by exact continued fractions.


\begin{theorem}[Density one in the trace lattice]
\label{thm:trace-density-one}
Fix a real quadratic orbit of primitive discriminant $\Delta$,
and let $x$ have partial quotients tending to infinity.
Uniformly over real intervals $I$ of length $X$, for
$X$ sufficiently large,
\[
 \#\{\,m\in I\cap\Z: m\equiv\Delta\bmod 2,\ 
           |1-m/\sqrt\Delta|\notin\mathcal L_\tau(x)\,\}
       \ll_\tau \frac{X}{(\log X)^{\delta_\tau}}.
\]
Here $\delta_\tau=1/(4h_f)$ is permissible, where $h_f$ is
the order of the ray class group of $K=\Q(\sqrt\Delta)$ for
the finite modulus $f\mathcal O_K$ and $\Delta=f^2D_K$.
The common spectrum of all orbits at fixed $\Delta$ therefore
has the same leading counting coefficient as its parity lattice.
\end{theorem}
\begin{proof}
Write $K=\Q(\sqrt\Delta)$, $\Delta=f^2D_K$ with $D_K$ the
field discriminant, and $\mathcal O_\Delta=\Z+f\mathcal O_K$.
Put $\mathfrak c=\mathfrak c_\tau$ in the ordinary invertible ideal
class group. Choose a representative $(A,A\tau_0)$ with
$\gcd(A,2\Delta)=1$ using the leading-coefficient choice in
Theorem~\ref{thm:trace-common-progression}. It is prime to the conductor.
Extend it to $\mathcal O_K$ and consider its ray class of modulus
$f\mathcal O_K$, with no infinite part.
Extension and contraction give the correspondence of ideals
prime to the conductor \cite[Theorem~3.8]{ConradConductor}.
For completeness, any principal ray relation $(\gamma)$ respects
the ordinary order class: at primes away from $f$ the two orders
agree locally, and at a rational prime $\ell\mid f$,
$\gamma\in1+f(\mathcal O_K\otimes\Z_\ell)$ belongs to
$(\mathcal O_\Delta\otimes\Z_\ell)^\times$.
Thus every ideal of this ray class contracts, in the fractional
ideal correspondence prime to the conductor, to the same order class.

The positive Dirichlet density of prime ideals in a ray class
\cite[Chapter~VII, Section~13]{NeukirchANT} gives a sum of
inverse norms to the power $s$ asymptotic to
$h_f^{-1}\log(1/(s-1))$ as $s\downarrow1$.
Prime ideals of residue degree two have a bounded contribution,
by comparison with $\sum_p p^{-2s}$.
Discard the primes over $2\Delta$ and choose one of the eligible
ideals over each remaining rational prime. Since at most two
eligible ideals lie over any one $p$, the resulting set
$\mathcal P$ satisfies
$\sum_{p\in\mathcal P}p^{-s}\ge
 (1/(2h_f)+o(1))\log(1/(s-1))$.
In particular,
\begin{equation}\label{eq:trace-primes-reciprocal}
 \sum_{\substack{p\in\mathcal P\\p\le Y}}p^{-1}
           \ge\left(\frac{1}{2h_f}+o(1)\right)\log\log Y.
\end{equation}
For this last deduction take $s=1+1/\log Y$. The elementary
prime-counting upper bound $\pi(u)\ll u/\log u$ and partial
summation bound the tail over all $p>Y$ by
$O(\int_Y^\infty u^{-s}/\log u\,du)=O(1)$, using
$v=\log u/\log Y$.
For each $p\in\mathcal P$ let $\mathfrak P_p$ be the selected
degree-one prime of $\mathcal O_K$ and let
$r_p$ be the reduction of $\sqrt\Delta$ modulo $\mathfrak P_p$.
Its contraction has norm $p$ and class $\mathfrak c$.
If $m\equiv r_p\bmod p$, then
$p\mid N_m=(m^2-\Delta)/4$ and $p\nmid m$, since $p\nmid2\Delta$.
Furthermore $(p,(-m+\sqrt\Delta)/2)$ is that contraction:
the generators lie in the contraction and its quotient is
$\mathbb F_p$, using $\mathcal O_\Delta=\Z[(-m+\sqrt\Delta)/2]$.
Taking the divisor $p$ in Theorem~\ref{thm:trace-all-norms}
therefore certifies this trace.

It remains to count integers missing all the selected residues.
Take $k=\lceil B\log\log X\rceil$ for a large fixed absolute
constant $B$, put $Y=X^{1/(4k)}$, and choose $m$ uniformly
from the indicated parity class in $I$.
Write $\mu=\sum_{p\in\mathcal P,\,p\le Y}1/p$.
For a selected subset with product $q$, the joint condition
is one residue class modulo $q$ by the Chinese remainder
theorem. All primes are odd, so after prescribing the parity
it is one residue class modulo $2q$. Its count in $I$
is $X/(2q)+O(1)$, with an absolute error uniform in the
subset and the location of $I$. The full parity count is
$X/2+O(1)$. Dividing gives a joint hit probability
$1/q+O(1/X)$ with the same uniformity. Apply the upper Bonferroni bound through order $2k$.
The sum of these progression-count errors is
$O((2k+1)Y^{2k}/X)=O((2k+1)X^{-1/2})$.
For independent events of probabilities $1/p$, the same
truncated sum is at most the no-hit probability plus the
$(2k+1)$st elementary symmetric sum, by the lower Bonferroni
bound at order $2k+1$. Thus the desired no-hit probability is at most
\[
 e^{-\mu}+\frac{\mu^{2k+1}}{(2k+1)!}
                  +O\bigl((2k+1)X^{-1/2}\bigr).
\]
The elementary prime-counting upper bound also gives
$\mu\ll\log\log X$. With $B$ large enough, the factorial term
is $O((\log X)^{-1})$, by $n!\ge(n/e)^n$.
Equation~\eqref{eq:trace-primes-reciprocal} and
$\log\log Y\sim\log\log X$ give
$e^{-\mu}\ll_\tau(\log X)^{-1/(4h_f)}$.
Every absent trace has no hit, proving the estimate.
There are finitely many orbits at fixed $\Delta$, as shown
in Theorem~\ref{thm:trace-common-progression}; take their
finite union of exceptions.
\end{proof}

\begin{corollary}[Leading count for the common spectrum]
\label{cor:trace-common-density}
Let a finite nonempty family of orbits lie in the same
field $\Q(\sqrt d)$, $d$ positive and squarefree, with primitive
discriminants $\Delta_i=d c_i^2$. Put $g=\gcd_i(c_i)$.
Uniformly in $T\ge0$, the number of distinct values in the
common spectrum lying in $[T,T+H]$ is
\[
 g\sqrt d\,H+O\!\left(\frac{H}{(\log H)^{\delta}}\right),
 \qquad \delta=\min_i\delta_{\tau_i}>0.
\]
The implied constant may depend on the family.
Moreover every interval of length $H$ in the integer
parameter $u$ of the common parity lattice contains a run
of $\gg (\log H)^\delta$ consecutive admissible lattice points
all of whose values are shared.
\end{corollary}
\begin{proof}
Equality of two values forces equality of the rational ratios
$m_i/c_i$: the other sign would require their sum to be
$2\sqrt d$. The ratios common to all denominator lattices
are exactly $u/g$ with $u\in\Z$. Write $c_i=g k_i$.
At least one $k_i$ is odd, and the conditions
$k_i u\equiv d c_i^2\bmod2$ together say exactly
$u\equiv dg\bmod2$. The candidate common values are therefore
$|1-u/(g\sqrt d)|$ for this single parity. The map from $u$
to the value is injective. The condition of lying in
$[T,T+H]$ cuts out two intervals in $u$, each of length
$g\sqrt d\,H$ and with disjoint interiors, so the total
candidate count is $g\sqrt d\,H+O(1)$.
For each $i$, absent values among these candidates correspond
injectively to exceptions of Theorem~\ref{thm:trace-density-one}
in the two scaled $m_i$-intervals, each of length
$k_i g\sqrt d\,H$. The location-uniform estimate proves the
asserted count.

The same argument applied directly in any $u$-interval
of length $H$ leaves at most $O(H/(\log H)^\delta)$
missing parity-lattice points. There are $H/2+O(1)$
candidate points. If $E$ are missing, the others split into
at most $E+1$ runs, so at least one run has length
$\gg (\log H)^\delta$.
\end{proof}
The ray class result supplies the prime residues; the trace
density follows from the finite inclusion--exclusion estimate.
For instance at $\Delta=60$, $m=4$ gives $N_m=-11$, so it
occurs in the monic orbit and is absent from the nonprincipal
orbit of $\sqrt{15}/3$, by Corollary~\ref{cor:prime-norm-orbit}
and the nonprincipal example above. The residue and parity
steps are audited in \texttt{DensityCube.lean}; the class-field
and density arguments are paper proofs.


\section{Reductions, obstructions, and conditional frameworks}\label{sec:new}

This chapter assembles the material directed at
Problem~\ref{prob:main} itself: the elementary reductions and
conditional implications, the systematic obstruction analysis, the
conditional frameworks, and the exclusion map with the surviving
targets. The material is diagnostic: the verdicts concern specific
constructions; they are not meta-theorems that no proof along the
general line can exist.

\subsection{The irrationality-measure reduction}

\begin{lemma}[Invariance of the irrationality measure]\label{lem:invariance}
Let $x\in\R$ be irrational, $r\in\Q$, and $s\in\Q\setminus\{0\}$. Then
\[
\irr(x+r)=\irr(x),\qquad \irr(sx)=\irr(x),\qquad \irr(1/x)=\irr(x).
\]
\end{lemma}

\begin{proof}
We give a denominator-preserving comparison up to fixed factors. If infinitely many pairs approximate an irrational real at a fixed positive exponent, their denominators are unbounded: for each fixed denominator the numerators lie in a bounded interval. Fix $\nu>0$ realized infinitely often and $0<\varepsilon<\nu$.

\emph{Translation.} Write $r=a/b$ with $a\in\Z$, $b\in\N_{>0}$. Each approximation $|x-p/q|<q^{-\nu}$ gives
\[
\left|x+r-\frac{bp+aq}{bq}\right|<q^{-\nu}
=b^\nu(bq)^{-\nu}<(bq)^{-(\nu-\varepsilon)}
\]
for all sufficiently large $q$. The new denominators $bq$ are unbounded, so $\irr(x+r)\ge\nu-\varepsilon$. Taking suprema and then letting $\varepsilon\downarrow0$ gives $\irr(x+r)\ge\irr(x)$, also when the supremum is infinite. Apply the same argument with $-r$ for the reverse inequality. No assertion that reduced denominators divide $q$ is required.

\emph{Scaling.} Write $s=u/v$ with $u\in\Z\setminus\{0\}$ and $v\in\N_{>0}$. Then
\[
\left|sx-\frac{up}{vq}\right|<|s|q^{-\nu}
=|s|v^\nu(vq)^{-\nu}<(vq)^{-(\nu-\varepsilon)}
\]
eventually. This proves $\irr(sx)\ge\irr(x)$; scaling by $1/s$ gives the opposite inequality.

\emph{Reciprocal.} Since $x\ne0$, along the approximating pairs $p/q\to x$ and hence $|p|\asymp q$, with $p\ne0$ eventually. Put $q'=|p|$ and $p'=\operatorname{sgn}(p)q$. Then $p'/q'=q/p$ and
\[
\left|\frac1x-\frac{p'}{q'}\right|
=\frac{q}{|x|\,|p|}\left|x-\frac pq\right|
<C(q')^{-\nu}<(q')^{-(\nu-\varepsilon)}
\]
eventually, for a constant $C>0$ depending only on $x$ and $\nu$. Thus $\irr(1/x)\ge\irr(x)$; applying this inequality to $1/x$ proves equality.
\end{proof}

\begin{lemma}[Classical; Bundschuh {\cite{Bundschuh1971}}, see also {\cite{CNU2024}}]\label{lem:mue}
$\irr(\ee)=2$. More generally, $\irr(\ee^r)=2$ for every $r\in\Q\setminus\{0\}$.
\end{lemma}

\begin{theorem}[Measure reduction]\label{thm:measure-reduction}
If any one of $e+\pi$, $e-\pi$, $e\pi$, $e/\pi$ is rational, then $\irr(\pi)=2$.
\end{theorem}

\begin{proof}
We treat the four cases using Lemmas~\ref{lem:invariance} and~\ref{lem:mue}.
\begin{enumerate}
\item If $e+\pi=r\in\Q$, then $\pi=r-e$, so $\irr(\pi)=\irr(-e)=\irr(e)=2$ (translation and scaling by $-1$).
\item If $e-\pi=r\in\Q$, then $\pi=e-r$, so $\irr(\pi)=\irr(e)=2$.
\item If $e\pi=r\in\Q$, then $r\ne0$ and $\pi=r/e$, so $\irr(\pi)=\irr(1/e)=\irr(e)=2$ (scaling and reciprocal).
\item If $e/\pi=r\in\Q$, then $r\ne0$ and $\pi=e/r$, so $\irr(\pi)=\irr(e)=2$.
\end{enumerate}
In every case $\irr(\pi)=2$.
\end{proof}

\begin{corollary}\label{cor:four-irrational}
If $\irr(\pi)>2$, then all four numbers $e+\pi$, $e-\pi$, $e\pi$, $e/\pi$ are irrational.
\end{corollary}

\begin{proof}
Contrapositive of Theorem~\ref{thm:measure-reduction}.
\end{proof}

\begin{remark}[Status and honest assessment]
The published upper bounds for $\irr(\pi)$ are: $\le42$ (Mahler 1953), $\le20.6$ (Mignotte 1974), $\le19.89$ (Chudnovsky 1982), $\le8.016045$ (Hata 1993), $\le7.606308$ (Salikhov 2008), $\le7.10320534$ (Zeilberger--Zudilin 2020) \cite{Alekseyev2011,ZZ2020}; a 2026 preprint of Bai claims $<7.101862832357$ \cite{Bai2026}. Whether $\irr(\pi)=2$ is unknown; a preprint of Carella claims it \cite{CarellaMeasure} but is not accepted and is not used here. Since almost all real numbers have measure $2$ and $\irr(\pi)=2$ is the expected value, the sufficient target $\irr(\pi)>2$ is probably false. The value of Theorem~\ref{thm:measure-reduction} is structural: it connects the problem to an active research line, and any proof via Diophantine approximation must confront the measure. Note also the boundary of the method: Lemma~\ref{lem:invariance} concerns rational translations only; invariance under algebraic translation is unknown, so the reduction cannot even be started under the weaker hypothesis ``$e+\pi$ algebraic.''
\end{remark}

\subsection{A continued-fraction necessary condition}

We use two classical results: Serret's theorem (two reals have eventually equal continued-fraction tails iff they lie in the same $\mathrm{PGL}(2,\Z)$-orbit) \cite{Panti2017}, and Raney's theorem (for an integer-coefficient linear fractional transformation $Mx=(ax+b)/(cx+d)$ with $ad-bc\ne0$, the partial quotients of $Mx$ are computed from those of $x$ by a finite-state transducer) \cite{Raney1973}.

\begin{theorem}\label{thm:cf}
If $e+\pi=p/q$ with $p,q\in\N$, then the sequence of partial quotients of $\pi$ is the image, under a finite-state transducer determined by $p,q$, of the explicitly known partial quotients of $\ee$.
\end{theorem}

\begin{proof}
From $e+\pi=p/q$ we obtain $\pi=(p-qe)/q=(-q\cdot e+p)/q$, a linear fractional transformation of $e$ with integer coefficients and determinant $-q^2\ne0$. Raney's theorem applies directly. The partial quotients of $\ee$ are explicitly known: $\ee=[2;1,2,1,1,4,1,1,6,1,1,8,\dots]$, with the pattern $1,2k,1$ repeating.
\end{proof}

\begin{proposition}\label{prop:not-integer}
$e+\pi\notin\Z$.
\end{proposition}

\begin{proof}
From the rigorous decimal enclosures $2.71828<e<2.71829$ and $3.14159<\pi<3.14160$ we obtain $5.85987<e+\pi<5.85989$, an interval containing no integer.
\end{proof}

\begin{remark}[The obstruction]
Theorem~\ref{thm:cf} is a strong structural necessary condition: a finite-state transduction of the explicit sequence of $\ee$ is highly structured. But no structural theorem about the partial quotients of $\pi$ is known (not even their unboundedness is proved), so no contradiction can be derived. Numerically the continued fraction of $\pi$ matches no such transduction for small $p/q$ (Appendix~\ref{sec:computation}), but computation is not proof.
\end{remark}

\subsection{Two expected-false necessary conditions: Hurwitzian numbers and \texorpdfstring{$E$}{E}-values}\label{sec:expected-false}

The measure reduction (Theorem~\ref{thm:measure-reduction}) and the continued-fraction condition (Theorem~\ref{thm:cf}) can be sharpened into two \emph{single} necessary conditions, each a proposition expected to be \emph{false}. This reorganization, suggested by external review, is the cleanest structural framing of the four combinations: a useful necessary condition should be expected false, so that proving it false would yield a contradiction. Both conditions below are classical consequences, not new theorems; their value is organizational.

\paragraph{Hurwitzian numbers.} A \emph{Hurwitzian number} is a real whose continued fraction, from some index on, is formed by interleaving finitely many rational polynomials taking positive integer values, at least one nonconstant. Perron's theorem (see \cite[p.~141]{Stambul}) states that Hurwitzian numbers are preserved under non-degenerate integer-coefficient M\"obius transformations. Since $\ee=[2;1,2,1,1,4,1,1,6,\dots]$ is a Hurwitzian number, rationality of any of $e\pm\pi$, $e\pi$, $e/\pi$ forces $\pi$ to be a Hurwitzian number (each is a M\"obius transformation of $e$ with rational, hence integer-clearable, coefficients). Thus the four cases unify into the single target
\[
\text{``$\pi$ is a Hurwitzian number''}\qquad(\text{expected false}).
\]
This is \emph{stronger} than Theorem~\ref{thm:measure-reduction}: Hurwitzian numbers have polynomially growing partial quotients, hence irrationality measure $2$, so $\irr(\pi)=2$ is a corollary. \textbf{Honest boundary:} no current technique proves that $\pi$ is \emph{not} a Hurwitzian number; even whether the partial quotients of $\pi$ are unbounded is open \cite{Hildebrand2024}. The condition is a reformulation, not a reduction in difficulty.

\paragraph{Siegel \texorpdfstring{$E$}{E}/$G$-values.} Fischler and Rivoal \cite{FischlerRivoal2016} define the ring $\mathbf{E}$ of values of Siegel $E$-functions at algebraic points and the ring $\mathbf{G}$ of values of analytic continuations of $G$-functions at algebraic points. Known: $\overline{\Q}\subset\mathbf{E}$, the units of $\mathbf{E}$ contain $\overline{\Q}^*$ and $\exp(\overline{\Q})$ (so $e,1/e$ are units), and $\pi,1/\pi\in\mathbf{G}$ \cite{FischlerRivoal2014}. They conjecture $\mathbf{E}\cap\mathbf{G}=\overline{\Q}$ \cite{FischlerRivoal2024}. If any of $e\pm\pi,e\pi,e/\pi$ equals an algebraic number $a$, then $\pi$ equals $a-e$, $e-a$, $ae^{-1}$, or $ea^{-1}$ respectively, all in $\mathbf{E}$; hence $\pi\in\mathbf{E}\cap\mathbf{G}=\overline{\Q}$, contradicting Lindemann. So the single sufficient hypothesis
\[
\text{``$\pi\notin\mathbf{E}$''}\qquad(\text{expected true})
\]
would resolve all four cases.
The intersection conjecture and the membership question for $\pi$ are
open. Fischler and Rivoal \cite[Conjecture~1.2]{FischlerRivoal2026}
conjecture that the units of $\mathbf E$ are exactly the numbers
$\alpha e^\beta$, with $\alpha,\beta\in\Qbar$ and $\alpha\ne0$.
To infer that $\pi$ is a non-unit from this conjecture requires excluding
$\pi=\alpha e^\beta$; in particular the case $\beta=1$ would require
excluding an algebraic ratio $\pi/e$.

\paragraph{A conditional consequence of the restricted unit theorem.}
Assume Theorem~1.3 of \cite{FischlerRivoal2026}, which classifies the
units among finite algebraic linear combinations of exponentials of
algebraic numbers. If $a=e+\pi$ or $a=e-\pi$ were algebraic, then $a\ne0$
and $\pi=\pm(a-e)$ would belong to $\mathbf E$. By
Lindemann--Weierstrass, this sum of nonzero terms at the two distinct
exponents $0,1$ cannot equal $\alpha e^\beta$ with
$\alpha,\beta\in\Qbar$ and $\alpha\ne0$. Indeed, if $\beta$ is distinct
from $0,1$, the proposed representation gives a relation between three
exponentials of distinct algebraic numbers. If $\beta$ is $0$ or $1$,
one of the two original coefficients is still nonzero after collecting
terms. Thus the restricted unit theorem yields $1/\pi\notin\mathbf E$
under either algebraicity hypothesis. Equivalently, $1/\pi\in\mathbf E$
would imply that both $e+\pi$ and $e-\pi$ are transcendental, conditional
on the cited preprint theorem. The membership premise is open.

\begin{remark}[Why these are the right targets]
Both conditions are expected false ($\pi$ Hurwitz) or expected true ($\pi\notin\mathbf{E}$) in the direction that would give a contradiction, unlike $\irr(\pi)>2$ or Flint Hills divergence, which are expected \emph{false} and hence cannot serve as attack routes (Section~\ref{sec:conclusion}). The factorial-Cantor digit criterion of Proposition~\ref{prop:window} is, by contrast, only a coordinate rewrite of $\pi=r-e$; the equivalence with $e+\pi\in\Q$ was independently recorded by Yu \cite{Yu2026}. These observations reorganize the necessary conditions; they supply no proof.
\end{remark}

\subsection{Exact rigidity of the series obstruction}

\begin{proposition}[Window condition]\label{prop:window}
Suppose $e+\pi=p/q$ with $p,q\in\N$. For every integer $n\ge q$ there exists $B_n\in\Z$ such that
\[
B_n-\frac1n<n!\,\pi<B_n.
\]
Equivalently, the fractional part satisfies $\{n!\pi\}=1-t_n\in(1-1/n,1)$, where $t_n=\sum_{k>n}n!/k!$.
\end{proposition}

\begin{proof}
Since $n\ge q$, $q\mid n!$, so $n!(e+\pi)=n!p/q\in\Z$. Write $n!e=A_n+t_n$ with $A_n=\sum_{k=0}^n n!/k!\in\Z$ and $t_n=\sum_{k=n+1}^\infty n!/k!$. Subtracting, $n!\pi=B_n-t_n$ with $B_n=n!p/q-A_n\in\Z$.

We estimate $t_n$. Clearly $t_n>1/(n+1)>0$. For the upper bound,
\[
t_n=\frac1{n+1}\Bigl(1+\frac1{n+2}+\frac1{(n+2)(n+3)}+\cdots\Bigr)
<\frac1{n+1}\sum_{j=0}^\infty\frac1{(n+1)^j}=\frac1n.
\]
Thus $B_n-1/n<n!\pi<B_n$, and since $0<t_n<1/n<1$, the fractional part of $n!\pi$ is $1-t_n$.
\end{proof}

\begin{proposition}[Rigidity recurrence]\label{prop:recurrence}
The tails satisfy the exact recurrence $t_{n+1}=(n+1)t_n-1$ for all $n\ge1$.
\end{proposition}

\begin{proof}
Multiply the series by $n+1$:
\[
(n+1)t_n=\sum_{k=n+1}^{\infty}\frac{(n+1)!}{k!}=1+\sum_{k=n+2}^{\infty}\frac{(n+1)!}{k!}=1+t_{n+1},
\]
where the $k=n+1$ term equals $(n+1)!/(n+1)!=1$. Rearranging gives the claim.
\end{proof}

\begin{remark}
Under the hypothesis of Proposition~\ref{prop:window}, the infinitely many integrality conditions $n!\pi+t_n\in\Z$ ($n\ge q$) are not independent: by Proposition~\ref{prop:recurrence} they are all determined by the single value $\{q!\pi\}$. The hypothesis is thus equivalent to a one-parameter orbit hitting integers at every step --- heuristically of probability zero, but internally consistent. The distribution of $\{n!\pi\}$ is completely unknown (even its density in $[0,1]$ is open), so no contradiction is available.
\end{remark}

\begin{remark}[Direction of the factorial criterion --- correction (review)]\label{rem:factorial-direction}
The formalized factorial criterion is an \emph{equivalence}: $e+\pi\in\Q \iff n\{-n!\pi\}\to1$. Hence \textbf{proving the right-hand limit equals $1$ would imply \emph{rationality}, not irrationality}. To prove irrationality one must \emph{exclude} that limit. A clean sufficient target is
\[
\{n!\pi\}\le\tfrac12 \quad\text{for infinitely many }n,
\]
because on those indices $\{-n!\pi\}=1-\{n!\pi\}\ge\tfrac12$, so $n\{-n!\pi\}\ge n/2\to\infty$ and cannot tend to $1$. This is a precise, weaker goal than full distributional control of $\{n!\pi\}$. \textbf{Boundary:} the project currently has \emph{no} tool to control the factorial fractional parts of $\pi$; this must be pursued as an independent exploration line, and the fact that the criterion is formalized does \emph{not} mean the problem is close to a proof.
\end{remark}

\subsection{The four exponentials conjecture: application blocked}

Recall the four exponentials conjecture: if $x_1,x_2$ and $y_1,y_2$ are respectively $\Q$-linearly independent, then at least one of the four numbers $\ee^{x_iy_j}$ is transcendental. To refute $e+\pi\in\Q$ by this conjecture one would need, under that hypothesis, a $2\times2$ matrix all of whose exponential entries are algebraic.

\begin{proposition}\label{prop:four-exp}
Assume $e+\pi=r\in\Q$. Then no $2\times2$ matrix of the form $(\ee^{x_iy_j})$ with all four entries algebraic can be constructed from the known algebraic exponential values; the standard application of the four exponentials conjecture is blocked.
\end{proposition}

\begin{proof}
First we inventory the arithmetic nature of the relevant exponential values under the hypothesis.
\begin{enumerate}
\item $\ee^{ir}$ is transcendental: $ir$ is a nonzero algebraic number, so Theorem~\ref{thm:HL} applies.
\item From Euler's identity and $\pi=r-e$: $\ee^{i(r-e)}=-1$, hence $\ee^{ie}=-\ee^{ir}$, which is transcendental by item 1.
\item $\ee$ and $\ee^\pi$ are transcendental (Theorems~\ref{thm:HL} and~\ref{thm:epi}).
\end{enumerate}
Thus the hypothesis creates no new algebraic exponential values: the available ones remain $\ee^{qi\pi}$ for $q\in\Q$ (roots of unity, including $\ee^0=1$).

Now suppose all four products $x_iy_j$ lay in $\Q\cdot i\pi$, so that all four entries are known algebraic values. Then
\[
\frac{y_1}{y_2}=\frac{x_1y_1}{x_1y_2}\in\Q,
\]
contradicting the $\Q$-linear independence of $y_1,y_2$. Hence no all-algebraic matrix can be built from the available values.
\end{proof}

\begin{remark}
Proposition~\ref{prop:four-exp} blocks the standard application path; it is not a meta-theorem that the conjecture is useless in principle. An identical argument for the (proved) six exponentials theorem shows that at most one of the six values can be controlled as a known algebraic number: among $x_1,x_2,x_3$ at most one lies in $\Q\cdot i\pi$, and for that one, algebraicity of both $\ee^{x_1y_1}$ and $\ee^{x_1y_2}$ would force $y_1,y_2\in\Q$, contradicting their independence. This is far from the ``five algebraic force the sixth transcendental'' configuration.
\end{remark}

\subsection{Translation of the \texorpdfstring{$\irr(\pi)$}{mu(pi)} upper-bound program}

\begin{theorem}[Translation identity]\label{thm:translation}
Let $(L_n)$ be integer linear forms $L_n=a_n+b_n\pi$ ($a_n,b_n\in\Z$) with $|L_n|\le C\rho^n$ and $|b_n|\sim D^n$, where $0<\rho<1<D$. The standard argument yields the upper bound $\irr(\pi)\le M:=1+\log D/\log(1/\rho)$. If $e+\pi=p_0/q_0\in\Q$, the same forms translate into rational approximations to $\ee$ with approximation exponent
\[
\irr_0=\frac{\log(D/\rho)}{\log D}=1+\frac{1}{M-1}.
\]
Since $M\ge2$, one always has $\irr_0\le2$, consistent with $\irr(\ee)=2$; and $\irr_0>2\iff M<2$, which contradicts the irrationality of $\pi$. Hence no contradiction can be extracted from this program under the hypothesis.
\end{theorem}

\begin{proof}
\emph{Step 1: the upper bound.} For any rational $p/q$,
\[
b_n(q\pi-p)-qL_n=-(b_np+a_nq)\in\Z.
\]
Choose $n$ with $q|L_n|<1/2$ (possible since $|L_n|\to0$). If the integer is nonzero, $|b_n|\cdot|q\pi-p|\ge1-q|L_n|\ge1/2$, so $|\pi-p/q|\ge1/(2q|b_n|)$; if it is zero, $|\pi-p/q|=|L_n|/|b_n|$ and the bound holds a fortiori. Taking $n\approx\log q/\log(1/\rho)$ gives $|b_n|\approx D^n=q^{\log D/\log(1/\rho)}$, whence $|\pi-p/q|\ge c\,q^{-1-\log D/\log(1/\rho)}$ for all large $q$, i.e.\ $\irr(\pi)\le M$.

\emph{Step 2: translation.} Under $\pi=p_0/q_0-e$,
\[
q_0L_n=(q_0a_n+p_0b_n)-q_0b_n\,e,
\]
so
\[
\Bigl|e-\frac{q_0a_n+p_0b_n}{q_0b_n}\Bigr|=\frac{|L_n|}{|b_n|}\approx\Bigl(\frac\rho D\Bigr)^n,
\]
with denominators $Q_n=q_0|b_n|\sim D^n$. The exponent of these approximations is
\[
\irr_0=\frac{\log(D/\rho)}{\log D}=1+\frac{\log(1/\rho)}{\log D}=1+\frac1{M-1}.
\]

\emph{Step 3: consistency and vacuity.} Since $\pi$ is irrational, $\irr(\pi)\ge2$, so every valid upper bound satisfies $M\ge2$, hence $\irr_0=1+1/(M-1)\le2$, compatible with $\irr(\ee)=2$ (Lemma~\ref{lem:mue}). Moreover $\irr_0>2\iff M<2$, impossible. Thus the translated approximations never violate the known measure of $\ee$, and the criterion is vacuous.
\end{proof}

\begin{remark}
Numerically: $M=7.10320534$ gives $\irr_0\approx1.16385$; Bai's claimed $M=7.101862832357$ \cite{Bai2026} gives $\irr_0\approx1.16388$. The identity makes explicit that the entire $\irr(\pi)$ upper-bound program is internally consistent with the rationality hypothesis.
\end{remark}

\subsection{Scaling closure and the \texorpdfstring{$\wp$}{wp}-value subquestion}

Let $S$ be the set of values $\wp_\Omega(p)$, where $\Omega$ is a CM lattice with algebraic invariants and $p\in\Qbar\setminus\{0\}$ is not a pole of $\wp_\Omega$.

\begin{proposition}[Scaling closure]\label{prop:scaling}
The set $S$ is closed under multiplication by nonzero algebraic numbers. Consequently, ``$\pi\in S$'' is equivalent to ``$\pi$ is an algebraic multiple of some algebraic-point $\wp$-value,'' i.e., to the existence of an algebraic-coefficient linear relation between $\pi$ and some element of $S$.
\end{proposition}

\begin{proof}
Let $\lambda\in\Qbar\setminus\{0\}$. From the series definition,
\[
\wp_{\lambda\Omega}(\lambda z)=\frac1{\lambda^2z^2}+\sum_{\omega'\in\Omega\setminus\{0\}}\Bigl(\frac1{(\lambda z-\lambda\omega')^2}-\frac1{(\lambda\omega')^2}\Bigr)=\lambda^{-2}\wp_\Omega(z).
\]
The invariants transform as $g_2(\lambda\Omega)=\lambda^{-4}g_2(\Omega)$ and $g_3(\lambda\Omega)=\lambda^{-6}g_3(\Omega)$, which remain algebraic for algebraic $\lambda$; complex multiplication is preserved under scaling; and $\lambda p$ is algebraic when $p$ is. Hence if $v=\wp_\Omega(p)\in S$ and $a\in\Qbar\setminus\{0\}$, taking $\lambda=a^{-1/2}\in\Qbar$ gives $a\cdot v=\wp_{\lambda\Omega}(\lambda p)\in S$. The equivalence is immediate: if $\pi=a v$ with $a$ algebraic then $\pi\in S$; conversely if $\pi\in S$ take $v=\pi$ and $a=1$.
\end{proof}

\begin{corollary}[Conditional reduction; conditional on {\cite{Bertolin2026}}]\label{cor:bertolin}
Assume Bertolin's theorem: for a lattice $\Omega$ with algebraic invariants and CM, if $t_1,\dots,t_s$ are $\Q$-linearly independent algebraic numbers and $p_1,\dots,p_n$ are $k$-linearly independent algebraic numbers ($k=\operatorname{End}(\mathcal E)\otimes\Q$), then $\ee^{t_1},\dots,\ee^{t_s},\wp(p_1),\dots,\wp(p_n)$ are algebraically independent over $\Qbar$. If $\pi=\wp_\Omega(p)$ for some algebraic $p\ne0$, then $\ee$ and $\pi$ are algebraically independent, and $e+\pi$ is transcendental.
\end{corollary}

\begin{proof}
Take $s=n=1$, $t_1=1$, $p_1=p$. A single nonzero algebraic number is $k$-linearly independent, and $1$ is $\Q$-linearly independent. Bertolin's theorem gives algebraic independence of $\ee$ and $\wp_\Omega(p)=\pi$ over $\Qbar$, hence over $\Q$. If $e+\pi$ were algebraic with minimal polynomial $m(t)\in\Q[t]$, then $m(x+y)\in\Q[x,y]$ would be a nonzero polynomial vanishing at $(e,\pi)$, contradiction. Thus $e+\pi$ is transcendental.
\end{proof}

\begin{remark}
The subquestion ``$\pi\in S$?'' is open. The Grothendieck--André period conjecture predicts $\pi\notin S$ (periods and algebraic-point values are motivically separated); Schneider's theorem (Theorem~\ref{thm:schneider}) shows every element of $S$ is transcendental, so no separation from $\pi$ arises there. An affirmative answer together with the preprint hypothesis in Corollary~\ref{cor:bertolin} would imply a solution of Problem~\ref{prob:main}. Huber--W\"ustholz \cite{HuberWustholz2022} have completed the theory of \emph{linear} relations of $1$-periods; the present subquestion concerns \emph{algebraic} relations, which remain conjectural.
\end{remark}

\subsection{The Flint Hills bridge}\label{sec:flint-bridge}

The Flint Hills series is $\sum_{n\ge1}1/(n^3\sin^2n)$; the classical convergence question (see also the later preprint claim in \cite{OpenAI2026}) is recorded in \cite{Alekseyev2011,FormalConjectures}; generalized Flint Hills series are studied in \cite{Dan2026}. We first record the precise published results of Alekseyev \cite{Alekseyev2011}, verified by full-text reading.

\begin{theorem}[Alekseyev, published]\label{thm:alekseyev}
For $u,v>0$:
\begin{enumerate}
\item the terms satisfy $1/(n^u|\sin n|^v)=O(n^{-(u-(\irr(\pi)-1)v-\varepsilon)})$ for every $\varepsilon>0$;
\item if the terms tend to $0$, then $\irr(\pi)\le1+u/v$; in particular, convergence of the Flint Hills series implies $\irr(\pi)\le5/2$;
\item if $\irr(\pi)<1+(u-1)/v$, the series converges. For $(u,v)=(3,2)$ this threshold is $1+(3-1)/2=2$, which is never satisfied since $\irr(\pi)\ge2$; the sufficient condition is vacuous. Alekseyev explicitly notes that divergence of the Flint Hills series implies no nontrivial bound.
\end{enumerate}
\end{theorem}

\begin{remark}[Correction of the record]
An intermediate round of this research attributed the direction ``$\irr(\pi)<5/2\Rightarrow$ convergence'' to Alekseyev, following a secondary source. Full-text verification showed the attribution incorrect: Alekseyev's sufficient condition is the vacuous $\irr(\pi)<2$. The direction $\irr(\pi)<5/2\Rightarrow$ convergence is due to Meiburg \cite{Meiburg2022}, a preprint (math.NT, cited by subsequent literature, no peer-reviewed publication found). The correction is permanent. The bridge below was originally stated conditionally on Meiburg; the Flint Hills case of Meiburg's criterion has since been machine-verified (Section~\ref{sec:lean-f2}), so the bridge is kernel-checked relative to the explicitly listed classical axioms.
\end{remark}

\begin{theorem}[Meiburg, preprint; only the Flint Hills case is formalized in Lean]\label{thm:meiburg}
For any sine-like periodic function with irrational period $\alpha$, $v\ge1$, and $\irr(\alpha)<1+u/v$, the series $\sum n^{-u}P(n)^{-v}$ converges. In particular, $\irr(\pi)<5/2$ implies convergence of the Flint Hills series. The edge case $\irr(\pi)=5/2$ is undecidable by the measure alone: Meiburg constructs $\alpha$ with $\irr(\alpha)=5/2$ for which the series diverges.
\end{theorem}

\begin{remark}[Verification status]
The proof architecture was verified in full across two rounds: (i) the sparsity lemma (good approximation denominators grow as $\Omega(n^{1/(1-\varepsilon_2)})$), proved by forming the difference quotient of two close good approximations and deriving an approximation of exponent $>\irr(\alpha)$, contradicting the definition; the multiplicity-bounding step is compressed but structurally complete; (ii) the partition of integers by approximation exponent into finitely many classes, each contributing a convergent subsum; the key linear function satisfies $f(\irr(\alpha))=1+u/v-\irr(\alpha)$ (independently recomputed) with derivative $1/(v(\irr(\alpha)-1))-1\le0$ (recomputed from the definition; the parsed manuscript contained a typesetting artifact at this formula). No error was found. The general statement remains a preprint; however, the Flint Hills case needed here has been fully machine-verified in Lean (Section~\ref{sec:lean-f2}), so the bridge below no longer depends on the preprint.
\end{remark}

\begin{theorem}[Flint Hills bridge; machine-verified]\label{thm:flint}
If $e+\pi\in\Q$, then the Flint Hills series converges. Contrapositively, if the Flint Hills series diverges, then $\irr(\pi)\ge5/2>2$, and all four combinations $e+\pi$, $e-\pi$, $e\pi$, $e/\pi$ are irrational.
\end{theorem}

\begin{proof}
If $e+\pi\in\Q$, Theorem~\ref{thm:measure-reduction} gives $\irr(\pi)=2<5/2$; Theorem~\ref{thm:meiburg} then yields convergence. For the contrapositive: divergence implies $\irr(\pi)\ge5/2$ (contrapositive of Theorem~\ref{thm:meiburg}), hence $\irr(\pi)>2$, and Corollary~\ref{cor:four-irrational} applies.
\end{proof}

\begin{remark}
Since $\irr(\pi)=2$ is expected, the Flint Hills series is expected to converge, so the divergence target is probably false. The bridge gives a useful implication, but its direction matters: divergence would exclude rationality, whereas convergence alone gives no contradiction. In particular, convergence $\Rightarrow\irr(\pi)\le5/2$ and $e+\pi\in\Q\Rightarrow\irr(\pi)=2$ are compatible. The verified Meiburg derivation supplies the displayed rationality-to-convergence implication relative to the listed classical assumptions; no universal non-derivability statement is asserted.
\end{remark}

\subsection*{Systematic attempts and obstruction analysis}

Each route below was carried to a rigorous terminal verdict. The verdicts concern the specific constructions; they are not meta-theorems that no proof along the general line can exist.

\subsection{Route A: series and integer sequences}\label{sec:obstructions}

The construction and its necessary condition are Proposition~\ref{prop:window}; the rigidity is Proposition~\ref{prop:recurrence}. \emph{Verdict: blocked.} The distribution of $\{n!\pi\}$ is completely unknown; even density in $[0,1]$ is open. A refutation of the window condition would prove Problem~\ref{prob:main}; the condition is a necessary one and is not known to fail (it is not claimed to be equivalent in difficulty). The denominator-clearing mechanism is specific to the exponential series: the known fast series for $\pi$ (Machin-type arctangent, Chudnovsky) have partial-sum denominators that are products of powers and odd integers, and multiplying the tail by them does not produce integers --- so this particular factorial-clearing argument does not transfer directly. This does not exclude series proofs of the irrationality of $\pi$ or other denominator-clearing constructions.

\subsection{Route B: Niven-type integrals}\label{sub:route-b}

Transplant the polynomial $f(x)=x^n(a-bx)^n/n!$ to the endpoint $\alpha=e+\pi=a/b$. The first half succeeds: $f(\alpha-x)=f(x)$ (since $a=b\alpha$), so $f^{(k)}(0),f^{(k)}(\alpha)\in\Z$ exactly as in Theorem~\ref{thm:niven}; endpoint derivative integrality holds for \emph{any} rational endpoint.

\emph{Verdict: blocked at the kernel.} Niven's proof uses $K=\sin x$, relying on $K(0)=K(\pi)=0$, $K'(0)=1$, $K'(\pi)=-1$, and $K''=-K$. For $\alpha=e+\pi$: $\sin\alpha=-\sin e\ne0$, since $0<e<\pi$ implies $\sin e>0$. Thus the endpoint cancellation used in Niven's proof fails for this kernel. Forcing $K(x)=\sin(\pi x/\alpha)$ gives the correct zeros but $K''=-(\pi/\alpha)^2K$, so the telescoping construction acquires coefficients involving powers of $(\alpha/\pi)^2=(1+e/\pi)^2$, whose arithmetic nature is unknown; integrality collapses. No known kernel vanishes at $e+\pi$ with arithmetically controlled derivatives.

\subsection{Route C: algebraic and transcendental methods}\label{sub:route-c}

\emph{Lindemann--Weierstrass attempt.} Suppose $S=e+\pi$ is algebraic. From $\pi=S-e$ and Euler's identity,
\[
\ee^{iS}\cdot\ee^{-ie}=\ee^{i(S-e)}=\ee^{i\pi}=-1,\qquad\text{so}\qquad \ee^{iS}=-\ee^{ie}.
\]
The left exponent $iS$ is nonzero algebraic, so $\ee^{iS}$ is transcendental (Theorem~\ref{thm:HL}); the right exponent $ie$ is transcendental, so the theorem gives no information about $\ee^{ie}$. The relation ties an exponential at an algebraic point to one at a transcendental point --- outside every proved theorem. With $S=e+\pi$, the identity holds without an algebraicity assumption on $S$; the extra hypothesis that $S$ is algebraic produces no contradiction by this argument.

\emph{Baker's theory.} Baker's theory of linear forms in logarithms \cite{Baker} handles $\beta_0+\sum\beta_j\log\alpha_j$ with algebraic $\alpha_j,\beta_j$. Here $\pi=\log(-1)/i$ fits the framework, but $\ee$ has no known representation as a logarithm of an algebraic number. \emph{Verification note:} the popular assertion ``$\ee$ is not the logarithm of any algebraic number'' is not a theorem: $\ee=\log\alpha$ with $\alpha$ algebraic is equivalent to $\ee^\ee$ being algebraic, and the transcendence of $\ee^\ee$ is itself open. The accurate statement is ``no known representation.''

\emph{Non-closure.} $\pi$ and $1-\pi$ are both transcendental with sum $1$ (immediate from Theorem~\ref{thm:HL}); individual transcendence logically cannot yield irrationality of the sum. Independence information is indispensable. (Unaccepted preprints claiming irrationality of $e\pi$ \cite{CarellaEpi} are recorded as negative leads only and are not used.)

\emph{Verdict: blocked.} The cited transcendence theorems do not yield a contradiction in this direct application. No universal impossibility claim is made about additive constructions.

\subsection{Route D: Diophantine approximation}

Criterion: if integers $P_n$ and positive integers $Q_n$ satisfy $P_n/Q_n\ne\alpha$ and $|Q_n\alpha-P_n|\to0$, then $\alpha$ is irrational. Indeed, for $\alpha=A/B$ a nonzero such linear form has absolute value at least $1/B$.

Combining convergents $a_n/b_n$ of $\ee$ and $c_n/d_n$ of $\pi$ with the product denominator gives
\[
\left|b_nd_n(e+\pi)-(a_nd_n+c_nb_n)\right|
<\frac{d_n}{b_n}+\frac{b_n}{d_n}.
\]
The right side is at least $2$. This means only that this triangle-inequality estimate does not establish smallness: an upper bound bounded away from zero is not a lower bound for the actual error. It does not exclude cancellation of signed errors, reduction of the combined fraction, or a better common-denominator construction. If $L_n=\operatorname{lcm}(b_n,d_n)$, a sharper elementary upper bound for the corresponding integer linear form is
\[
L_n\left(b_n^{-2}+d_n^{-2}\right).
\]
The present argument does not show that this tends to zero, nor that the form is nonzero infinitely often. Either issue requires further work. In particular, no necessity of $\irr(\pi)>2$ follows from the failure of the displayed estimate. Theorem~\ref{thm:measure-reduction} provides a separate sufficient condition, not a classification of all approximation proofs.

\subsection{Bridge identities: four candidates examined}\label{sub:bridges}

\begin{enumerate}
\item \emph{Gauss integral.} $\int_0^\infty \ee^{-x^2}dx=\sqrt\pi/2$ (symbolically verified), i.e.\ $\pi=4I^2$. Obstruction: quadratic in $I$; the paradigm requires integer \emph{linear} forms in $1,\pi$; no linearization with controlled denominators is known, and $I$ has no factorial-clearing series.
\item \emph{Stirling.} $\lim_{n\to\infty}n!\ee^n/(n^n\sqrt{2\pi n})=1$ (symbolically verified). Obstruction: an asymptotic identity, not an exact algebraic relation; under $e=r-\pi$ it becomes a tautology about powers of $r-\pi$; the error term has no Diophantine control.
\item \emph{Mixed integral.} $\int_0^\pi \ee^x\sin x\,dx=(\ee^\pi+1)/2$ (symbolically verified). Produces $\ee^\pi$, already known transcendental; no new information.
\item \emph{Cosine integral.} $\int_0^\infty\cos x/(1+x^2)\,dx=\pi/(2\ee)$ (symbolically verified). Under the hypothesis, $\ee=r/(1+2J)$ with $J$ the integral: a multiplicative relation with an arithmetically uncontrolled integral.
\end{enumerate}

\begin{remark}[Mechanistic synthesis]
The exponential and period identities examined here linking $\ee$ and $\pi$ --- Euler's identity, $\ee^\pi=(-1)^{-i}$, the partial fraction expansion $\pi\coth(\pi z)=1/z+\sum_{n\ge1}2z/(z^2+n^2)$ (whose Fourier expansion involves $\ee^{2\pi a}$), and modular forms ($\ee^{2\pi i\tau}$) --- illustrate several different ways these constants can occur together. None of the four tested integral or asymptotic identities supplies the required small nonzero integer linear forms. The number $\ee=\ee^1$ is an exponential value at an \emph{algebraic} point; no known identity connects algebraic-point exponential values to periods additively. This is the mechanistic explanation of the missing bridge. (An observation about the examined constructions, not a theorem about all possible identities.) A contrasting success: Powers \cite{Powers2025} proves algebraic independence of $\gamma^{(n)}+\delta^{(n)}/\ee$ combinations using Hardy's equation --- two individually unknown constants can be conquered when a structural bridge identity exists.
\end{remark}

\subsection{Mahler classification}\label{sub:mahler}

It is classical that $\ee$, $\pi$, and $\log2$ are not $U$-numbers (for instance, finite irrationality measure excludes Liouville numbers, and Baker's results on Mahler's classification exclude $U$-numbers \cite{Baker}). Mahler classes are invariant under rational M\"obius transformations (heights distort by bounded factors depending on the transformation and degree; the defining $\limsup$ exponents are unchanged) \cite{Baker,HuberWustholz2022}. The hypothesis $e+\pi\in\Q$ would force $\pi$ into the same class as $\ee$; both are already known non-$U$. \emph{Verdict: no separation.}

\subsection*{Conditional frameworks and the frontier}

\subsection{Schanuel's conjecture: complete conditional proof}\label{sec:conditional}

\begin{theorem}[Schanuel, conjectural]\label{thm:schanuel}
If $z_1,\dots,z_n\in\C$ are linearly independent over $\Q$, then $\trdeg_\Q\Q(z_1,\dots,z_n,\ee^{z_1},\dots,\ee^{z_n})\ge n$.
\end{theorem}

\begin{theorem}[Conditional]\label{thm:schanuel-consequence}
Assume Theorem~\ref{thm:schanuel}. Then $\ee$ and $\pi$ are algebraically independent over $\Q$, and $e+\pi$ is transcendental.
\end{theorem}

\begin{proof}
\emph{Step 1: linear independence.} Take $z_1=1$, $z_2=i\pi$. If $a+b\cdot i\pi=0$ with $a,b\in\Q$ and $b\ne0$, then $\pi=ai/b$ is algebraic, contradicting Theorem~\ref{thm:HL}. Hence $b=0$ and $a=0$.

\emph{Step 2: transcendence degree.} By the conjecture, $\trdeg_\Q\Q(1,i\pi,\ee,\ee^{i\pi})\ge2$. Since $\ee^{i\pi}=-1$, this field is $\Q(\ee,i\pi)$, so $\trdeg_\Q\Q(\ee,i\pi)\ge2$. Adjoining the algebraic number $i$ does not change transcendence degree: $\Q(\ee,i\pi,i)=\Q(\ee,\pi,i)$, and both are obtained by adjoining the algebraic element $i$ to the respective fields. Hence $\trdeg_\Q\Q(\ee,\pi)=\trdeg_\Q\Q(\ee,i\pi)\ge2$. Thus $\ee$ and $\pi$ are algebraically independent.

\emph{Step 3: transcendence of the sum.} If $e+\pi$ were algebraic with minimal polynomial $m(t)\in\Q[t]$, then $m(x+y)\in\Q[x,y]$ is nonzero and vanishes at $(\ee,\pi)$, contradicting algebraic independence. Hence $e+\pi$ is transcendental (and irrational).
\end{proof}

\begin{remark}[Sufficient assumption hierarchy]
The conjecture is proved only for $n=1$ (equivalent to Hermite--Lindemann); the general case remains open \cite{Baker,MO33817} (while its function-field analogue was established by Ax \cite{Ax1971}). The instance at $(1,i\pi)$ is \emph{equivalent} to algebraic independence of $\ee,\pi$; the hierarchy is
\[
\text{Schanuel}(1,i\pi)\iff \ee,\pi\text{ algebraically independent}\implies e+\pi\text{ transcendental}\implies e+\pi\text{ irrational}.
\]
This is a sufficient assumption hierarchy, not a minimality theorem. The argument does not show that irrationality of $e+\pi$ requires algebraic independence or this Schanuel instance. The four-exponentials application tested in Proposition~\ref{prop:four-exp} does not supply the required input. No strictly weaker named conjecture is known to imply the instance; proving $\ee,\pi$ algebraically independent is a strictly stronger goal than merely excluding $e+\pi\in\Q$. We therefore call it ``a clean sufficient hypothesis,'' not ``the minimal unproven assumption.''
\end{remark}

\subsection{Ax--Schanuel and arithmetic geometry}\label{sub:ax}

Ax proved the differential/function-field analogue in 1971 \cite{Ax1971}. Mok--Pila--Tsimerman (2019) proved Ax--Schanuel for Shimura varieties \cite{MPT2019}; Bakker--Tsimerman extended it to variations of polarized Hodge structures \cite{BakkerTsimerman2019}. These are function-field results, not equivalent to the number-field conjecture. Published analysis notes that the naive number-theoretic analogues of the functional results are false: the first counterexample comes from periods, whose Riemann--Legendre relations are quadratic and cannot be traced to the relevant hull structures \cite{CambridgeFF}. This is documented evidence that periods and function values obey different structural laws.

\subsection{Zilber's model-theoretic program}\label{sub:zilber}

Zilber axiomatized exponential fields satisfying Schanuel's conjecture and proved uniqueness (up to isomorphism) in cardinality continuum. If that field were isomorphic to $\C$, the conjecture would follow. The isomorphism question is open; o-minimality is a property of ordered structures, and the real structure underlying $\C_{\exp}$ is not o-minimal (the integers are definable), so the mature model theory of the real exponential field (Wilkie's theorem) does not transfer \cite{MO33817}.

\subsection{The motivic frontier (2026)}\label{sub:motivic}

Bertolin's preprint \cite{Bertolin2026} (Corollary~\ref{cor:bertolin}; see also the Schanuel-style conjecture for $1$-motives in \cite{BertolinMotives}) simultaneously extends Lindemann--Weierstrass and Philippon--W\"ustholz (1983); it is a special case of the semi-elliptic LW conjecture, which follows from the Grothendieck--André period conjecture for $1$-motives. The proof mechanism: for the $1$-motive $M=[u:\Z\to\mathbb G_m^s\times\mathcal E^n]$ with $u(1)=(\ee^{t_i},P_j)$, the Lie algebra of the unipotent radical splits as $B\oplus Z(1)$; polynomial relations have degree at most $1$ and separate into elliptic and toric parts, excluded by Philippon--W\"ustholz and Lindemann--Weierstrass respectively. The interface with Problem~\ref{prob:main} is exactly the subquestion $\pi\in S$ (Proposition~\ref{prop:scaling}): the framework's value side accepts algebraic-point exponential values (including $\ee=\ee^1$) and algebraic-point $\wp$-values, while $\pi$ is a period. Waldschmidt \cite{Waldschmidt2025} proves a strong Schanuel property for $\ee^z,\wp,\zeta,\sigma$ and Serre functions for \emph{almost all} tuples --- a measure-theoretic result that cannot be specialized to specific constants. The Legendre relation expresses $\pi i$ as a bilinear combination of periods and quasi-periods \cite{DLMF23}: $\pi$ has exact period expressions, while $\ee$ has no known expression in the period ring --- the gap is precisely localized.

\subsection{The exclusion map}\label{sub:exclusion}

\begin{center}
\small\begin{tabularx}{\textwidth}{@{}p{0.25\textwidth}p{0.18\textwidth}X@{}}
\toprule
Option & Verdict & Obstruction coordinate \\
\midrule
Series route & blocked & distribution of $\{n!\pi\}$ unknown (Prop.~\ref{prop:window}, \ref{prop:recurrence}) \\
Integral route & blocked & no controlled kernel at $e+\pi$ (\S\ref{sub:route-b}) \\
Transcendence theorems, direct & blocked & bridge $\ee^{iS}=-\ee^{ie}$ outside all proved theorems \\
Approximation route & estimate insufficient & cancellation and denominator savings unexamined \\
$\irr(\pi)$ program translation & consistent & $\irr_0=1+1/(M-1)\le2$ (Thm.~\ref{thm:translation}) \\
Mahler classification & no separation & $\ee,\pi$ both non-$U$ \\
Bridge identities (4 candidates) & inconclusive & no controlled integer forms obtained \\
Four/six exponentials & application blocked & Prop.~\ref{prop:four-exp} \\
Conditional frameworks & unproven conjectures & sufficient input: Schanuel$(1,i\pi)$ instance \\
Motivic framework (2026) & open subquestion & $\pi\in S$? (Prop.~\ref{prop:scaling}), expected negative \\
\bottomrule
\end{tabularx}
\end{center}

\subsection{Surviving targets and their logical direction}

The first three sufficient conditions recorded earlier remain logically valid as implications: $\irr(\pi)>2$; divergence of the Flint Hills series; and membership $\pi\in S$ together with the cited conditional motivic theorem. \textbf{However, the first two are expected to be \emph{false} and should be abandoned as attack routes.} The external claim $\irr(\pi)=2$ (OpenAI, \cite{OpenAI2026}, with a Lean formalization but no independent specialist verification as of October 2026) would, if confirmed, directly negate both $\irr(\pi)>2$ and Flint Hills divergence; even if that claim fails, $\irr(\pi)=2$ is the universally expected value, so neither condition is a plausible target. The third, $\pi\in S$, is likewise expected false and does not lower the difficulty. A useful necessary condition must be expected \emph{false} in the direction that yields a contradiction; the two such conditions isolated in Section~\ref{sec:expected-false} (``$\pi$ is a Hurwitzian number'', ``$\pi\in\mathbf{E}$'') satisfy this criterion, whereas the three above do not. None of these conditions is established here. A new identity or approximation construction is a research proposal, not a proposition that should be called ``expected false'' without justification.

The supplement gives more explicit targets. Retain its notation $\delta_j=Q_je-P_j$, and for each positive integer $B$ put
\[
M_j(B)=\operatorname{round}(BQ_j\pi),\quad
C_j=P_jQ_{j+1}-P_{j+1}Q_j=2(-1)^{j+1},
\]
\[
T_j(B)=M_j(B)Q_{j+1}-M_{j+1}(B)Q_j,\qquad
D_j(B)=T_j(B)+BC_j.
\]
By the recorded adjacent-index criterion, $B(e+\pi)\in\Z$ if and only if $D_j(B)=0$ eventually, and also if and only if it vanishes infinitely often. Consequently
\[
e+\pi\notin\Q
\iff \forall B\ge1\ \forall N\ \exists j\ge N:\ D_j(B)\ne0.
\]
This follows directly from the existing forward implication and the infinite-match converse. The equivalence is now also kernel-checked in \texttt{EpiFormal/\allowbreak{}Escape.lean} (\S\ref{sec:lean-escape}), using the proved selected-convergent error estimate and only base axioms. Equivalently (\S\ref{sec:lean-escape}), for every $B$ all sufficiently large $j$ are mismatches; the two forms are interchangeable because infinitely many matches would force eventual matching. This replaces a demand for an explicit last match by an asymptotic non-vanishing statement, and does not reduce the logical difficulty of the open problem.

Similarly, write $R_j(B)=BQ_j\pi-\operatorname{round}(BQ_j\pi)$. It would suffice to prove, for every $B\ge1$, that $R_j(B)=-B\delta_j$ for only finitely many $j$ (equivalently, $R_j(B)\ne-B\delta_j$ for all sufficiently large $j$); by the recorded criterion, infinitely many failures are in fact equivalent, since infinitely many matches would force eventual matching. The previously proposed positive-product estimate $R_j(B)\delta_j>0$ infinitely often is a stronger sufficient condition. A separate sufficient condition is $\|BQ_j\pi\|\ge\varepsilon_B$ infinitely often for some $\varepsilon_B>0$, since $\delta_j\to0$. Neither estimate is established for $\pi$.

The factorial criterion has the direction
\[
e+\pi\in\Q\iff n\{-n!\pi\}\longrightarrow1.
\]
To prove irrationality by this route one must \emph{disprove} this limit; proving it would prove rationality. For example, infinitely many indices with $\{n!\pi\}\le1/2$ would suffice: $\pi$ is irrational, so these fractional parts are positive and $n\{-n!\pi\}\ge n/2$ along that subsequence. This escape condition is sufficient, not claimed equivalent, and is unproved here.

\subsection{Adaptive windows and factorial digit quantization}

\begin{theorem}[An adaptive exclusion strip]\label{thm:adaptive-strip}
Let $x$ be irrational, let $p/d$ be one of its reduced convergents,
and let $s$ be its next convergent denominator. Put $q=Q_j$ and
$P=P_j$. Let $\ell=\gcd(qd,dP+qp)$. If $B(e+x)=A\in\Z$ for a
positive integer $B$ and
\begin{equation}\label{eq:adaptive-budget}
 B\left(\frac{d}{2jq}+\frac{q}{s}\right)\le\ell,
\end{equation}
then $p/d=A/B-P/q$, in particular $d\mid Bq$.
For $j>2B^2$ put
\[
 t_\pm(j,B)=\frac{j}{B}\left(1\pm\sqrt{1-\frac{2B^2}{j}}\right).
\]
Under the same rationality assumption there is no convergent
denominator $d$ of $x$ in $[qt_-(j,B),qt_+(j,B)]$.
In particular, the simpler strip $[2Bq,jq/B]$ is forbidden whenever
it is nonempty.
\end{theorem}
\begin{proof}
The integer $Z=Aqd-B(Pd+qp)$ is divisible by $\ell$ and satisfies
\[
 |Z|\le B\bigl(d|qe-P|+q|dx-p|\bigr)
 < B\left(\frac{d}{2jq}+\frac{q}{s}\right).
\]
Here the second error bound is the standard next-denominator bound
for regular continued fractions. Under \eqref{eq:adaptive-budget}
one has $|Z/\ell|<1$, hence $Z=0$ and reduction gives $d\mid Bq$.
If $d=qt$ is in the displayed strip, then $d>q$ and $s>d$, and
\[
 B\left(\frac{t}{2j}+\frac1t\right)\le1
 \quad\Longleftrightarrow\quad t_-(j,B)\le t\le t_+(j,B).
\]
The smaller root is $2B/(1+\sqrt{1-2B^2/j})>B$, so $d>Bq$,
contrary to divisibility. For the simpler strip the two terms of the
budget are each at most one half after multiplication by $B$, while
$d\ge2Bq>Bq$. This also covers $j=2B^2$.
\end{proof}
\begin{corollary}[A denominator-only escape test]\label{cor:strip-escape}
If along $j\to\infty$ one can choose convergent denominators $d$ of
$\pi$ with $d/Q_j\to\infty$ and $d/(jQ_j)\to0$, then $e+\pi$ is
irrational. For a fixed $B$, a single pair $(d,s)$ obeying
\eqref{eq:adaptive-budget} and $d\nmid BQ_j$ excludes every rational
value with denominator $B$.
\end{corollary}
A joint check of two consecutive $\pi$ denominators can exploit
the full cancellation factor $\ell$ without requiring two witnesses
in a single Dirichlet window.
For any fixed $B$ the identity in Theorem~\ref{thm:adaptive-strip}
also gives an exact check on the alleged numerator. The infinite
interior-strip condition for $\pi$ is open here. The finite data in
\texttt{verify\_research\_advance.py} check both the joint budget and
its simpler strip on a strict Machin enclosure and on rational-shift
controls. The strict integer step and the two-half budget are
audited in \texttt{EpiFormal/ResearchAdvance.lean} (integrated into
\texttt{EpiFormal.lean} and \texttt{Audit.lean}).


\begin{theorem}[Factorial digit quantization]\label{thm:factorial-quantization}
Let $x\in\R$, $t_n=\{n!x\}$, $u_n=nt_n$, and
$a_{n+1}=\lfloor(n+1)t_n\rfloor$. If $(u_n)$ is bounded, its set of
subsequential limits is exactly the set of integers occurring
infinitely often among the $a_{n+1}$. For every integer $k\ge0$,
\[
 \limsup_{n\to\infty}|u_n-k|\in\Z_{\ge0}\cup\{+\infty\},
 \qquad
 \limsup_{n\to\infty}|u_n-k|<1
       \quad\Longleftrightarrow\quad x-ke\in\Q.
\]
In the latter case $a_{n+1}=k$ eventually and $u_n\to k$.
\end{theorem}
\begin{proof}
Since $(n+1)n!=(n+1)!$, the fractional parts satisfy
$(n+1)t_n=a_{n+1}+t_{n+1}$. Consequently
\[
 u_n-a_{n+1}=\frac{n t_{n+1}-a_{n+1}}{n+1}.
\]
If $u_n$ stays bounded, then so do $a_{n+1}$ and $nt_{n+1}$,
so the difference is $O(1/n)$. A bounded sequence of integers has
precisely its infinitely recurrent values as subsequential limits.
This also quantizes the limsup of the distance to $k$. If $u_n$ is
unbounded, that limsup is $+\infty$.

A limsup smaller than $1$ forces the sole recurrent digit to be $k$.
Iteration of the fractional-part recurrence, with its bounded
remainder divided by the product of successive bases, then gives
$t_n=k\sum_{\ell>n}n!/\ell!$ for all sufficiently large $n$.
The finite exponential sum multiplied by $n!$ is an integer, so
$n!(x-ke)\in\Z$. Conversely, if $x-ke$ is rational, the same
expression equals $t_n$ for all large $n$: its value is nonnegative
and tends to zero, and $n!(x-ke)$ is an integer. The expression gives
$nt_n\to k$, proving the theorem.
\end{proof}
For $x=-\pi$ and $k=1$, irrationality of $e+\pi$ is equivalent to
$\limsup_n|n\{-n!\pi\}-1|\ge1$ and also to infinitely many factorial
digits different from $1$. Every fixed separation strictly between
zero and one can thus be used in the subsequence criterion.
This strengthens the quantitative formulation without establishing
which alternative holds for $\pi$. The integer separation step is
audited in \texttt{ResearchAdvance.lean}; the certificate tests the
digit rule at $\pi$ and at the controls $x=ke+r$, $r\in\Q$.


\section{Conclusion and open problems}\label{sec:conclusion}

This section records what the assembled results decide, the
companion metric context, and the research priorities. Whether the
actual $\pi$ violates the slope condition of
Corollary~\ref{cor:recover-slope} remains unresolved.
\paragraph{Companion metric results.} The revised companion
\cite{CompanionGauge} gives exact Hausdorff dimension $1-\alpha$ for
prescribed admissible limits and sublevels of $j^\alpha\|Q_jx\|$ when
$0\le\alpha<1$, an exact Hausdorff gauge and three-way positive-level
trichotomy for each positive critical spectral level (on both $\R$ and
$\R/\Z$, and for general polynomial-growth recurrences $r\ge1$), and
exact Hausdorff dimensions ($1$ and $1/2$) for Hermite continued-fraction
window constraints; these metric theorems are now kernel-checked in
Lean~4 (base axioms only, integrated into \texttt{EpiFormal.lean} and
\texttt{Audit.lean}). In particular the unweighted zero-set has
dimension one, whereas the weighted zero-set at exponent one
equals the countable group $L$; the counterexample to their
equality is kernel-checked and integrated in
\texttt{EpiFormal/UnweightedCorrection.lean}. Neither dimension
calculation decides membership for the prescribed point $\pi$.
\subsection{Research priorities and limitations}

The immediate priorities are to study exact identities for $D_j(B)$ and the signed rounding errors, then seek estimates using a specific analytic representation of $\pi$. The recurrence for $Q_j$ and the universal bound $|R_j(B)|\le1/2$ alone cannot exclude matching: substituting $x=A/B-e$ for $\pi$ supplies real examples with eventual matching. Generic or almost-everywhere distribution theorems cannot be specialized to $\pi$ without an additional argument.

The obstruction analysis is limited to the displayed constructions. It proves neither exhaustion of elementary mathematics nor a universal no-go theorem for series, integrals, Pad\'e approximation, or integer linear forms. A finite no-go search must retain its quantified family and parameter bounds. The algebraic-independence reduction is a clean sufficient framework, but it asks for a stronger conclusion than irrationality of the sum. The present contribution is an auditable collection of reductions and method-specific limitations, not a resolution of the central problem.

\paragraph{Current research paths.}
The eventual rounding mismatch for each fixed denominator $B$ remains
the arithmetic target. One sufficient route is to produce infinitely many
windows with two primitive, distinct witnesses whose denominators are
coprime. The factorial route instead requires a fixed separation from~$1$
of $n\{-n!\pi\}$ along an infinite subsequence. The signed linear-form
route uses the common divisor and the next convergent denominator of
$\pi$. The adaptive-strip test also uses these denominators directly
and has a sufficient strip independent of parity. The finite
certificates establish neither the infinite strip condition nor the
factorial digit escape for $\pi$.

The quadratic-orbit route has a fixed-target test: vanishing
$c_\varphi(\pi)$ would separate $\pi$ from the rational affine models.
The finite-phase, trace and rate results establish the model-side
conclusions. The vanishing condition still requires information about the
continued-fraction digits of $\pi$. The model-side trace test now includes conductor primes,
and the density theorem gives the leading common count across
quadratic orders in a fixed field. Obtaining extremal formulas
uniform in the rational denominator remains a separate problem. The companion
paper's finite-state gauge and frequency theorems concern the corresponding
metric problem on sets of targets.

\section*{Use of AI and formal verification}
This work was carried out with substantial assistance from
generative AI systems (including ChatGPT and Codex), both for
developing and drafting mathematical arguments and text,
preparing the numerical certificates, and writing the Lean~4
formalization. The results marked as kernel-checked were verified
by the Lean~4 kernel, and their axiom dependencies were audited
with \texttt{\#print axioms} in \texttt{Audit.lean} (937 audited
declarations project-wide across 134 Lean files as of 2026-11-22, recorded in
\texttt{verification/\allowbreak audit\_\allowbreak final\_\allowbreak 20261122.log};
914 depend only on the Lean base axioms \texttt{propext},
\texttt{Classical.choice}, and \texttt{Quot.sound}, and 23 use the
registered transcendence/irrationality-measure whitelist axioms
\texttt{e\_transcendental}, \texttt{pi\_transcendental}, and
\texttt{mu\_e\_eq\_two}; zero project-added axioms and zero
\texttt{sorryAx}). The remaining
paper proofs---including the window and spectrum computations of
Sections~\ref{sec:windows}--\ref{sec:quadratic} and the obstruction
analysis---are not machine-checked; they have been checked by
line-by-line re-derivation within the same AI-assisted workflow
and cross-checked against the reproducible scripts and Lean
audits of Appendices~\ref{sec:computation}--\ref{app:repro}; they
have not yet been examined by an independent human referee.
\begin{thebibliography}{99}

\bibitem{Hermite1873} C.~Hermite, \emph{Sur la fonction exponentielle}, C.~R.~Acad.~Sci.~Paris~77 (1873), 18--24, 74--79, 226--233, 285--293.
\bibitem{Lindemann1882} F.~Lindemann, \emph{\"Uber die Zahl $\pi$}, Math.~Ann.~20 (1882), 213--225.
\bibitem{OEIS} OEIS Wiki, \emph{Pi+e}, \url{https://oeis.org/wiki/Pi\%2Be}. (No nonzero integer pair $m,n$ with $m\pi+n\ee$ of known rationality status.)
\bibitem{Roth1955} K.~F.~Roth, \emph{Rational approximations to algebraic numbers}, Mathematika 2 (1955), 1--20.
\bibitem{GS1934} A.~O.~Gelfond and Th.~Schneider, independent solutions of Hilbert's seventh problem, 1934--1935.
\bibitem{Duke} W.~Duke, \emph{The zeros of the Weierstrass $\wp$-function}, preprint, \url{https://www.math.ucla.edu/~wdduke/preprints/zeros.pdf}. (Statement of Schneider's elliptic theorem.)
\bibitem{Nesterenko1996} Yu.~V.~Nesterenko, \emph{Modular functions and transcendence questions}, Sb.~Math.~187 (1996), 1319--1348.
\bibitem{Tijdemann1971} R.~Tijdeman, \emph{On the algebraic independence of certain numbers}, Indag.~Math.~33 (1971), 146--162. (The precise ``at least two of $\ee,\pi,\ee^\pi$'' corollary is not verified from the original; the safe attribution for $\pi,\ee^\pi$ algebraic independence is Nesterenko \cite{Nesterenko1996}.)
\bibitem{CNU2024} Z.~Liang and Y.~Wang, \emph{Irrationality measures of several types of irrational numbers}, J.~Capital Normal Univ.~(Nat.~Sci.), 2024. (Theorem 1: $\irr(\ee^r)=2$ for $r\in\Q^*$; Bundschuh's result.)
\bibitem{ZZ2020} D.~Zeilberger and W.~Zudilin, \emph{The irrationality measure of $\pi$ is at most $7.103205334137\dots$}, Moscow J.~Combin.~Number Theory~9 (2020), 407--419, arXiv:1912.06345.
\bibitem{Alekseyev2011} M.~A.~Alekseyev, \emph{On convergence of the Flint Hills series}, arXiv:1104.5100 (2011). Full text read; Theorems 2, 5, Corollaries 3, 4, 6 verified.
\bibitem{Bai2026} Y.~Bai, \emph{The irrationality measure of $\pi$ is at most 7.101862832357}, arXiv:2609.11276 (2026). Preprint; abstract verified; local-minimizer claim only.
\bibitem{CarellaMeasure} N.~A.~Carella, \emph{Irrationality Measure of Pi}, arXiv:1902.08817. Preprint, not accepted; recorded as a negative lead only.
\bibitem{CarellaEpi} N.~A.~Carella, \emph{The Product $e\pi$ Is Irrational}, arXiv:1706.08394. Preprint, not accepted; public criticism by J.~D'Aurizio on Math StackExchange.
\bibitem{Panti2017} G.~Panti, \emph{Slow continued fractions, transducers, and the Serret theorem}, arXiv:1706.00698; J.~Number Theory (2018).
\bibitem{Raney1973} G.~N.~Raney, \emph{On continued fractions and finite automata}, Math.~Ann.~206 (1973), 265--283.
\bibitem{Baker} A.~Baker, \emph{Transcendental Number Theory}, Cambridge, 1975. (Linear forms in logarithms; Fields Medal 1970.)
\bibitem{Bertolin2026} C.~Bertolin, \emph{Algebraic independence of the exponential and Weierstrass $\wp$-functions}, arXiv:2609.15294 (2026). Preprint; abstract, introduction, and proof opening read in the research audit. Related published work: \emph{Third kind elliptic integrals and 1-motives}, J.~Pure Appl.~Algebra 224 (2020), no.~10, 106396, arXiv:1905.07247.
\bibitem{HuberWustholz2022} A.~Huber and G.~W\"ustholz, \emph{Transcendence and Linear Relations of $1$-Periods}, Cambridge Tracts in Mathematics 227, Cambridge Univ.~Press, 2022.

\bibitem{FormalConjectures} google-deepmind/formal-conjectures, issue \#3484, \emph{Convergence of the Flint Hills series}, closed 2026-03-17; Lean statement in repository; formalizability rating 1/5.
\bibitem{Meiburg2022} A.~Meiburg, \emph{Bounds on Irrationality Measures and the Flint-Hills Series}, arXiv:2208.13356 (2022). Preprint (math.NT); full text parsed and proof structure verified (sparsity lemma, partition argument, $f(\mu)$ algebra recomputed).
\bibitem{Powers2025} M.~R.~Powers, \emph{On the Algebraic Independence of a Set of Generalized Constants}, arXiv:2508.12123. Preprint.
\bibitem{Ax1971} J.~Ax, \emph{On Schanuel's conjectures and Skolem's method}, Ann.~of Math.~(2) 93 (1971), 252--268.
\bibitem{MPT2019} N.~Mok, J.~Pila and J.~Tsimerman, \emph{Ax-Schanuel for Shimura varieties}, Ann.~of Math.~(2) 189 (2019), 945--978.
\bibitem{BakkerTsimerman2019} B.~Bakker and J.~Tsimerman, \emph{The Ax-Schanuel conjecture for variations of Hodge structures}, Invent.~Math.~217 (2019), 77--94, arXiv:1712.05088.
\bibitem{Bundschuh1971} P.~Bundschuh, \emph{Irrationalit\"atsma\ss e f\"ur $e^a$, $a\ne0$ rational oder Liouville-Zahl}, Math.~Ann.~192 (1971), 229--242, doi:10.1007/bf02052875.
\bibitem{CambridgeFF} D.~Bertrand, \emph{Schanuel's conjecture for non-isoconstant elliptic curves over function fields}, in: \emph{Model Theory with Applications to Algebra and Analysis}, Cambridge Univ.~Press, 2008, pp.~41--62, doi:10.1017/cbo9780511735226.004. (Periods counterexample; Riemann--Legendre relations quadratic.)
\bibitem{MO33817} MathOverflow, \emph{Work on independence of $\pi$ and $e$}, question 33817, \url{https://mathoverflow.net/questions/33817}. (Answers verified via API: Schanuel derivation; Zilber program.)
\bibitem{Waldschmidt2025} M.~Waldschmidt, \emph{Schanuel Property for Elliptic and Quasi-Elliptic Functions}, arXiv:2504.14041. Preprint; almost-all result.
\bibitem{DLMF23} NIST DLMF, \S23.2, \emph{Weierstrass Elliptic Functions: Definitions and Periodic Properties}, \url{https://dlmf.nist.gov/23.2}. (Quasi-periodicity verified; exact Legendre constants not extracted.)
\bibitem{Stambul} P.~Stambul, \emph{A generalization of Perron's theorem about Hurwitzian numbers}, Acta Arith.~80 (1997), no.~2, 141--148, doi:10.4064/aa-80-2-141-148.
\bibitem{MatthewsWalters1970} K.~R.~Matthews and R.~F.~C.~Walters, \emph{Some properties of the continued fraction expansion of $(m/n)e^{1/q}$}, Math.~Proc.~Cambridge Philos.~Soc.~67 (1970), no.~1, 67--74, doi:10.1017/S0305004100057108.
\bibitem{Davis1945} C.~S.~Davis, \emph{On some simple continued fractions connected with $e$}, J.~London Math.~Soc.~s1-20 (1945), no.~4, 194--198, doi:10.1112/jlms/s1-20.4.194. (Metadata verified via CrossRef; paywalled full text not opened.)
\bibitem{Davis1978} C.~S.~Davis, \emph{Rational approximations to $e$}, J.~Austral.~Math.~Soc.~25 (1978), no.~4, 497--502, doi:10.1017/S1446788700021480.
\bibitem{Walters1968} R.~F.~C.~Walters, \emph{Alternative derivation of some regular continued fractions}, J.~Austral.~Math.~Soc.~8 (1968), no.~2, 205--212, doi:10.1017/S1446788700005255.
\bibitem{Perron1954} O.~Perron, \emph{Die Lehre von den Kettenbr\"uchen}, 3.~verb.~u.~erw.~Aufl., Band~I: \emph{Elementare Kettenbr\"uche}, B.~G.~Teubner, Stuttgart, 1954. (The digitized table of contents in the German National Library gives p.~119 as the opening page of \S33.)
\bibitem{Komatsu1999JTNB} T.~Komatsu, \emph{On inhomogeneous Diophantine approximation with some quasi-periodic expressions, II}, J.~Th\'eor.~Nombres Bordeaux~11 (1999), no.~2, 331--344, doi:10.5802/jtnb.254.
\bibitem{Komatsu2003Acta} T.~Komatsu, \emph{On Hurwitzian and Tasoev's continued fractions}, Acta Arith.~107 (2003), 161--177, doi:10.4064/aa107-2-4.
\bibitem{Komatsu2003Tsukuba} T.~Komatsu, \emph{Simple continued fraction expansions of some values of certain hypergeometric functions}, Tsukuba J.~Math.~27 (2003), no.~1, 161--173, doi:10.21099/tkbjm/1496164567.
\bibitem{Komatsu2007Czech} T.~Komatsu, \emph{Hurwitz continued fractions with confluent hypergeometric functions}, Czechoslovak Math.~J.~57~(132) (2007), no.~3, 919--932, doi:10.1007/s10587-007-0085-1.
\bibitem{Komatsu2008Sarajevo} T.~Komatsu, \emph{More on Hurwitz and Tasoev continued fractions}, Sarajevo J.~Math.~4~(17) (2008), no.~2, 155--180, doi:10.5644/sjm.04.2.02.
\bibitem{Komatsu2005Monatsh} T.~Komatsu, \emph{Hurwitz and Tasoev continued fractions}, Monatsh.\ Math.~145 (2005), 47--60, doi:10.1007/s00605-004-0281-0.
\bibitem{Komatsu2009JNT} T.~Komatsu, \emph{Shrinking the period length of quasi-periodic continued fractions}, J.~Number Theory~129 (2009), no.~2, 358--366, doi:10.1016/j.jnt.2008.08.004. (Abstract verified; paywalled full text not opened.)
\bibitem{KalininKomatsu2026} N.~Kalinin and T.~Komatsu, \emph{Error terms for continued fractions of $e^{1/s}$ and $\sqrt{v/u}\tanh(1/\sqrt{uv})$}, arXiv:2604.09726 (2026). Preprint.
\bibitem{HBCM2020} C.~Havens, S.~Barbero, U.~Cerruti, and N.~Murru, \emph{Linear fractional transformations and nonlinear leaping convergents of some continued fractions}, Res.~Number Theory~6 (2020), no.~1, art.~11, doi:10.1007/s40993-020-0187-5, arXiv:2002.12644.
\bibitem{Dan2026} Y.~Dan, \emph{Critical convergence and Hausdorff measures for generalized Flint Hills series}, arXiv:2609.12338 (2026). Preprint; abstract verified; studies generalized Flint Hills series.
\bibitem{Conrad} K.~Conrad, \emph{Irrationality of $\pi$ and $e$}, lecture notes, \url{https://kconrad.math.uconn.edu/blurbs/analysis/irrational.pdf}.
\bibitem{Zhu2024} Y.~Zhu, \emph{Transcendental Numbers: Algebraic Independence}, USTC Press, 2024. (Chinese monograph.)
\bibitem{Brownawell1974} W.~D.~Brownawell, \emph{The algebraic independence of certain numbers related by the exponential function}, J.~Number Theory 6 (1974), 22--31. (Source of the $\ee^{\pi^2}$ conditional result.)
\bibitem{Lang1966} S.~Lang, \emph{Introduction to Transcendental Numbers}, Addison-Wesley, 1966; and \emph{Nombres transcendants}, S\'eminaire Bourbaki 1964/66, exp.~305. (Six exponentials theorem.)
\bibitem{Ramachandra1968} K.~Ramachandra, \emph{Contributions to the theory of transcendental numbers (I)}, Acta Arith.~14 (1968), 65--72. (Six exponentials theorem.)
\bibitem{BertolinMotives} C.~Bertolin, \emph{A conjecture in Schanuel style for 1-motives}, arXiv:2509.08700v3 (2026), Prop.~5.4.
\bibitem{Hetyei2014} G.~Hetyei, \emph{Hurwitzian continued fractions containing a repeated constant and an arithmetic progression}, SIAM J.~Discrete Math.~28 (2014), no.~2, 962--985, doi:10.1137/130926092, arXiv:1211.2494.
\bibitem{FischlerRivoal2016} S.~Fischler and T.~Rivoal, \emph{Arithmetic theory of $E$-operators}, J.~\'Ec.~polytech.~Math.~3 (2016), 31--65. (Defines $\mathbf{E}$, $\mathbf{G}$; units contain $\overline{\Q}^*$ and $\exp(\overline{\Q})$; arXiv:1406.5995.)
\bibitem{FischlerRivoal2014} S.~Fischler and T.~Rivoal, \emph{On the values of $G$-functions}, Comment.~Math.~Helv.~89 (2014). ($\pi,1/\pi\in\mathbf{G}$; arXiv:1103.6022.)
\bibitem{FischlerRivoal2024} S.~Fischler and T.~Rivoal, \emph{On the linear independence of values of $G$-functions}, J.~Number Theory 261 (2024). (Conjecture $\mathbf{E}\cap\mathbf{G}=\overline{\Q}$; arXiv:2301.13518.)
\bibitem{FischlerRivoal2026} S.~Fischler and T.~Rivoal, \emph{The ring of values of $E$-functions is not a field}, arXiv:2610.16603 (2026). Preprint, unreviewed; $\mathbf{E}$ not a field, unit conjecture $\alpha e^\beta$.
\bibitem{ParkkonenPaulin2011} J.~Parkkonen and F.~Paulin, \emph{Spiraling spectra of geodesic lines in negatively curved manifolds}, Math.~Z.~268 (2011), no.~1--2, 101--142, doi:10.1007/s00209-010-0662-0; Erratum, Math.~Z.~276 (2014), no.~3--4, 1215--1216, doi:10.1007/s00209-013-1239-5.
\bibitem{BugeaudQuad} Y.~Bugeaud, \emph{On the quadratic Lagrange spectrum}, Math.~Z.~276 (2014), no.~3--4, 985--999, doi:10.1007/s00209-013-1230-1.
\bibitem{Pejkovic2016} T.~Pejkovi\'c, \emph{Quadratic Lagrange spectrum}, Math.~Z.~283 (2016), no.~3--4, 861--869, doi:10.1007/s00209-016-1624-y.
\bibitem{Lin2018} X.~Lin, \emph{Quadratic Lagrange spectrum: I}, Math.~Z.~289 (2018), no.~1--2, 515--533, doi:10.1007/s00209-017-1964-2, arXiv:1703.04847v2.
\bibitem{Hildebrand2024} A.~Jin, S.~Singh, Z.~Zhang and A.~J.~Hildebrand, \emph{On the continued fraction expansion of almost all real numbers}, Involve~18 (2025), 583--600, doi:10.2140/involve.2025.18.583.
\bibitem{Yu2026} R.~Yu, \emph{Tail Criteria, No-Go Audits, and Ap\'ery-Type Certificate Obstructions for the Irrationality of $e+\pi$}, arXiv:2606.17303v1 (2026). Preprint; abstract checked.
\bibitem{OpenAI2026} OpenAI, \emph{The irrationality exponent of $\pi$ is 2}, preprint dated 2026-09-24, repository \texttt{openai/math} (public initial commit 2026-10-06), with Lean formalization \texttt{lean/OAI/NumberTheory/PiExponent/}. The selected external Lean statement concerns the exponent. The paper also derives convergence of the Flint Hills series. The external development was not independently compiled here.

\bibitem{CompanionGauge} \emph{Discrete spectra, exact Hausdorff gauges and dimension phase transitions for Diophantine approximation along the convergent denominators of $e$}, companion manuscript, November 2026, Sect.~6.

\bibitem{ConradConductor}
K.~Conrad, \emph{The conductor ideal of an order}, expository notes,
\url{https://kconrad.math.uconn.edu/blurbs/gradnumthy/conductor.pdf}.
\bibitem{NeukirchANT}
J.~Neukirch, \emph{Algebraic Number Theory}, Grundlehren der
mathematischen Wissenschaften, vol.~322, Springer-Verlag, 1999.
\end{thebibliography}

\clearpage
\appendix

\section{Computational verification and certificates}\label{sec:computation}

\subsection{Symbolic machine verification}

Twelve key algebraic steps were verified with sympy 1.14.0 (script \texttt{verify\_round3.py}); all passed: Vieta's root identity; the six pair solves of Theorem~\ref{thm:four}; the identity $\irr_0=1+1/(M-1)$ of Theorem~\ref{thm:translation}; the termwise recurrence identity $(n+1)\cdot n!/k!=(n+1)!/k!$; the Gauss integral; the mixed integral; the Stirling limit. The cosine integral was verified separately in \texttt{verify\_round4.py}. These verify algebraic identities, not kernel-checked formal proofs; transcendence facts are used as cited theorems.

\subsection{Interval-arithmetic certificate (450 digits)}

Using directed-rounding interval arithmetic (mpmath, 450 decimal digits), $e+\pi$ was enclosed in an interval of width $<9.13\times10^{-451}$. Two certificates were produced (\texttt{certificate\_\allowbreak{}round4.json}):
\begin{enumerate}
\item For every $q\le1000$, the open interval $(Lq,Uq)$ contains no integer; hence $e+\pi\ne p/q$ for all $q\le1000$, rigorously.
\item The window condition of Proposition~\ref{prop:window} is rigorously violated at $n=2,4$ and all $7\le n\le60$ except $n=34$; every $q\le60$ has a certified first violation (for $q=34$ it is $n=35$).
\end{enumerate}
These are finite-range interval certificates, conditional on the stated implementation and directed-rounding computations, not progress on Problem~\ref{prob:main}; their value is methodological --- necessary conditions converted into reproducible certificates.

\begin{proposition}[Certified separation on a denominator range]\label{prop:farey-gap}
Set $s=e+\pi$. For all integers $a,b$ with $1\le b<3635695100$,
\[
  |bs-a|>1/5000000000 .
\]
Consequently $D_j(b)\ne0$ for all such $b$ and all $j\ge8$.
\end{proposition}
\begin{proof}
Write $(p,q)=(16687677659,2847787561)$ and
$(u,v)=(4617039282,787907539)$. Direct multiplication gives $uq-pv=1$.
The rational bounds for $\pi$ at 20 decimal places in mathlib and
$|e-363916618873/133877442384|\le10^{-20}$ give
$q s-p>2\cdot10^{-10}$ and $u-vs>2\cdot10^{-10}$ by rational arithmetic.
Set $m=ub-va$, $n=qa-pb$. Then $(a,b)=m(p,q)+n(u,v)$.
Since $0<b<q+v$, exactly one coefficient is positive and the other is
nonpositive. In $bs-a=m(qs-p)-n(u-vs)$ there is consequently no
cancellation. For $j\ge8$ one has $Q_{j+1}\ge Q_9=10622799089$ and
$|\delta_{j+1}|<1/(2(j+1)Q_{j+1})$. Hence a match would give, by
Proposition~\ref{prop:finite-match},
\[
 \|bs\|\le\frac{1/2+b|\delta_{j+1}|}{Q_{j+1}}
 \le\frac{1/2+3635695100/Q_9}{Q_9}<2\cdot10^{-10}.
\]
\end{proof}
The four declarations in the accompanying Lean certificate were rechecked against the local pinned project. Their axiom reports contain only the three Lean base axioms listed in Appendix~\ref{sec:lean}. The denominator range is finite even though the range of indices is
infinite. The Lean source is \texttt{DenominatorCertificate.lean}.

\begin{proposition}[Consequence of a rational affine relation]\label{prop:nest-transfer}
If $\pi=r+se$ with $r,s\in\Q$ and $s\ne0$, then
$e,\exp(e),\Gamma(1/4)$ are algebraically independent over $\Q$.
This includes the hypothesis $e+\pi\in\Q$.
\end{proposition}
\begin{proof}
Put $G=\Gamma(1/4)$ and $K=\Q(e,\exp(e),G)$. Choose $d\ge1$ with
$dr,ds\in\Z$. Then $\pi\in K$ and
$(\exp\pi)^d=e^{dr}\exp(e)^{ds}\in K$, including for negative exponents.
The algebraic extension $K(\exp\pi)$ contains
$\pi,\exp(\pi),G$, which are algebraically independent by
Nesterenko's theorem \cite{Nesterenko1996}. Consequently $K$ has
transcendence degree at least three. Its three generators are therefore
algebraically independent.
\end{proof}
This consequence uses a proved property specific to $\pi$. It supplies a
contradiction only if an independent algebraic dependence of the resulting
triple is proved. It is an application of Nesterenko's theorem.

\subsection{Continued-fraction experiment}

For $q=1,\dots,12$, with $p$ the nearest integer to $(e+\pi)q$, the continued fraction of $p/q-e$ was compared with that of $\pi$ at 200-digit precision, using
\[
\pi=[3;7,15,1,292,1,1,1,2,1,3,1,14,2,\dots].
\] All candidates diverge at partial-quotient position $1$ or $2$ ($q=7$, the best approximation among these, matches two positions). This visualizes the mechanism of Theorem~\ref{thm:cf}: approximation quality determines match length; the finitely many tested candidates disagree with $\pi$ at an identifiable step (numerically); this does not quantify over all rational candidates.

\subsection{Flint Hills numerics}

The partial sum is $S(10000)=30.31451083$ (numerical evidence, not proof). \textbf{Correction (2026-11-14):} an earlier version of this section stated that $S(10000)$ ``agrees with the literature's claimed limit $30.3144\dots$''. That comparison is withdrawn: all terms of the series are positive, so if the series converges to a finite limit $L$, then necessarily $L>S(10000)=30.31451083$; the quoted value $30.3144\dots$ is \emph{smaller} than the partial sum and therefore cannot be the limit. The largest term occurs at $n=355$ --- the numerator of the famous convergent $355/113$ of $\pi$ --- with value $24.6$ and $\sin(355)\approx-3.014\times10^{-5}$. The series' behavior is governed entirely by the Diophantine approximation properties of $\pi$, corroborating the logic of Theorem~\ref{thm:flint}. Dan \cite{Dan2026} studies critical convergence for generalized Flint Hills series. The classical convergence claim and the scope of its external verification are discussed in \S\ref{sec:lean-external}.

\subsection{Chowla--Selberg exploration (closed)}

The identity $\varpi=\Gamma(1/4)^2/(2\sqrt{2\pi})$ (lemniscatic constant) yields $\pi=\Gamma(1/4)^4/(8\varpi^2)$, verified at 120 digits (error $7.7\times10^{-121}$). Substituting into the subquestion $\pi=a\cdot\wp_\Omega(p)$ gives $\Gamma(1/4)^4=8a\varpi^2\wp_\Omega(p)$; but $\varpi^2=\Gamma(1/4)^4/(8\pi)$ is the same identity, so the substitution is circular and produces no new constraint. Nesterenko's theorem does not cover $\varpi$ or $\wp(p)$. \emph{Verdict: closed, with precise reason.} (The identity's literature source was not independently verified; it is recorded as numerically verified only.)

\section{Machine-Verified Formalization in Lean}\label{sec:lean}

This section records the formalization of the paper's derivations in the Lean~4 proof assistant (Lean 4.35.0 with mathlib v4.35.0). We distinguish three trust levels throughout: (i) \emph{kernel-checked} results, verified by the Lean kernel and audited with \texttt{\#print axioms}; (ii) \emph{human-verified preprint proof structures}, which were originally admitted as explicit axioms and flagged as such --- the Meiburg axiom has since been eliminated by a complete machine proof; (iii) \emph{symbolic/numeric auxiliary checks}, which are not proofs. The formalization verifies the \emph{logical correctness of the derivations}, not the open Problem~\ref{prob:main}.

\subsection{The irrationality measure: a non-degenerate definition}\label{sec:lean-def}

The skeleton supplied with the research archive contained a degenerate definition: allowing rational numerators gives infinitely many approximations even at $q=1$, because infinitely many rationals lie in $(x-1,x+1)$. Thus every nonnegative exponent is realized and the extended-real supremum becomes $+\infty$ for every $x$; this does not measure Diophantine approximation. We corrected this in two ways, both essential.

\begin{definition}[Approximation pairs and measure]\label{def:lean-measure}
For $x\in\R$ and $\nu\in\R$, the set of \emph{$\nu$-approximation pairs} is
\[
A(x,\nu)=\{(p,q)\in\Z\times\N: q>0,\ |x-p/q|<q^{-\nu}\}.
\]
The \emph{approximation-exponent set} is $S(x)=\{\nu\in\R_{\ge0}: A(x,\nu)\text{ is infinite}\}$, and the \emph{irrationality measure} is
\[
\irr(x)=\sup\,\iota(S(x))\ \in\ [0,\infty],
\]
where $\iota:\R_{\ge0}\hookrightarrow[0,\infty]$ is the embedding into the extended nonnegative reals. For irrational $x$ this agrees with Definition~\ref{def:measure}. For rational $x=a/b$, the exact pairs $(ka,kb)$ make every exponent admissible, so this formal extension has value $+\infty$. No assertion that it equals $1$ on rationals is used.
\end{definition}

Two design choices are forced. The numerator must range over $\Z$ (not $\Q$) to avoid the degeneracy above. The value must live in $[0,\infty]$ (Lean's \texttt{ENNReal}), not $\R$: a Liouville number has $\irr(x)=\infty$, and the supremum of an unbounded set in $\R$ does not exist. With \texttt{ENNReal}, $\sup$ is always defined, and the comparison lemma below avoids all boundedness side-conditions.

\begin{lemma}[$\varepsilon$-absorption comparison]\label{lem:eps-abs}
If for every $\nu\in S(y)$ and every $0<\varepsilon\le\nu$ one has $\nu-\varepsilon\in S(x)$, then $\irr(y)\le\irr(x)$.
\end{lemma}
\begin{proof}
In \texttt{ENNReal} this is \texttt{le\_\allowbreak{}of\_\allowbreak{}forall\_\allowbreak{}pos\_\allowbreak{}le\_\allowbreak{}add}: when $\irr(x)=\infty$ the claim is trivial; otherwise approximate $\nu$ from below by $\nu-\varepsilon'$ with $\varepsilon'=\tfrac12\min(\varepsilon,\nu)$, use $\nu-\varepsilon'\in S(x)$ to bound $\nu\le\irr(x)+\varepsilon$, and let $\varepsilon\to0$. No nonemptiness or boundedness hypothesis is needed.
\end{proof}

The invariance lemmas (translation, scaling, reciprocal) and Theorem~\ref{thm:measure-reduction} are all kernel-checked; their proofs use a denominator-unboundedness lemma (fixed $q$ forces $p$ into a finite interval, and vice versa) together with constant absorption $C\,Q^{-\nu}\le Q^{-(\nu-\varepsilon)}$ on the cofinite subset $Q>T$. We do not repeat them here; the audit is summarized in \S\ref{sec:lean-status}.

\subsection{Status of the formalization}\label{sec:lean-status}

The earlier 9098-job, 9099-job, and 9101-job baselines were built and
reproduced separately. In the current fully integrated sources
(2026-11-22), \texttt{lake build EpiFormal} completed cleanly
(9179 jobs, exit code~0) and \texttt{lake env lean Audit.lean}
(recorded in \texttt{verification/\allowbreak audit\_\allowbreak final\_\allowbreak 20261122.log}
and reconciled by \texttt{verification/\allowbreak check\_\allowbreak audit\_\allowbreak reconcile.py})
reports 937 audited declarations across 134 scanned Lean files with zero \texttt{sorryAx} and zero
project-added axioms. Of these, 914 use only the Lean base axioms
(\texttt{propext}, \texttt{Classical.choice}, \texttt{Quot.sound},
including 5 axiom-free declarations) and 23 also use the three
whitelisted classical inputs (\texttt{e\_transcendental},
\texttt{pi\_transcendental}, \texttt{mu\_e\_eq\_two}).
In particular, every Lean module referenced in this paper---including
\texttt{EpiFormal/\allowbreak GeneralEndpoint\allowbreak Rigidity.lean}
(30 base-only declarations),
\texttt{EpiFormal/\allowbreak TrichotomyAnd\allowbreak PolyExtension.lean}
(45 base-only declarations),
\texttt{EpiFormal/\allowbreak AffineWindow\allowbreak Assembly.lean}
(7 base-only declarations),
\texttt{EpiFormal/\allowbreak Denominator\allowbreak Certificate.lean},
and the eight previously standalone helper modules
\texttt{GoldenHeadCutoff.lean}, \texttt{GoldenAffineEndpoint.lean},
\texttt{AffineWindowAudit.lean}, \texttt{DigitTailBound.lean},
\texttt{DensityCube.lean}, \texttt{OrbitAndFlow.lean},
\texttt{ResearchAdvance.lean}, and \texttt{TraceMarkov.lean}
(29 base-only declarations combined)---is integrated into
\texttt{EpiFormal.lean} and audited inside \texttt{Audit.lean}.

Some declarations use whitelisted classical inputs such
as \texttt{pi\_transcendental} or \texttt{mu\_e\_eq\_two}.
The Meiburg criterion, formerly an axiom, is now a theorem (Section~\ref{sec:lean-f2-remaining}). In particular the Flint Hills bridge (Theorem~\ref{thm:flint}) is kernel-checked with dependencies exactly $\{\texttt{mu\_\allowbreak{}e\_\allowbreak{}eq\_\allowbreak{}two},\texttt{pi\_\allowbreak{}transcendental}\}$. The Alekseyev direction (convergence $\Rightarrow\irr(\pi)\le5/2$) is now machine-verified via a nearest-integer argument (Section~\ref{sec:alekseyev}, theorem \texttt{alekseyev\_\allowbreak{}bound}, depending only on $\{\texttt{pi\_\allowbreak{}transcendental}\}$ plus base axioms); it is not needed for the bridge. For these principal bridge results, the remaining non-base assumptions are \texttt{mu\_\allowbreak{}e\_\allowbreak{}eq\_\allowbreak{}two} ($\irr(\ee)=2$, Bundschuh) and the transcendence of $\ee$ and $\pi$ (Hermite--Lindemann).

\subsection{Formalizing Meiburg's criterion: the sparsity algebra}\label{sec:lean-f2}

We describe the formalization that replaces the axiom \texttt{meiburg\_criterion} by a theorem. The key simplification is that the bridge needs only the special case $\irr(\pi)=2$, not the full statement for arbitrary $\irr(\alpha)<1+u/v$. This permits a direct dyadic-block counting argument rather than Meiburg's full fine partition. The development contains two \emph{separate} formalizations. The first is a sparsity algebra (good approximations cannot cluster), recorded in \texttt{EpiFormal/\allowbreak{}Meiburg.lean}; it is a self-contained verification of Meiburg's sparsity lemma and is \emph{not} used by the convergence proof. The second, in \texttt{EpiFormal/\allowbreak{}MeiburgConvergence.lean}, is a direct dyadic-block counting plus geometric-summation argument that proves the Flint Hills case without invoking the sparsity layer. Both are kernel-checked (no \texttt{sorry}); we record them separately to avoid implying a dependency that does not exist.

\subsubsection{From the measure bound to finiteness and eventual lower bounds}

\begin{lemma}[Good approximations are finite above the measure]\label{lem:f2-finite}
If $\irr(x)\le M$ (with $M\in\R_{\ge0}$) and $\nu>M$, then $A(x,\nu)$ is finite.
\end{lemma}
\begin{proof}
Suppose $A(x,\nu)$ were infinite. Since $\nu>M\ge0$, we have $\nu>0$, so $\nu$ defines an element of $S(x)$. By definition of $\irr$ as a supremum, $\nu\le\irr(x)\le M$, contradicting $\nu>M$. (In Lean: \texttt{by\_contra}, embed $\nu$ into \texttt{NNReal}, use \texttt{le\_sSup}, and cast the \texttt{ENNReal} inequality back to $\R$ via \texttt{coe\_le\_coe}.)
\end{proof}

\begin{lemma}[Eventual approximation lower bound, denominator form]\label{lem:f2-eventual}
If $\irr(x)\le M$ and $\varepsilon>0$, then there exists $N$ such that for all $n\ge N$ and all $m\in\Z$,
\[
|x-m/n|\ge n^{-(M+\varepsilon)}.
\]
\end{lemma}
\begin{proof}
Apply Lemma~\ref{lem:f2-finite} with $\nu=M+\varepsilon$: the set $A(x,M+\varepsilon)$ is finite, so its projection to denominators is a finite set $S$ of naturals. Take $N=\max S+1$ (or $N=1$ if $S=\varnothing$). If some $n\ge N$ violated the bound, then $(m,n)\in A(x,M+\varepsilon)$, so $n\in S$ and $n\le\max S<N$, a contradiction.
\end{proof}

\begin{remark}[Why the denominator form matters]\label{rem:f2-denom}
Lemma~\ref{lem:f2-eventual} uses $n$ as the \emph{denominator}. This is the direction that follows directly from the definition of $\irr$. The reverse direction needed for the Alekseyev bound (Theorem~\ref{thm:flint}, necessary side) involves $|\sin n|\le|n-m\pi|=|m|\,|\pi-n/m|$, where $n$ is the \emph{numerator}; conflating the two is the algebraic error $|n-m\pi|\ne n|\pi-m/n|$ that blocked an earlier attempt. The present lemmas deliberately stay on the measure-upper-bound side, where no such swap is needed.
\end{remark}

\subsubsection{The difference-quotient error bound}

The sparsity lemma rests on a simple algebraic observation: two close approximations produce a new approximation whose error is controlled by the gap of the denominators.

\begin{lemma}[Difference-quotient error bound]\label{lem:f2-dq}
Let $q_1<q_2$ be positive integers, $p_1,p_2\in\Z$, and suppose $|x-p_i/q_i|\le A_i$ for $i=1,2$. Then
\[
\left|x-\frac{p_2-p_1}{q_2-q_1}\right|\le\frac{q_1 A_1+q_2 A_2}{q_2-q_1}.
\]
\end{lemma}
\begin{proof}
The key identity is
\[
x-\frac{p_2-p_1}{q_2-q_1}=\frac{q_2(x-p_2/q_2)-q_1(x-p_1/q_1)}{q_2-q_1},
\]
verified by clearing denominators (in Lean: \texttt{field\_simp}, \texttt{push\_cast}, \texttt{ring}). Taking absolute values and applying the triangle inequality,
\[
|q_2(x-p_2/q_2)-q_1(x-p_1/q_1)|\le q_2|x-p_2/q_2|+q_1|x-p_1/q_1|\le q_2 A_2+q_1 A_1,
\]
where we used $q_i>0$ to pull the factors out of the absolute values. Dividing by $q_2-q_1>0$ gives the claim.
\end{proof}

Specializing to good approximations yields a strict bound.

\begin{lemma}[Good-approximation difference quotient]\label{lem:f2-good-dq}
If $|x-p_i/q_i|<q_i^{-\nu}$ for $i=1,2$ and $q_1<q_2$, then
\[
\left|x-\frac{p_2-p_1}{q_2-q_1}\right|<\frac{q_1^{1-\nu}+q_2^{1-\nu}}{q_2-q_1}.
\]
\end{lemma}
\begin{proof}
Apply Lemma~\ref{lem:f2-dq} with $A_i=q_i^{-\nu}$ and use $q_i\cdot q_i^{-\nu}=q_i^{1-\nu}$ (the law $q^{1-\nu}=q^1\cdot q^{-\nu}$, i.e.\ \texttt{rpow\_add}). The strict inequality is preserved because both input bounds are strict and the operations are monotone.
\end{proof}

\subsubsection{The exponent estimate and the sparsity core}

The next lemma is the real-analytic heart: under a numerical condition, the difference quotient is an approximation of \emph{strictly higher} exponent.

\begin{lemma}[Sparsity exponent estimate]\label{lem:f2-exp}
Let $\nu>1$, $\delta>0$, $1\le q_1<q_2\le2q_1$, and $Q=q_2-q_1>0$. If
\begin{equation}\label{eq:f2-cond}
2\,q_1^{1-\nu}<Q^{1-\nu-\delta},
\end{equation}
then
\[
\frac{q_1^{1-\nu}+q_2^{1-\nu}}{Q}<Q^{-\nu-\delta}.
\]
\end{lemma}
\begin{proof}
Since $\nu>1$, the exponent $1-\nu$ is negative; the function $t\mapsto t^{1-\nu}$ is antitone on $(0,\infty)$, so $q_1\le q_2$ gives $q_2^{1-\nu}\le q_1^{1-\nu}$. Hence the numerator is at most $2q_1^{1-\nu}$. For the denominator, $Q^{1-\nu-\delta}=Q^{-\nu-\delta}\cdot Q$ (again \texttt{rpow\_add}). Thus
\[
\frac{q_1^{1-\nu}+q_2^{1-\nu}}{Q}\le\frac{2q_1^{1-\nu}}{Q}<\frac{Q^{1-\nu-\delta}}{Q}=Q^{-\nu-\delta},
\]
where the strict step is exactly condition~\eqref{eq:f2-cond} divided by $Q>0$.
\end{proof}

Combining the two previous lemmas gives the sparsity core.

\begin{lemma}[Sparsity core: the difference quotient is better]\label{lem:f2-core}
Let $\nu>1$, $\delta>0$, $2\le q_1<q_2\le2q_1$, and suppose $(p_i,q_i)$ are $\nu$-good approximations of $x$ (i.e.\ $|x-p_i/q_i|<q_i^{-\nu}$). If $Q=q_2-q_1$ satisfies~\eqref{eq:f2-cond}, then
\[
\left|x-\frac{p_2-p_1}{Q}\right|<Q^{-(\nu+\delta)}.
\]
That is, the difference quotient is a $(\nu+\delta)$-good approximation.
\end{lemma}
\begin{proof}
Lemma~\ref{lem:f2-good-dq} bounds the left side by $(q_1^{1-\nu}+q_2^{1-\nu})/Q$; Lemma~\ref{lem:f2-exp} bounds that by $Q^{-\nu-\delta}=Q^{-(\nu+\delta)}$. Chaining the two strict inequalities gives the result. (In Lean the only subtlety is that the denominator is written both as $(q_2-q_1:\N)$ cast to $\R$ and as $(q_2:\R)-(q_1:\R)$; these are definitionally equal via \texttt{Nat.cast\_sub}, and \texttt{simpa} reconciles them.)
\end{proof}

The meaning of Lemma~\ref{lem:f2-core} is exactly Meiburg's sparsity principle: if two $\nu$-good approximations have $\delta$-close denominators (so $Q$ is small relative to $q_1$) and $q_1$ is large enough for~\eqref{eq:f2-cond} to hold, their difference quotient is a $(\nu+\delta)$-good approximation. When $\nu$ is chosen just below $\irr(x)$, the exponent $\nu+\delta$ exceeds $\irr(x)$, so by Lemma~\ref{lem:f2-finite} only finitely many such difference quotients exist. This forces the good denominators to be sparse.

\subsubsection{The multiplicity bound (fiber finiteness)}

To turn ``finitely many difference quotients'' into a density bound on denominators, one must control how many pairs $(p_1,q_1),(p_2,q_2)$ can yield the \emph{same} difference $(p_2-p_1,q_2-q_1)$. The first ingredient is that a fixed denominator admits only finitely many good numerators.

\begin{lemma}[Fixed-denominator fiber is finite]\label{lem:f2-fiber}
For any $d\in\N_{>0}$, any $\nu\in\R$, and any $x\in\R$, the set
\[
\{p\in\Z:|x-p/d|<d^{-\nu}\}
\]
is finite.
\end{lemma}
\begin{proof}
Multiplying by $d>0$, the condition is $|dx-p|<d^{1-\nu}$, i.e.
\[
dx-d^{1-\nu}<p<dx+d^{1-\nu}.
\]
Thus $p$ lies in an open interval of length $2d^{1-\nu}$, hence in the closed integer interval $[\lfloor L\rfloor,\lceil U\rceil]$ with $L=dx-d^{1-\nu}$, $U=dx+d^{1-\nu}$, which is finite. (In Lean: \texttt{Set.finite\_Icc}, with the floor/ceil bounds from \texttt{Int.floor\_le} and \texttt{Int.le\_ceil}.)
\end{proof}

Lemma~\ref{lem:f2-fiber} is the foundation of Meiburg's $Q_{\max}$ multiplicity bound: a fixed difference denominator $Q$ admits only finitely many difference numerators, so each difference quotient arises from a bounded number of pairs once the goodness and closeness constraints are imposed.

\subsection{Completion of the convergence argument}\label{sec:lean-f2-remaining}

The seven lemmas above (Lemmas~\ref{lem:f2-finite}--\ref{lem:f2-fiber}) form the algebraic core of the \emph{separate} sparsity formalization; they are not invoked by the convergence proof. The convergence itself is completed by the direct dyadic counting and geometric summation below, which replaces the Meiburg axiom by a theorem. Both developments are kernel-checked with no \texttt{sorry}. We record the completed steps, preserving the reasoning.

\begin{lemma}[Circle spacing]\label{lem:f2-spacing}
If $\irr(\pi)\le M<s$ with $s>1$, there is $c>0$ such that $\|h/\pi\|\ge c\,h^{1-s}$ for every positive integer $h$, where $\|x\|$ is the distance from $x$ to the nearest integer.
\end{lemma}
\begin{proof}
By Lemma~\ref{lem:f2-finite} the set of $s$-good approximations is finite; taking the minimum of $b^s|\pi-a/b|$ over these exceptions (positive by irrationality of $\pi$) and of $1$ gives a uniform constant $C>0$ with $|\pi-a/b|\ge C b^{-s}$ for all $a,b$. Let $j$ be the nearest integer to $h/\pi$. If $j\ge1$, then $\|h/\pi\|=(j/\pi)|\pi-h/j|\ge(C/\pi)j^{1-s}\ge(C/\pi)h^{1-s}$ since $j\le h$ and $1-s<0$. If $j=0$, then $\|h/\pi\|=h/\pi\ge h^{1-s}/\pi$. Take $c=\min(C,1)/\pi$. The two cases (in particular $j=0$) are handled separately in the formalization.
\end{proof}

\begin{lemma}[Dyadic-block counting]\label{lem:f2-count}
With $\delta=c(2N)^{1-s}$, the number of $n\in[N,2N)$ with $\|n/\pi\|\le t$ is at most $2(\lfloor t/\delta\rfloor+1)$.
\end{lemma}
\begin{proof}
Each point is at distance $\ge\delta$ from $0$ (Lemma~\ref{lem:f2-spacing}), and two distinct points are at circular distance $\ge\delta$ since $1\le|n-m|<2N$. Representing points in $[-1/2,1/2)$ and sorting each side of $0$, the $i$-th point is at distance $\ge i\delta$. The formalization uses a floor-injection $x\mapsto\lfloor x/\delta\rfloor$ (injective because two points in the same length-$\delta$ interval would be $<\delta$ apart) rather than an explicit sort, bounding each side by $\lfloor t/\delta\rfloor+1$.
\end{proof}

\begin{lemma}[Jordan bound and block sum]\label{lem:f2-blocksum}
$|\sin n|\ge2\|n/\pi\|$, and $\sum_{N\le n<2N}\frac1{n^3\sin^2 n}\le C_s N^{2s-5}$ for a constant $C_s$ independent of $N$.
\end{lemma}
\begin{proof}
Writing $n=\pi y+j\pi$ with $|y|=\|n/\pi\|\le1/2$, periodicity gives $|\sin n|=|\sin(\pi y)|\ge(2/\pi)|\pi y|=2\|n/\pi\|$ by Jordan's inequality. Combining with the counting lemma and the bound $\sum_{x\in S}x^{-2}\le(2/\delta^2)\sum_{i\ge1}(i+1)^{-2}$ for a $\delta$-separated set $S$ (proved by reindexing along the injection $x\mapsto\lfloor x/\delta\rfloor-1$), the block sum is $\le C_s N^{2s-5}$.
\end{proof}

\begin{theorem}[Convergence under the exponent threshold]\label{thm:f2-conv}
If $\irr(\pi)<5/2$, the Flint Hills series converges.
\end{theorem}
\begin{proof}
Fix $s$ with $\irr(\pi)<s<5/2$. The partial sums are monotone (nonnegative terms). Partitioning $[1,\infty)$ into dyadic blocks $[2^k,2^{k+1})$ and applying Lemma~\ref{lem:f2-blocksum}, the partial sums are bounded by $\sum_k C_s 2^{k(2s-5)}$, a convergent geometric series since $2s-5<0$. Monotone bounded partial sums converge.
\end{proof}

\paragraph{Wiring and axiom audit.} The declaration \texttt{meiburg\_\allowbreak{}criterion} is now a \emph{theorem}, proved by Theorem~\ref{thm:f2-conv}. The axiom audit (\texttt{\#print axioms}) confirms:
\begin{itemize}
\item \texttt{meiburg\_\allowbreak{}criterion}, \texttt{flint\_\allowbreak{}hills\_\allowbreak{}converges\_\allowbreak{}of\_\allowbreak{}mu\_\allowbreak{}lt}, \texttt{circle\_\allowbreak{}spacing}, \texttt{dyadic\_\allowbreak{}block\_\allowbreak{}count}, \texttt{block\_\allowbreak{}sum\_\allowbreak{}bound}: depend only on \texttt{pi\_transcendental} and Lean's base axioms, with no \texttt{sorryAx}.
\item \texttt{proposition\_\allowbreak{}H'}: depends on \texttt{mu\_e\_eq\_two}, \texttt{pi\_transcendental} and Lean's base axioms. \textbf{The Meiburg preprint axiom has been eliminated.}
\end{itemize}
At that stage the project built with zero errors (9056 jobs); the later audit scope is reported in \S\ref{sec:lean-status}. The audited declarations for this argument have no \texttt{sorryAx}. The Alekseyev direction (Section~\ref{sec:alekseyev}) is now machine-verified as well (theorem \texttt{alekseyev\_\allowbreak{}bound}); the earlier density obstruction is preserved there as a historical record.

\subsection{The Alekseyev direction: from a recorded obstruction to a machine-verified bound}\label{sec:alekseyev}

The converse implication (Flint Hills converges $\Rightarrow\irr(\pi)\le5/2$, Alekseyev's corollary) is \emph{not} needed for the bridge $H'$. An earlier version of this section recorded it as a research-level gap: the naive termwise argument fails (the precise obstruction is preserved below as a permanent record), and the density argument of Alekseyev's original proof had not been formalized. \textbf{Update (2026-11-14):} a revised supplement supplied a correct paper proof via nearest integers that bypasses the obstruction entirely, and we have now formalized it in full. The bound $\irr(\pi)\le5/2$ is therefore \emph{machine-verified} (file \texttt{EpiFormal/\allowbreak{}Alekseyev.lean}, theorem \texttt{alekseyev\_\allowbreak{}bound}, no \texttt{sorryAx}), conditional on the convergence of the Flint Hills series, which remains an open problem. Three auxiliary lemmas were already machine-verified and are retained:

\begin{lemma}[nonvanishing of $\sin n$] For $n\ge1$, $\sin n\ne0$.
\end{lemma}
\begin{proof} If $\sin n=0$ then $n=m\pi$ for some integer $m\ne0$, so $\pi=n/m\in\Q$, contradicting the transcendence of $\pi$.
\end{proof}

\begin{lemma}[eventual sine lower bound]\label{lem:sin-lower} If $a_n:=1/(n^3\sin^2 n)\to0$, then for every $\varepsilon>0$ there is $N$ such that $|\sin n|\ge n^{-3/2-\varepsilon}$ for all $n\ge N$.
\end{lemma}
\begin{proof} Since $a_n\to0$, eventually $a_n<1$, i.e.\ $n^3\sin^2 n>1$, so $|\sin n|>n^{-3/2}\ge n^{-3/2-\varepsilon}$ for $n\ge1$. The step $|\sin n|>n^{-3/2}$ from $\sin^2 n>n^{-3}$ uses $|x|=\sqrt{x^2}$ and monotonicity of $t\mapsto t^{1/2}$; the final inequality uses $n\ge1$ and $-\tfrac32\ge-\tfrac32-\varepsilon$.
\end{proof}

\begin{lemma}[factorization] For $m\ne0$, $|n-m\pi|=|m|\,|\pi-n/m|$.
\end{lemma}
\begin{proof} $n-m\pi=-m(\pi-n/m)$; take absolute values.
\end{proof}

\paragraph{The obstruction (historical record, preserved).} Combining Lemma~\ref{lem:sin-lower} with the periodicity bound $|\sin n|\le|n-m\pi|$ (formalized as \texttt{abs\_\allowbreak{}sin\_\allowbreak{}int\_\allowbreak{}le}) and the factorization gives
\[
|\pi-n/m|\ \ge\ \frac{n^{-3/2-\varepsilon}}{|m|}.
\]
To deduce the measure bound $\irr(\pi)\le5/2$ one would need $|\pi-n/m|\ge C\,|m|^{-5/2-\varepsilon}$, which is equivalent to $n^{-3/2-\varepsilon}\ge C\,|m|^{-3/2-\varepsilon}$, i.e.\ $n\le C'|m|$ (the exponent is negative). But the lower bound of Lemma~\ref{lem:sin-lower} holds for \emph{all} large $n$, including $n>\pi|m|$, where the inequality reverses. Hence the termwise argument ``$\sin$ lower bound $\Rightarrow\irr(\pi)\le5/2$'' is \textbf{invalid as stated}. This is the precise form of the numerator--denominator obstruction noted in the seventh research round. An earlier version of this section concluded that formalizing Alekseyev's density argument was a genuine research-level gap; that conclusion is now superseded by the nearest-integer argument below, which needs no density control at all. The failed route is kept on record.

\paragraph{The nearest-integer argument (machine-verified).} The revised supplement (2026-11-14) closes the gap on paper by restricting the sine lower bound to the single nearest integer $p=\operatorname{round}(q\pi)$ for each denominator $q$, which automatically satisfies $p\asymp q$. We formalized the argument in full.

\begin{lemma}[nearest-integer facts]\label{lem:nearest} For $q\ge1$, put $p=\operatorname{round}(q\pi)$. Then $|p-q\pi|\le\tfrac12$, $p\le4q$, and $p\to\infty$ as $q\to\infty$.
\end{lemma}
\begin{proof} $|p-q\pi|\le\tfrac12$ is the defining property of rounding. Since $\pi<4$, $p\le q\pi+\tfrac12<4q+\tfrac12<4q+1$, and integrality gives $p\le4q$. Finally $p>q\pi-\tfrac12\to\infty$.
\end{proof}

\begin{lemma}[uniform divergence]\label{lem:uniform} If the Flint Hills series converges, then for every $B>0$ there is $N$ such that $n^3\sin^2n>B$ for all $n\ge N$; consequently $p^{3/2}|\sin p|\to\infty$ along $p=\operatorname{round}(q\pi)$.
\end{lemma}
\begin{proof} Convergence gives $a_n\to0$, so eventually $a_n<1/B$, i.e.\ $n^3\sin^2n>B$. Taking square roots, $n^{3/2}|\sin n|>\sqrt B$ eventually; since $p\to\infty$ (Lemma~\ref{lem:nearest}), the claim follows.
\end{proof}

\begin{theorem}[Alekseyev direction; machine-verified]\label{thm:alekseyev-bound} If the Flint Hills series converges, then $\irr(\pi)\le5/2$.
\end{theorem}
\begin{proof} Fix $\varepsilon>0$. By Lemma~\ref{lem:uniform} with $B=4^{3/2}\cdot4^{3/2}=64$, for all large $q$, writing $p=\operatorname{round}(q\pi)$,
\[
p^{3/2}|\sin p|>4^{3/2}.
\]
For \emph{every} integer $a$, the nearest-integer property and the periodicity bound give $|q\pi-a|\ge|q\pi-p|\ge|\sin p|$, while $p\le4q$ gives $q^{3/2}\ge4^{-3/2}p^{3/2}$. Hence, uniformly in $a$,
\[
q^{5/2}\,|\pi-a/q|=q^{3/2}|q\pi-a|\ \ge\ q^{3/2}|q\pi-p|\ \ge\ q^{3/2}|\sin p|\ \ge\ 4^{-3/2}p^{3/2}|\sin p|\ >\ 1.
\]
Thus eventually $|\pi-a/q|>q^{-5/2}\ge q^{-(5/2+\varepsilon)}$ for all integers $a$ (using $q\ge1$). For any $\nu>5/2$, apply this with $\varepsilon=\nu-5/2$: the set of $\nu$-good approximation pairs is then contained in finitely many denominators, hence finite (each fixed denominator admits only finitely many good numerators). Therefore no $\nu>5/2$ lies in the approximation-exponent set of $\pi$, and $\irr(\pi)\le5/2$.
\end{proof}

The formalization adds six declarations to \texttt{EpiFormal/\allowbreak{}Alekseyev.lean}: \texttt{nearest\_\allowbreak{}int\_\allowbreak{}le\_\allowbreak{}four\_\allowbreak{}mul}, \texttt{nearest\_\allowbreak{}int\_\allowbreak{}tendsto\_\allowbreak{}atTop}, \texttt{term\_\allowbreak{}tendsto\_\allowbreak{}eventually\_\allowbreak{}gt}, \texttt{alekseyev\_\allowbreak{}eventual\_\allowbreak{}lower}, \texttt{apprPairs\_\allowbreak{}finite\_\allowbreak{}of\_\allowbreak{}flintHills}, and \texttt{alekseyev\_\allowbreak{}bound}. The axiom audit confirms that all six carry no \texttt{sorryAx}; \texttt{alekseyev\_\allowbreak{}bound} depends only on \texttt{pi\_transcendental} and Lean's base axioms. The key point is that the lower bound is applied only at $p=\operatorname{round}(q\pi)$, where $p\asymp q$ holds automatically, so replacing $q^{3/2}$ by $p^{3/2}$ costs only the constant $4^{-3/2}$; no density argument is needed.

\paragraph{Honest status.} This formalization is pure engineering over already-established mathematics: it removes a preprint axiom and raises the trust level of the bridge, but adds no new information about Problem~\ref{prob:main}. Theorem~\ref{thm:alekseyev-bound} is conditional on the convergence of the Flint Hills series, which is itself open. The bridge remains the conditional statement ``$e+\pi\in\Q\Rightarrow$ Flint Hills converges''; convergence of the series is not itself a contradiction, so the irrationality of $e+\pi$ remains open.

\subsection{The revised supplement S1--S7 and Propositions 3--4: full joint compilation}\label{sec:lean-m6}

A revised supplement (2026-11-14) supplied paper proofs of a family of structural results around the selected continued-fraction convergents of $\ee$: the classification lattice (S1, Theorem~2), the quantitative gap (S2, Corollary~3), the finite spectrum (S3 Lemma~5 / S4 Theorem~6), signed clusters and realization (S5, Theorems~7--8), defect-tail equivalence and bijection (S6, Theorem~9 / Corollary~10), the two-index criterion (S7, Theorem~11), and the signed and factorial criteria (Propositions~3--4). These were developed as nine concurrent subtasks (A--I) and then independently re-compiled, foundation-normalized, and jointly compiled by the reviewer. We record the outcome with the same honesty discipline as above.

\paragraph{Selected convergents and their errors.}
The formalized integer infrastructure includes the Hermite recurrence,
the adjacent determinant of absolute value two, odd denominators, and
adjacent coprimality. The error estimates for
$\delta_j=Q_j e-P_j$ are the theorems \texttt{deltaSel\_tendsto\_zero}
and \texttt{abs\_deltaSel\_lt}. Both depend only on
\texttt{propext}, \texttt{Classical.choice}, and \texttt{Quot.sound}.
The proof of the error bound uses the Hermite integral representation.

\begin{itemize}
\item \textbf{S1/S2 (B, \texttt{Classification.lean}, 20 decls).} $j\|Q_j x\|\to0\iff x\in L:=\Z+\Z(\ee+1)/2$ (Supplement Theorem~2), and $x\notin L\Rightarrow\limsup_j j\|Q_j x\|\ge\tfrac14$ (Supplement Corollary~3). The proof uses the $\pm2$ determinant to construct integers $a,b$ with $2M_j=aP_j+bQ_j$ (no divisibility division), propagates by the common recurrence, and divides by $Q_j$. The classification and gap use the two proved error estimates. Their audited closure contains only base axioms.
\item \textbf{S3/S4 (C, \texttt{Spectrum.lean}, 30 decls).} The finite spectrum of $j\|Q_j x\|$ is exactly $\{k/4:k\in\N\}$ (Lemma~5 / Theorem~6). \emph{No axiom dependence beyond Lean's base axioms}: the core spectrum argument never uses the $\delta$ error theorems, and the isolated $\delta$-dependent lemma was removed during normalization. (Subtask C uses the placeholder convention $P_0=Q_0=0$, differing from A's $P_0=Q_0=1$; the two agree for $j\ge1$, and C's construction depends on its own $j=0$ convention, so C retains an independent namespace foundation mathematically equivalent to A on $j\ge1$.)
\item \textbf{S5 (D, \texttt{SignedClusters.lean}, 16 decls).} If $S(x)<\infty$ then the cluster set of $j\varepsilon_j$ is $\{k/4:k\text{ taken infinitely often by the integer defect}\}$ (Theorem~7), and every bounded integer sequence is realized as the eventual defect sequence of some $x$ with $S(x)<\infty$ (Theorem~8). Both depend only on base axioms (no $\delta$ error theorems). Theorem~8 is an existence construction freely choosing $x$; it does \emph{not} determine the defect of $x=B\pi$.
\item \textbf{S6 (E, \texttt{DefectFibres.lean}, 32 decls).} For $x,y$ with $S$ finite, equal defect tails $\iff x-y\in L$ (Theorem~9); the defect-tail map is a \emph{bijection} on $F/L$ (Corollary~10, injectivity \emph{and} surjectivity); a fine-scale ratio and a control example ($x=6\pm\ee$). Injectivity depends on \texttt{deltaSel\_\allowbreak{}tendsto\_\allowbreak{}zero}; \texttt{deltaSel\_\allowbreak{}ne\_\allowbreak{}zero} uses the whitelist axiom \texttt{e\_\allowbreak{}transcendental}. \textbf{Surjectivity is now wired} (\texttt{corollary10\_\allowbreak{}surjective}): it invokes D's Theorem~8 realization \texttt{xOfSeq}, bridges the identical defect definitions ($Mseq=Msel=\operatorname{round}(Q_j x)$), and proves the constructed real lies in $F$ via the residual bound; it depends \emph{only on base axioms} (no $\delta$ error theorems), so the full bijection is kernel-checked.
\item \textbf{S7 (F, \texttt{AdjacentWitness.lean}, 32 decls).} A two-index three-way equivalence (Theorem~11) with quantitative bounds. Depends on \texttt{deltaSel\_\allowbreak{}tendsto\_\allowbreak{}zero} (some lemmas base-only). The target ``for every $B$ the matching set is finite'' remains \emph{open}; Theorem~11 is only an equivalent criterion.
\item \textbf{Proposition 3 (H, \texttt{SignedCriterion.lean}, 17 decls).} A signed criterion $e+\pi\in\Q\iff\exists B>0$ such that a signed-residual identity holds eventually (two-way equivalence). Depends on \texttt{deltaSel\_\allowbreak{}tendsto\_\allowbreak{}zero} (converse base-only). The sign estimate (infinitely many positive products for each $B$) is \emph{not established} and is recorded as open.
\item \textbf{Proposition 4 (G, \texttt{FactorialCriterion.lean}, 21 decls).} A factorial criterion using only the exponential series, \emph{with no axioms beyond Lean's base axioms} and no dependence on the selected-convergent foundation. The right-hand limit $n\{-n!\pi\}\to1$ is \emph{not asserted}; it is equivalent to rationality, so an irrationality proof must exclude it.
\end{itemize}
Subtask I provides numerical/symbolic cross-checks (four \texttt{sympy}/\texttt{mpmath} scripts); these are auxiliary, \emph{not} kernel proofs, exclude only the integer case, and do not prove any open problem.

\paragraph{Verification scope.}
A declaration's use of a selected-convergent error theorem is a proof
dependency. The two error theorems and the classification and gap
declarations listed in the companion's verification section have only
the three standard Lean axioms in their audited closure. The subsequent
window-period, digit-shift and denominator-classification results are
given as paper proofs. Problem~\ref{prob:main} remains open.

\subsection{External developments in 2026 and a conditionally verified bridge}\label{sec:lean-external}

We record the relevant external developments of 2026, strictly separating established results from claims that have not been independently verified, and we describe a new conditional bridge that \emph{is} kernel-checked in our project.

\paragraph{The OpenAI $\irr(\pi)=2$ claim (unverified).} OpenAI dated its preprint \emph{The Irrationality Exponent of $\pi$ Is 2} 24 September 2026. The preprint appears in its \texttt{openai/\allowbreak{}math} repository, with a Lean formalization (\texttt{OAI/\allowbreak{}NumberTheory/\allowbreak{}PiExponent/\allowbreak{}}, toolchain Lean v4.34.1, mathlib4 pinned to commit \texttt{d13f23b}, depending on more than thirty external libraries). The proof uses an algebraic-geometric route (admissible-matrix interpolation, cohomology, jets). The repository's scope note (\url{https://github.com/openai/math/blob/main/lean/docs/017.md}, consulted 14 November 2026) describes formalization of the exact exponent and explicitly places the paper's Flint Hills convergence consequence outside that selected formal statement. \textbf{Verification status:} as of October 2026, independent specialist verification of the complete proof had not been reported, and the repository's own README cautions that some unformalized results may have issues. We therefore treat $\irr(\pi)=2$ as a \emph{claim}, not an established theorem. This project has not independently compiled that external proof with its pinned dependencies; \textbf{we do not claim to have independently verified OpenAI's proof}.

\paragraph{Other 2026 developments.} (i) A claimed proof of Flint Hills convergence (arXiv:2502.03474) was \emph{withdrawn} on 14 September 2026; the authors acknowledged that the partial-sum argument does not establish convergence and that a numerical error ($-4$ versus $-30.31$) was present---the same pitfall as our withdrawn comparison in \S\ref{sec:computation}. We do not settle the status of Flint Hills convergence independently of the external claim above. (ii) A no-go audit for $e+\pi$ (arXiv:2606.17303, \cite{Yu2026}) proves exact equivalences between $e+\pi\in\Q$ and eventual factorial-arithmetic phenomena (a ceiling recurrence, a factorial-Cantor digit condition, a divisibility criterion), closely related in spirit to our Proposition~4, and finds no non-circular Ap\'ery-type mechanism within tested low-complexity families; its bounded audit does not rule out all Ap\'ery-type certificates (preprint, not peer-reviewed). (iii) Bertolin's 1-motive program (arXiv:2509.08700, v3 June 2026) \cite{BertolinMotives} gives a torsion-point case of the Grothendieck--Andr\'e period conjecture for 1-motives associated with a CM elliptic curve with algebraic invariants (Proposition~5.4). The $(1,i\pi)$ Schanuel instance is outside that result. (iv) The cited Bai preprint reports the upper bound $\irr(\pi)\le 7.101862832357$ (Bai 2026, arXiv:2609.11276). Schanuel's conjecture and the algebraic independence of $e$ and $\pi$ remain open.

\paragraph{A conditionally verified bridge (kernel-checked).} Independently of OpenAI's heavy proof, we formalized and kernel-checked the \emph{interface} between their claim and our project. In \texttt{EpiFormal/\allowbreak{}OpenAI\_\allowbreak{}mu\_\allowbreak{}pi\_\allowbreak{}bridge.lean} we (1) ingest OpenAI's exact definitions (\texttt{GoodRationalApproximations}, \texttt{ApproximationExponents}, \texttt{irrationalityExponent}, \texttt{EventualLowerBound}) verbatim; (2) prove that their eventual lower bound for $\pi$ implies our $\irr(\pi)\le 2$ (a genuine reconciliation of their $\mathbb R$-valued $\sup$ definition with our $\mathbb{R}_{\ge0}^\infty$-valued measure, via a global rectangle bound on good approximation pairs); and (3) compose with our already-verified Meiburg criterion to obtain:
\begin{theorem}[Conditional bridge; machine-verified]\label{thm:conditional-bridge}
If $\pi$ satisfies the eventual lower bound $\forall\nu>2\,\exists Q\,\forall p,q\ (q\ge Q\Rightarrow q^{-\nu}\le|\pi-p/q|)$ (the form of OpenAI's $\irr(\pi)=2$ claim), then the Flint Hills series converges.
\end{theorem}
The axiom audit of this bridge shows dependencies only on \texttt{pi\_transcendental} and Lean's base axioms, with \emph{no} \texttt{sorryAx}; the reconciliation lemma itself uses only base axioms. The build at that stage was \textbf{9065 jobs} (the final build is reported in \S\ref{sec:lean-status}). \textbf{The premise is OpenAI's unverified claim}; the theorem is conditional. If that claim is ever independently verified, this bridge immediately yields Flint Hills convergence inside our project, the bridge $H'$ is automatically satisfied, and the two surviving targets $\irr(\pi)>2$ and Flint Hills divergence are both closed. The irrationality of $e+\pi$ itself remains open regardless.

\paragraph{Honest status.} This bridge is a verified \emph{conditional} statement, not a proof of $\irr(\pi)=2$ and not a proof of Flint Hills convergence. It demonstrates that our project's verified Meiburg machinery is the ready-made interface that would convert a verified $\irr(\pi)=2$ into Flint Hills convergence. In a period of numerous AI-assisted number-theoretic claims whose independent verification is pending, the value of kernel-checked, axiom-audited formalization is precisely that it provides reproducible, trustworthy conditional infrastructure.

\subsection{Isolating the missing hypothesis: algebraic independence of \texorpdfstring{$e$}{e} and \texorpdfstring{$\pi$}{pi}}\label{sec:lean-ai}

The obstruction analysis of \S\ref{sec:obstructions} and the 2026 literature survey above converge on a single structural diagnosis: the strongest available tools (the determinant/interpolation method behind the $\irr(\log r)=2$ results, and the period machinery) live on the \emph{multiplicative} side---logarithms and periods---whereas $e+\pi$ additively mixes an exponential value $e^1$ with a period $\pi$. The cleanest hypothesis that bridges the two sides is the algebraic independence of $e$ and $\pi$ over $\Q$, which is exactly what Schanuel's conjecture supplies for the pair $\{1,i\pi\}$ (via $e^1=e$ and $e^{i\pi}=-1$). We now \emph{isolate} this sufficient hypothesis as a machine-verified conditional reduction. Its stronger conclusion has a stronger premise; this is not an unconditional strengthening of the theorem in \S\ref{sec:unconditional}.

\begin{theorem}[Algebraic-independence reduction; machine-verified]\label{thm:ai-reduction}
Assume $e$ and $\pi$ are algebraically independent over $\Q$. Then all four elementary combinations
\[
e+\pi,\qquad e-\pi,\qquad e\cdot\pi,\qquad e/\pi
\]
are transcendental over $\Q$.
\end{theorem}

The conclusion under this stronger hypothesis is stronger than the conclusion of Theorem~\ref{thm:vieta} (``at least one of $e+\pi$, $e\pi$ is transcendental''): it closes \emph{all four} combinations at once, at the cost of upgrading the hypothesis from the separate transcendence of $e$ and $\pi$ to their algebraic independence---a statement essentially equivalent to the Schanuel instance needed for the problem.

\paragraph{Proof mechanism (formalized).} In \texttt{EpiFormal/\allowbreak{}AlgebraicIndependence.lean} we argue by contradiction. If a combination $c=f(e,\pi)$ were algebraic, let $m\in\Q[T]$ be a nonzero annihilating polynomial, $m(c)=0$. For the sum, difference, and product we substitute $T\mapsto X_0\pm X_1$ or $X_0X_1$ into $m$, obtaining a bivariate polynomial $P\in\Q[X_0,X_1]$ with $P(e,\pi)=m(c)=0$; $P\neq0$ because each substitution map admits a left inverse (evaluation at $(X,0)$ or $(X,1)$, reducing to the univariate identity), hence is injective. For the quotient $e/\pi$ we use homogenization: $P=\operatorname{homogenize}(m,\deg m)$ satisfies $P(e,\pi)=\pi^{\deg m}\,m(e/\pi)=0$ and $P\neq0$. In every case algebraic independence (injectivity of $\operatorname{aeval}_{(e,\pi)}$) forces $P=0$, a contradiction. The four theorems are named \texttt{transcendental\_\allowbreak{}add/\allowbreak{}sub/\allowbreak{}mul/\allowbreak{}div\_\allowbreak{}of\_\allowbreak{}algebraicIndependent}.

\paragraph{Axiom audit and honest status.} The kernel audit shows that all four theorems depend only on Lean's base axioms, with \emph{no} \texttt{sorryAx} and---notably---\emph{no} dependence on \texttt{e\_\allowbreak{}transcendental} or \texttt{pi\_\allowbreak{}transcendental},\allowbreak  since algebraic independence is a strictly stronger hypothesis that already implies both. \textbf{The premise is open:} the algebraic independence of $e$ and $\pi$ is not proved, and this module does not prove it. The result is a verified \emph{conditional} reduction that pins down precisely which single unproved hypothesis would close the four-combination transcendence question. This isolates one sufficient input. It does not prove that such an input is necessary, exclude other series or approximation constructions, or establish an equivalence between constructing an auxiliary function and proving Schanuel's conjecture; the memo's ``two-type auxiliary construction'' diagnosis is a heuristic judgment, not a proved equivalence, and must not be used to rule out other methods.

\subsection{The escape program and a machine-certified boundary}\label{sec:lean-escape}

The review of the source and proofs redirected the effort from adding further equivalent criteria to proving that $\pi$ \emph{escapes} the structure forced by rationality. Write $D_j(B):=T_j+B\,C_j$ with $M_j=\operatorname{round}(BQ_j\pi)$, $T_j=M_jQ_{j+1}-M_{j+1}Q_j$, $C_j=P_jQ_{j+1}-P_{j+1}Q_j=2(-1)^{j+1}$. By the three-way equivalence of \S\ref{sec:lean-m6}, $B(e+\pi)\in\Z$ iff $D_j(B)=0$ eventually, and also iff it vanishes infinitely often. Negating, the cleanest form of the target is
\[
e+\pi\notin\Q \iff \forall B\ge1,\ \exists N,\ \forall j\ge N:\quad D_j(B)\neq0,
\]
i.e.\ for every $B$ all sufficiently large $j$ are mismatches. The apparently weaker form $\forall B\ge1,\forall N,\exists j\ge N,\ D_j(B)\neq0$ (infinitely many mismatches) is \emph{equivalent}, not merely necessary: if it held while $D_j(B)=0$ infinitely often, the implication (a)$\Rightarrow$(b) would force $D_j(B)=0$ eventually, a contradiction. Thus an alternating match/mismatch sequence is \emph{not} a counterexample. In \texttt{EpiFormal/\allowbreak{}Escape.lean} we formalize the exact identity
\[
D_j(B)=\varepsilon_j Q_{j+1}-\varepsilon_{j+1}Q_j+B\,C_j,\qquad \varepsilon_j=M_j-BQ_j\pi,\quad|\varepsilon_j|\le\tfrac12,
\]
together with both equivalent restatements: \texttt{irrational\_iff\_inf\_mismatch} states the eventual form $\forall B\ge1,\forall^{\mathsf f}j,\ D_j(B)\neq0$, and \texttt{inf\_mismatch\_iff\_eventually} proves the equivalence of the eventual and infinitely-often forms for each $B$. The current axiom audit lists only the three standard Lean axioms for both declarations; the right-hand non-vanishing is \emph{not} proved. Proving it does not reduce the logical difficulty of the open problem.

\paragraph{A machine-certified boundary.} The review stressed that the generic conditions alone cannot exclude eventual matching. We now \emph{certify} this in \texttt{EpiFormal/\allowbreak{}EscapeBoundary.lean}: replacing $\pi$ by $\pi'=A/B-e$, all generic conditions (the $Q_j$ recurrence, $|\varepsilon|\le\tfrac12$, $\delta_j\to0$) still hold, yet
\[
\exists N,\ \forall j\ge N:\quad D'_j=0
\]
(\texttt{boundary\_eventual\_match}, kernel-checked using the proved error theorem, with only base axioms). Hence any argument using only the generic recurrence and rounding \emph{cannot} prove irrationality; the final step must inject a property specific to $\pi$. Numerically (mpmath, 300 digits) the identity holds to $<10^{-157}$ over $205$ pairs $(B,j)$, the true $\pi$ gives $D_j\neq0$ in $100\%$ of $j=0..40$ for $B=1..5$, and the counterexample $\pi'=6-e$ matches from $N=1$. These are corroborations, not proofs.

\paragraph{What is missing (three-route survey).} Parallel literature surveys locate the same gap: a \emph{subsequence lower bound / distribution} fact about $\pi$ that no current tool supplies. The irrationality-measure upper bound $\mu(\pi)\le7.1032$ controls worst-case approximation over \emph{all} large denominators and gives no lower bound on the sparse subsequence $\|BQ_j\pi\|$; it is logically compatible with eventual matching. The sign/size route would require $Q_{j+1}\gg Q_j^{5.1}$  infinitely often. The displayed recurrence rules this out: $Q_{j+1}/Q_j<4j+3$, while $Q_j$ grows at least exponentially. The factorial route needs $\{n!\pi\}\le\tfrac12$ infinitely often, a distribution fact of the same difficulty as $\{(3/2)^n\}$; note the criterion is an equivalence, so proving the limit $=1$ would imply \emph{rationality}---one must \emph{exclude} it. Each route is an independent exploration line; formalizing a criterion does not bring the proof closer.

\subsection{The \texorpdfstring{$D_j(B)$}{Dj(B)} reduction to a lattice-escape target}\label{sec:lean-reduction}

The escape program of \S\ref{sec:lean-escape} restates the target but does not advance it. We now record a \emph{quantitative reduction} that isolates the missing input into a single, precisely stated subproblem. This is a change of target, \emph{not} a reduction in difficulty; the terminal subproblem is of the same family as the original open problem. The full chain is kernel-checked in \texttt{EpiFormal/\allowbreak{}EscapeReduction.lean} (all seven audited declarations depend only on the three standard Lean axioms).

\begin{proposition}[Finite-index test for a fixed slope]
\label{prop:finite-match}
Replace $\pi$ by any real $x$ in $M_k,D_j(B)$. Fix $B\in\N_{>0}$,
put $d=\|B(e+x)\|$, and let $\delta_k=Q_ke-P_k$. For $j\ge1$ one has
\[
 D_j(B)=0\quad\Longrightarrow\quad
 d\le \frac{1/2+B|\delta_{j+1}|}{Q_{j+1}}.
\]
Conversely, $D_j(B)=0$ whenever
$Q_k d+B|\delta_k|<1/2$ for both $k=j,j+1$.
\end{proposition}
\begin{proof}
Write $U_k=M_k+BP_k$. A match says
$U_jQ_{j+1}=U_{j+1}Q_j$. Adjacent denominators are coprime, so
$U_k=AQ_k$ at both indices for one integer $A$. Subtracting gives
\[
 Q_k(B(e+x)-A)=BQ_kx-M_k+B\delta_k .
\]
Taking absolute values at $k=j+1$ proves the upper bound.
Conversely choose $A$ nearest to $B(e+x)$. At each of the two
indices the integer $AQ_k-BP_k$ differs from $BQ_kx$ by at most
$Q_kd+B|\delta_k|<1/2$. It must therefore equal $M_k$,
which gives a match. Strictness removes any dependence on rounding
ties in the converse.
\end{proof}
In particular, a certified bound $d\ge\eta>0$ rules out a match
at every index with $Q_{j+1}\eta>1/2+B|\delta_{j+1}|$.
This separates the finite-index error from the unsolved requirement
$d>0$ for every $B$ when $x=\pi$.

\paragraph{The reduction chain.} Retain the notation of \S\ref{sec:lean-escape}. The pivotal identity $D_j(B)=\varepsilon_jQ_{j+1}-\varepsilon_{j+1}Q_j+BC_j$ with $|\varepsilon_j|\le\tfrac12$ and $|C_j|=2$ gives, on a match $D_j(B)=0$,
\[
|\varepsilon_j(B)|\le\frac{\tfrac12Q_j+2B}{Q_{j+1}}
\qquad\text{(\texttt{match\_eps\_bound}).}
\]
Since $Q_{j+1}=(4j+2)Q_j+Q_{j-1}$ with $Q_{j-1}/Q_j\to0$ (\texttt{Qsel\_pred\_div\_tendsto\_zero}), one has $jQ_j/Q_{j+1}\to\tfrac14$ (\texttt{jQsel\_div\_Qsel\_succ\_tendsto}), hence
\[
j\,\frac{\tfrac12Q_j+2B}{Q_{j+1}}\longrightarrow\tfrac18
\qquad\text{(\texttt{sharp\_bound\_tendsto}).}
\]
In particular the right-hand side is eventually $\le\tfrac15$. Combining with the match bound, if $D_j(B)=0$ eventually then $j\,\|BQ_j\pi\|\le\tfrac15$ eventually (\texttt{match\_eventually\_eps\_le\_fifth}), where we use $\|BQ_j\pi\|=|\varepsilon_j(B)|$ (the sign convention of \texttt{Classification.eps} is opposite to \texttt{Escape.eps}, but the absolute values agree).

\paragraph{Collision with the classification gap.} The kernel-checked Corollary~3 of the classification (\texttt{Classification.\allowbreak{}corollary3\_\allowbreak{}gap}, \S\ref{sec:lean-m6}) states that for $x\notin L:=\Z+\Z\frac{e+1}{2}$ and any $c<\tfrac14$, it is \emph{not} the case that $j\|Q_jx\|\le c$ eventually. Taking $x=B\pi$ and $c=\tfrac15<\tfrac14$, the eventual bound above contradicts $B\pi\notin L$. Hence
\[
\forall B\ge1:\quad D_j(B)=0\ \text{eventually}\ \Longrightarrow\ B\pi\in L
\qquad\text{(\texttt{reduction\_to\_lattice}).}
\]
Equivalently, with the open target
\[
\mathrm{P1}:\qquad \forall B\ge1,\quad B\pi\notin L,
\]
one has $\mathrm{P1}\Rightarrow e+\pi\notin\Q$ (\texttt{P1\_imp\_irrational}, a \emph{conditional} theorem: its premise P1 is not proved).

\begin{proposition}[Rational span of the exceptional lattice]
\label{prop:rational-span}
For $L=\Z+\Z(e+1)/2$ one has
\[
\bigcup_{B\ge1} B^{-1}L=\Q+\Q e,\qquad B\in\N.
\]
Thus P1 is equivalent to linear independence of $1,e,\pi$ over $\Q$.
\end{proposition}
\begin{proof}
Each $B^{-1}L$ lies in $\Q+\Q e$. Conversely, for $x=r+se$ with
$r,s\in\Q$, choose $B\ge1$ clearing both denominators. Then
$Bx=B(r-s)+2Bs(e+1)/2\in L$. Since $e$ is irrational, every
nontrivial rational relation among $1,e,\pi$ has a nonzero
$\pi$-coefficient and can be solved for $\pi$. This proves the assertion.
In particular $B(e+\pi)=A$ implies
$B\pi=(A+B)-2B(e+1)/2\in L$ directly.
\end{proof}
The slope in the original question is fixed at $-1$; P1 excludes every
rational slope. For instance the substitute $x=e$ violates the
corresponding lattice exclusion although $e+x=2e$ is irrational.
The quantitative match estimate describes rounding errors, while the
proposition identifies the precise arithmetic premise.

\paragraph{Numerical baseline (corroboration, not proof).} A rebuilt baseline (\texttt{compute\_dj.py}, mpmath, dynamic precision $\mathrm{digits}(Q_j)+500$) verifies over $B=1..20$, $j=1..80$ (1600 pairs): the recurrence and determinant $C_j=2(-1)^{j+1}$ hold termwise; $D_j(B)=0$ occurs $0$ times for the true $\pi$; the identity holds to relative error $<10^{-100}$; the boundary counterexamples $\pi'=6-e$ ($B=1$) and $\pi'=\tfrac{11}{2}-e$ ($B=2$) match eventually as predicted. The upper-bound expression $j(\tfrac12Q_j+2B)/Q_{j+1}$ increases toward $\tfrac18$ (values $0.113,0.119,0.122,0.123$ at $j=5,10,20,40$ for $B=1$). The sign statistic $\operatorname{sgn}(R_j\delta_j)>0$ occurs at ratio $0.42$--$0.48$ over $j=1..200$, $B=1..5$, consistent with no bias (the sign route carries no signal). These are finite computations, not proofs.

\subsection{Feasibility audit of the attack routes}\label{sec:feasibility}

A structured audit in eight sub-investigations examined every route we could identify against the precise target isolated above, equivalently $\forall B\ge1,\ \limsup_j j\|BQ_j\pi\|>\tfrac18$. \textbf{None yields an unconditional advance; all confirm the problem is open.} We record the verdicts as a negative-result audit, in the spirit of Yu's no-go audit \cite{Yu2026}.

\paragraph{Routes toward P1.}
\emph{Quantitative gap (2C).} The classification gap is already \emph{uniform}: the quantitative lemma never uses $x\notin L$, which appears only as the terminal contradiction hypothesis. Hence $\limsup_j j\|Q_jx\|\ge\tfrac14$ for all $x\notin L$, independent of $\operatorname{dist}(x,L)$; a continuous $\varphi(\operatorname{dist})$ is unnecessary. The exact sufficient target is the sparse lower bound above. The known $\mu(\pi)\le7.1032$ bound decays to $0$ along the super-exponential $Q_j$ (gap $\sim10^{1010}$ at $j=80$), so it cannot supply the target. \emph{Baker linear forms (2A).} $e$ has no known representation as a logarithm of an algebraic number (asserting otherwise is itself open, tied to $e^e$); the target does not enter Baker's framework; the four/six-exponentials inventory is blocked exactly as Proposition~\ref{prop:four-exp}. \emph{Continued fractions (2B).} No known structural property of $\pi$'s partial quotients combines with the transducer theorem to force a contradiction; numerically the match length is constantly $1$--$2$ and does not grow with precision (no numerical handle). \emph{$\pi$-representations (Stage~3).} Machin/BBP/Chudnovsky produce no usable estimate; only the AGM/Legendre period representation passes the boundary test (``$\pi$ is a period''), but yields merely an interface diagram depending on the expected-false $\pi\in S$ and unreviewed preprints.

\paragraph{Fresh angles.}
\emph{Integral positivity.} Niven/Beukers-type integrals give irrationality or power-type approximation \emph{upper} bounds, never the required subsequence \emph{lower} bound (the Beukers paradigm sends the integral to $0$, the wrong direction). \emph{Nesterenko contradiction mining.} Seven candidate paths from $e+\pi\in\Q$ to a violation of a proved theorem all fail: the proved theorems control exponentials at algebraic points, algebraic-base powers, and the modular triple $\{\pi,e^\pi,\Gamma(1/4)\}$, but \emph{none controls an algebraic relation between $e$ and $\pi$}; the hypothesis falls precisely in this blank. \emph{Period/$E$-$G$ conditional theorems.} We obtain two precise conditional results: $(\pi\notin\mathbf{E})\Rightarrow$ all four combinations transcendental, and $(\mathbf{E}\cap\mathbf{G}=\overline{\Q})\Rightarrow\pi\notin\mathbf{E}$. The inclusions $\Q+\Q e\subseteq\mathbf E$ and $\pi\in\mathbf G$ give the chain $(\mathbf E\cap\mathbf G=\overline{\Q})\Rightarrow(\pi\notin\mathbf E)\Rightarrow\mathrm{P1}$. Schanuel at $(1,i\pi)$ gives P1 separately through algebraic independence. The restricted unit theorem of \cite{FischlerRivoal2026} gives the conditional inverse-membership implication in Section~\ref{sec:expected-false}. \emph{Literature scan (2024--2026).} No new result supplies the missing input; even the claimed $\mu(\pi)=2$ \cite{OpenAI2026}, if verified, is non-effective and gives no subsequence lower bound, hence is of the wrong type.

\paragraph{Phase~1 conclusion.} The adjacent-index criterion concerns the
single slope $-1$. The classification and gap identify P1 with the
stronger target $\limsup_j j\|BQ_j\pi\|>1/8$ for every positive integer
$B$: points of $L$ have limit zero and points outside have limsup at
least $1/4$. The hypothesis $\pi\notin\mathbf E$ suffices for P1.
The following dimension calculation describes the size of the
exceptional sets in the ambient real line.

\subsection{Hausdorff dimension for eventual approximation}
\label{sec:eventual-dimension}
\begin{theorem}\label{thm:eventual-dimension}
For the selected denominators $Q_j$ and $\tau\ge0$, put
\[
 F_\tau=\{x\in[0,1]:\ \|Q_jx\|=O(j^{-\tau})\}.
\]
Then $\dim_{\rm H}F_\tau=\max(1-\tau,0)$. In particular the
entire set on which $\limsup_j j\|Q_jx\|$ is finite has Hausdorff
dimension zero.
\end{theorem}
\begin{proof}
The case $\tau=0$ follows from $\|y\|\le1/2$. The recurrence gives
$4j+2\le Q_{j+1}/Q_j\le4j+3$, and hence
$\log Q_j=j\log j+O(j)$ by summation. First fix $0<\tau\le1$.
For positive integers $C,N$, require $\|Q_jx\|\le Cj^{-\tau}$
at every $j\ge N$. At level $j$ the allowable intervals have centres
$k/Q_j$ and radius $r_j=Cj^{-\tau}/Q_j$. Each meets at most
$1+2(r_j+r_{j+1})Q_{j+1}=O_C(j^{1-\tau})$ next-level
intervals, by counting grid points in the enlarged parent.
The number of chains through level $n$ is therefore at most
$\exp((1-\tau)n\log n+O_{C,N}(n))$. For $s>1-\tau$,
multiplying this bound by $(2r_n)^s$ gives a quantity tending to
zero. Countable union over $C,N$ yields the upper bound.

For $0<\tau<1$ take $C=1/8$ and start with one interval inside
$(0,1)$ at a sufficiently large level $N$. Each retained parent
contains at least $2(r_j-r_{j+1})Q_{j+1}-1$ whole child intervals.
This is bounded below by $c j^{1-\tau}$ with an absolute $c>0$
for all sufficiently large $j$. Choose exactly
$b_j=\lfloor c j^{1-\tau}\rfloor\ge2$ children per parent.
The nested intersection lies in $F_\tau$. Give the children equal
shares of their parent's mass. Each level-$n$ interval then has mass
$1/M_n$, where $M_n=\prod_{j=N}^{n-1}b_j$ and
$\log M_n=(1-\tau)n\log n+O(n)$.
For an interval of length $Q_{n+1}^{-1}\le h<Q_n^{-1}$, the
level-$n$ grid spacing and $2r_n<Q_n^{-1}$ imply that it meets
at most three retained intervals. Its mass is at most $3/M_n$.
For every $0<s<1-\tau$, the growth estimates give
$3/M_n\le Q_{n+1}^{-s}\le h^s$ for all large $n$.
Summing over any sufficiently fine interval cover of the intersection
shows its $s$-dimensional Hausdorff measure is positive. Letting
$s$ increase to $1-\tau$ proves the lower bound.
For $\tau\ge1$, use $0\in F_\tau\subseteq F_1$.
\end{proof}
This elementary covering argument uses only the growth of the
denominators. It also gives dimension zero for each finite level
set of the spectrum. It leaves the membership of the particular point
$\{\pi\}$ undecided. The corrected Phase-2 companion gives prescribed-limit sets at subcritical weights and a critical Hausdorff gauge. The rational-span proposition and dimension theorem in this
revision are proved on paper; the historical Lean compilation records
refer to the earlier declarations.

\section{Theorem--source--application audit trail}\label{app:audit}

This is a \emph{summary} audit table of the principal external facts, with verification method and application sites; it is not exhaustive. The declaration-level audit and the build/gap ledger are in the archive as \texttt{LEMMA\_AUDIT.md}, \texttt{PROOF\_GAPS.md}, and \texttt{FINAL\_REVIEW.md} (the earlier ``50-entry ledger'' is superseded by these files). Discipline: preprints are conditional; weak (popular-science) sources are flagged and, where possible, replaced by primary sources in the bibliography.

\begin{center}
\footnotesize\begin{tabularx}{\textwidth}{@{}p{0.19\textwidth}p{0.17\textwidth}p{0.27\textwidth}X@{}}
\toprule
Fact & Source & Verification & Applications \\
\midrule
$\ee,\pi$ transcendental & \cite{Hermite1873,Lindemann1882} & primary citations & Thms.~\ref{thm:vieta},\ref{thm:four},\ref{thm:schanuel-consequence} \\
Hermite--Lindemann / LW & \cite{Lindemann1882,Bertolin2026} & restatement in read paper & \S\ref{sub:route-c}; Prop.~\ref{prop:four-exp} \\
Gelfond--Schneider & \cite{GS1934} & classical statement & Thm.~\ref{thm:epi} \\
Schneider elliptic & \cite{Duke} & direct quotation & Prop.~\ref{prop:scaling} remark \\
$\irr(\ee^r)=2$ & \cite{CNU2024} & theorem quotation & Lem.~\ref{lem:mue}, Thm.~\ref{thm:measure-reduction} \\
$\irr(\pi)$ bounds history & \cite{Alekseyev2011,ZZ2020} & full text + journal record & \S\ref{sec:new}.1 remark; Thm.~\ref{thm:translation} \\
Serret / Raney & \cite{Panti2017,Raney1973} & abstract / publisher record & Thm.~\ref{thm:cf} \\
Alekseyev criteria & \cite{Alekseyev2011} & full-text close reading & Thm.~\ref{thm:alekseyev}; correction \\
Meiburg criteria & \cite{Meiburg2022} & proof structure verified; Flint Hills case machine-verified & Thms.~\ref{thm:meiburg},\ref{thm:flint} (classical axioms listed) \\
Bertolin theorem & \cite{Bertolin2026} & abstract+intro+proof opening & Cor.~\ref{cor:bertolin} (conditional) \\
Schanuel statement/status & \cite{Ax1971,Baker,MO33817} & primary literature + MO & Thms.~\ref{thm:schanuel},\ref{thm:schanuel-consequence} \\
Nesterenko / Tijdeman & \cite{Nesterenko1996,Tijdemann1971} & bibliography & \S\ref{sub:stronger-known}; \S\ref{sec:computation}.5 \\
non-$U$ facts & \cite{Baker} & monograph & \S\ref{sub:mahler} \\
Ax--Schanuel line & \cite{Ax1971,MPT2019,BakkerTsimerman2019,CambridgeFF} & primary literature + volume & \S\ref{sub:ax} \\
Zilber program & \cite{MO33817} & accepted answer via API & \S\ref{sub:zilber} \\
Integral identities & sympy & symbolic, 12/12 PASS & \S\ref{sub:bridges}; \S\ref{sec:computation}.1 \\
Interval certificates & mpmath iv & 450-digit, reproducible & \S\ref{sec:computation}.2 \\
Chowla--Selberg identity & numerical only & 120-digit, literature unverified & \S\ref{sec:computation}.5 (flagged) \\
Legendre relation & \cite{DLMF23} & framework verified, constants not & \S\ref{sub:motivic} (flagged) \\
\bottomrule
\end{tabularx}
\end{center}

Correction log: (1) Round 7 corrected the Round 6 attribution of the $\irr(\pi)<5/2$ convergence direction from Alekseyev to Meiburg (secondary-source error caught by primary-source reading); the Round 6 notes carry a permanent correction banner. (2) Round 9: a typesetting artifact in the parsed Meiburg manuscript ($f'$ formula) was not trusted; the derivative was recomputed from the definition of $f$ and matched the paper's conclusion. (3) 2026-11-14: the comparison of $S(10000)=30.31451083$ with the literature's claimed limit $30.3144\dots$ in \S\ref{sec:computation}.4 is withdrawn (positive terms force any finite limit to exceed the partial sum). (4) 2026-11-14: the Alekseyev direction, previously recorded as a research-level density gap (\S\ref{sec:alekseyev}), is closed by the nearest-integer argument of the revised supplement and fully machine-verified (\texttt{alekseyev\_\allowbreak{}bound}); the failed termwise route is preserved as a historical record.

\section{Reproducibility}\label{app:repro}

The authoritative source is
\texttt{paper1\_\allowbreak e\_\allowbreak plus\_\allowbreak pi/\allowbreak 01\_\allowbreak main\_\allowbreak paper/\allowbreak paper1\_\allowbreak e\_\allowbreak plus\_\allowbreak pi.tex};
the PDF is regenerated from it with two passes of \texttt{pdflatex}.
The reproducible numerical and symbolic verification suite resides in
\texttt{verification/} (together with the task-specific certificates under
\texttt{paper1\_\allowbreak e\_\allowbreak plus\_\allowbreak pi/\allowbreak 05\_\allowbreak research\_\allowbreak process/},
bundled under \texttt{verification/\allowbreak paper1\_\allowbreak companion\_\allowbreak scripts/}
in the self-contained delivery archive):
\texttt{verify\_\allowbreak round3.py} (12/12 SymPy~1.14.0 symbolic checks, log
\texttt{verify\_\allowbreak round3.log}), \texttt{verify\_\allowbreak round4.py}
(cosine integral and 450-digit \texttt{mpmath.iv} interval certificate,
outputs \texttt{certificate\_\allowbreak round4.json} and
\texttt{verify\_\allowbreak round4.log}),
\texttt{verify\_\allowbreak golden\_\allowbreak orbit\_\allowbreak model.py},
\texttt{verify\_\allowbreak golden\_\allowbreak b2\_\allowbreak exact.py},
\texttt{verify\_\allowbreak golden\_\allowbreak b3\_\allowbreak b4\_\allowbreak exact.py},
\texttt{verify\_\allowbreak golden\_\allowbreak b5\_\allowbreak b6\_\allowbreak multiscale.py},
\texttt{verify\_\allowbreak golden\_\allowbreak general\_\allowbreak b.py},
\texttt{check\_\allowbreak audit\_\allowbreak reconcile.py}, and the
companion certificates \texttt{golden\_\allowbreak finite\_\allowbreak spectrum.py},
\texttt{verify\_\allowbreak ten\_\allowbreak rounds.py},
\texttt{verify\_\allowbreak trace\_\allowbreak markov.py}, and
\texttt{verify\_\allowbreak research\_\allowbreak advance.py}
(earlier archival scripts \texttt{verify\_\allowbreak round2.py},
\texttt{verify\_\allowbreak round6.py}, and \texttt{compute\_\allowbreak dj.py}
are historical baseline runs not retained as active scripts). Their
symbolic or numerical output is not a Lean kernel proof.

The formal project is in \texttt{epi\_formal/}, pinned by
\texttt{lean-toolchain} (\texttt{leanprover/lean4:v4.35.0}) and the
dependency manifest (\texttt{lakefile.toml}, \texttt{lake-manifest.json}).
After installing the dependencies, run \texttt{lake build EpiFormal} and
\texttt{lake env lean Audit.lean} from the \texttt{epi\_formal/} directory,
then run \texttt{verification/\allowbreak check\_\allowbreak audit\_\allowbreak reconcile.py}
to check exact declaration bijection, zero \texttt{sorryAx}, zero
project-added axioms, and zero escape keywords. \S\ref{sec:lean-status}
gives the audit scope and build record. The escape formulations are
checked in \texttt{EpiFormal/\allowbreak Escape.lean} and
\texttt{EpiFormal/\allowbreak EscapeBoundary.lean}; see
\S\ref{sec:lean-escape}.

For Theorem~\ref{thm:golden-finite-shift} and
Corollary~\ref{cor:golden-finite-shift-examples}, the verification suite
includes \texttt{golden\_\allowbreak finite\_\allowbreak spectrum.py}
(run without arguments for the thirteen benchmark cases, or pass
\texttt{22 7 --cap 1/100 --all} to list every limit for that input
within the stated cap), while \texttt{GoldenHeadCutoff.lean} and the
other seven Paper~1 helper modules (\texttt{GoldenAffineEndpoint.lean},
\texttt{AffineWindowAudit.lean}, \texttt{DigitTailBound.lean},
\texttt{DensityCube.lean}, \texttt{OrbitAndFlow.lean},
\texttt{ResearchAdvance.lean}, and \texttt{TraceMarkov.lean}) are now
integrated directly inside \texttt{epi\_formal/\allowbreak EpiFormal/}
and audited in \texttt{Audit.lean}. The 2026-11-22 rerun built the full
integrated project (\texttt{9179 jobs}, exit code~0) and verified all
937 audited declarations across 134 Lean files (914 base-only + 23 whitelisted classical,
zero \texttt{sorryAx}, zero project-added axioms) in
\texttt{verification/\allowbreak audit\_\allowbreak final\_\allowbreak 20261122.log}.

\end{document}


